%============================================================
%  Bd. 13 - Ganze Zahlen
%============================================================

\documentclass[main.tex]{subfiles}

\ifSubfilesClassLoaded{
  \BandSetupStandalone{B13}
}{
}

\title{Bd. 13 - Ganze Zahlen}
\author{Martin Kunze}
\date{}
\setcounter{file}{13}

\begin{document}

\title{Bd. 13 - Ganze Zahlen}
\author{Martin Kunze}
\date{}
\hypersetup{pageanchor=false}
\maketitle
\hypersetup{pageanchor=true}
\setcounter{tocdepth}{3}
\tableofcontents
\setcounter{file}{13}
\renewcommand{\FormulaBandID}{13}

\chapter{Einleitung}

Die ganzen Zahlen entstehen aus Paaren natürlicher Zahlen. Das Paar
\((a,b)\) wird als formale Differenz \(a-b\) gelesen. Da verschiedene Paare
dieselbe Differenz beschreiben können, werden genau diejenigen Paare
identifiziert, deren Kreuzsummen übereinstimmen. Die allgemeine Theorie der
Äquivalenzrelationen und Quotienten aus Band~10 liefert anschließend die
Menge der ganzen Zahlen.

Nach der Konstruktion werden Nachfolger, Vorgänger und die zweiseitige
Induktion zunächst im Quotientenmodell bewiesen. Erst danach werden diese
Eigenschaften als erweiterte Peano-Axiome registriert. Addition,
Multiplikation und Negation werden stets erst nach einem Beweis ihrer
Repräsentantenunabhängigkeit eingeführt.

\chapter{Konstruktion der ganzen Zahlen}

\section{Differenzenpaare}

\FormulaDefDeltaK[Träger der Differenzenpaare]{%
  D_{\mathbb Z}\coloneqq\mathbb{N}\times\mathbb{N}%
}{IntegerPairCarrier}{%
  \DeltaRow{Mengen}{D_{\mathbb Z}}%
  \DeltaRow{Neue Symbole}{D_{\mathbb Z}}%
}

\FormulaDefDeltaK[Äquivalenz von Differenzenpaaren]{%
  a,b,c,d\in\mathbb{N}\vdash
  (a,b)\IntPairRel(c,d)\coloneqq a+d=b+c%
}{IntPairRelDef}{%
  \DeltaRow{Mengen}{a\dsep b\dsep c\dsep d}%
  \DeltaRow{Relation}{\IntPairRel}%
}

\FormulaDefDeltaK[Relation der Differenzenpaare]{%
  \IntRel\coloneqq
  \bigl\{((a,b),(c,d))\in
  D_{\mathbb Z}\times D_{\mathbb Z}\mid
  a+d=b+c\bigr\}%
}{IntRelationDef}{%
  \DeltaRow{Mengen}{a\dsep b\dsep c\dsep d\dsep D_{\mathbb Z}}%
  \DeltaRow{Relation}{\IntRel}%
  \DeltaRow{Neue Symbole}{\IntRel}%
}

\FormulaThmDeltaK[Charakterisierung der Differenzenrelation]{%
  a,b,c,d\in\mathbb{N}\vdash
  ((a,b),(c,d))\in\IntRel
  \leftrightarrow
  a+d=b+c%
}{IntRelationChar}{%
  \DeltaRow{Mengen}{a\dsep b\dsep c\dsep d}%
  \DeltaRow{Relation}{\IntRel}%
}
\begin{tabproofsplit}
  \proofpart{\(\vdash\)}
    \proofstep{1}{a,b,c,d\in\mathbb{N}}{\rA}
    \proofstep{2}{((a,b),(c,d))\in\IntRel}{\rA}
    \proofstep{}{%
      \IntRel=
      \{((u,v),(x,y))\in D_{\mathbb Z}\times D_{\mathbb Z}
      \mid u+y=v+x\}}{\FormulaRefAuto{IntRelationDef}}
    \proofstep{1,2}{a+d=b+c}{%
      \FormulaRefAuto{x \in \{x \in A \mid P(x)\}\vdash P(x)}{3,2}}
  \closeproofpart

  \proofpart{\(\dashv\)}
    \proofstep{1}{a,b,c,d\in\mathbb{N}}{\rA}
    \proofstep{2}{a+d=b+c}{\rA}
    \proofstep{1}{(a,b)\in D_{\mathbb Z}}{%
      \FormulaRefAuto{(a,b)\in A\times B\eqvdash a\in A\land b\in B}{1}}
    \proofstep{1}{(c,d)\in D_{\mathbb Z}}{%
      \FormulaRefAuto{(a,b)\in A\times B\eqvdash a\in A\land b\in B}{1}}
    \proofstep{1}{((a,b),(c,d))\in
      D_{\mathbb Z}\times D_{\mathbb Z}}{%
      \FormulaRefAuto{(a,b)\in A\times B\eqvdash a\in A\land b\in B}{3,4}}
    \proofstep{}{%
      \IntRel=
      \{((u,v),(x,y))\in D_{\mathbb Z}\times D_{\mathbb Z}
      \mid u+y=v+x\}}{\FormulaRefAuto{IntRelationDef}}
    \proofstep{1,2}{((a,b),(c,d))\in\IntRel}{%
      \FormulaRefAuto{x \in \{u \in A \mid P(u)\} \eqvdash
      x \in A \land P(x)}{5,2,6}}
  \closeproofpart
\end{tabproofsplit}

\FormulaThmDeltaK[Trägerbedingung der Differenzenrelation]{%
  \IntRel\subseteq D_{\mathbb Z}\times D_{\mathbb Z}%
}{IntRelationCarrier}{%
  \DeltaRow{Mengen}{D_{\mathbb Z}}%
  \DeltaRow{Relation}{\IntRel}%
}
\begin{tabproof}
  \proofstep{1}{x\in\IntRel}{\rA}
  \proofstep{}{%
    \IntRel=
    \{((a,b),(c,d))\in D_{\mathbb Z}\times D_{\mathbb Z}
    \mid a+d=b+c\}}{\FormulaRefAuto{IntRelationDef}}
  \proofstep{1}{x\in D_{\mathbb Z}\times D_{\mathbb Z}}{%
    \FormulaRefAuto{x \in \{x \in A \mid P(x)\}\vdash x\in A}{2,1}}
  \proofstep{}{\IntRel\subseteq D_{\mathbb Z}\times D_{\mathbb Z}}{%
    \FormulaRefAuto{A\subseteq B \coloneq
      \forall x\,(x\in A\rightarrow x\in B)}{\rUI{\rRI{1,3}}}}
\end{tabproof}

\FormulaThmDeltaK[Reflexivität auf Differenzenpaaren]{%
  a,b\in\mathbb{N}\vdash
  ((a,b),(a,b))\in\IntRel%
}{IntRelationReflexivePairs}{%
  \DeltaRow{Mengen}{a\dsep b}%
  \DeltaRow{Relation}{\IntRel}%
}
\begin{tabproof}
  \proofstep{1}{a,b\in\mathbb{N}}{\rA}
  \proofstep{1}{a+b=b+a}{\FormulaRefAuto{PeanoAddCommutative}{1}}
  \proofstep{1}{((a,b),(a,b))\in\IntRel}{%
    \FormulaRefAuto{IntRelationChar}{1,2}}
\end{tabproof}

\FormulaThmDeltaK[Symmetrie auf Differenzenpaaren]{%
  a,b,c,d\in\mathbb{N}\dsep
  ((a,b),(c,d))\in\IntRel
  \vdash ((c,d),(a,b))\in\IntRel%
}{IntRelationSymmetricPairs}{%
  \DeltaRow{Mengen}{a\dsep b\dsep c\dsep d}%
  \DeltaRow{Relation}{\IntRel}%
}
\begin{tabproofwide}
  \proofstepwidestar[1]{a,b,c,d\in\mathbb{N}}{\rA}
  \proofstepwidestar[2]{((a,b),(c,d))\in\IntRel}{\rA}
  \proofstepwide[1,2]{a+d}{=}{b+c}{%
    \FormulaRefAuto{IntRelationChar}{1,2}}
  \proofstepwide[1,2]{c+b}{=}{b+c}{%
    \FormulaRefAuto{PeanoAddCommutative}{1}}
  \proofstepwide[1,2]{b+c}{=}{a+d}{%
    \FormulaRefAuto{a=b\vdash b=a}{3}}
  \proofstepwide[1,2]{a+d}{=}{d+a}{%
    \FormulaRefAuto{PeanoAddCommutative}{1}}
  \proofstepwide[1,2]{c+b}{=}{d+a}{%
    \rIE{4,5,6}}
  \proofstepwidestar[1,2]{((c,d),(a,b))\in\IntRel}{%
    \FormulaRefAuto{IntRelationChar}{1,7}}
\end{tabproofwide}

\FormulaThmDeltaK[Transitivität auf Differenzenpaaren]{%
  \begin{aligned}[t]
    &a,b,c,d,e,f\in\mathbb{N}\dsep
      ((a,b),(c,d))\in\IntRel\dsep{}\\
    &((c,d),(e,f))\in\IntRel
      \vdash ((a,b),(e,f))\in\IntRel
  \end{aligned}%
}{IntRelationTransitivePairs}{%
  \DeltaRow{Mengen}{a\dsep b\dsep c\dsep d\dsep e\dsep f}%
  \DeltaRow{Relation}{\IntRel}%
}
\begin{tabproofsplitwide}
  \proofpartwideindR[Zusammenführen der Kreuzsummen]{%
    \begin{aligned}[t]
      &a,b,c,d,e,f\in\mathbb{N}\dsep
      a+d=b+c\dsep c+f=d+e\vdash{}\\
      &(a+f)+(c+d)=(b+e)+(c+d)
    \end{aligned}
  }{%
    a,b,c,d,e,f\in\mathbb{N}\dsep
    a+d=b+c\dsep c+f=d+e
    \vdash (a+f)+(c+d)=(b+e)+(c+d)
  }[IntRelationTransitiveRearrangement]
    \proofstepwidestar[1]{a,b,c,d,e,f\in\mathbb{N}}{\rA}
    \proofstepwidestar[2]{a+d=b+c}{\rA}
    \proofstepwidestar[3]{c+f=d+e}{\rA}
    \proofstepwide[1,2,3]{(a+d)+(c+f)}{=}{(b+c)+(d+e)}{\rIE{2,3}}
    \proofstepwide[1]{(a+d)+(c+f)}{=}{(a+f)+(c+d)}{%
      \FormulaRefAuto{PeanoAddInterchange}{1}}
    \proofstepwide[1]{(b+c)+(d+e)}{=}{(b+e)+(c+d)}{%
      \FormulaRefAuto{PeanoAddInterchange}{1}}
    \proofstepwide[1,2,3]{(a+f)+(c+d)}{=}{(b+e)+(c+d)}{%
      \rIE{4,5,6}}
  \closeproofpartwideindR

  \proofpartwideindR[Kürzen der gemeinsamen Summe]{%
    \begin{aligned}[t]
      &a,b,c,d,e,f\in\mathbb{N}\dsep
      (a+f)+(c+d)=(b+e)+(c+d)\\
      &\vdash a+f=b+e
    \end{aligned}
  }{%
    a,b,c,d,e,f\in\mathbb{N}\dsep
    (a+f)+(c+d)=(b+e)+(c+d)\vdash a+f=b+e
  }[IntRelationTransitiveCancellation]
    \proofstepwidestar[1]{a,b,c,d,e,f\in\mathbb{N}}{\rA}
    \proofstepwidestar[2]{(a+f)+(c+d)=(b+e)+(c+d)}{\rA}
    \proofstepwidestar[1]{a+f\in\mathbb{N}}{\FormulaRefAuto{PeanoAddClosure}{1}}
    \proofstepwidestar[1]{b+e\in\mathbb{N}}{\FormulaRefAuto{PeanoAddClosure}{1}}
    \proofstepwidestar[1]{c+d\in\mathbb{N}}{\FormulaRefAuto{PeanoAddClosure}{1}}
    \proofstepwidestar[1,2]{a+f=b+e}{%
      \FormulaRefAuto{PeanoAddRightCancellation}{3,4,5,2}}
  \closeproofpartwideindR

  \proofpartwide{Schluss}
    \proofstepwidestar[1]{a,b,c,d,e,f\in\mathbb{N}}{\rA}
    \proofstepwidestar[2]{((a,b),(c,d))\in\IntRel}{\rA}
    \proofstepwidestar[3]{((c,d),(e,f))\in\IntRel}{\rA}
    \proofstepwidestar[1,2]{a+d=b+c}{\FormulaRefAuto{IntRelationChar}{1,2}}
    \proofstepwidestar[1,3]{c+f=d+e}{\FormulaRefAuto{IntRelationChar}{1,3}}
    \proofstepwidestar[1,2,3]{(a+f)+(c+d)=(b+e)+(c+d)}{%
      \FormulaRefAuto{IntRelationTransitiveRearrangement}{1,4,5}}
    \proofstepwidestar[1,2,3]{a+f=b+e}{%
      \FormulaRefAuto{IntRelationTransitiveCancellation}{1,6}}
    \proofstepwidestar[1,2,3]{((a,b),(e,f))\in\IntRel}{%
      \FormulaRefAuto{IntRelationChar}{1,7}}
  \closeproofpartwide
\end{tabproofsplitwide}

\FormulaThmDeltaK[Die Differenzenrelation ist eine Äquivalenzrelation]{%
  \EqRel{\IntRel,D_{\mathbb Z}}%
}{IntegerPairEquivalence}{%
  \DeltaRow{Mengen}{D_{\mathbb Z}}%
  \DeltaRow{Relation}{\IntRel}%
}
\begin{tabproofsplitwide}
  \proofpartwideindR[Relation auf dem Paarträger]{%
    \IntRel\subseteq D_{\mathbb Z}\times D_{\mathbb Z}%
  }{%
    \IntRel\subseteq D_{\mathbb Z}\times D_{\mathbb Z}%
  }[IntegerPairEquivalenceCarrier]
    \proofstepwidestar[]{\IntRel\subseteq
      D_{\mathbb Z}\times D_{\mathbb Z}}{%
      \FormulaRefAuto{IntRelationCarrier}}
  \closeproofpartwideindR

  \proofpartwideindR[Reflexivität]{%
    x\in D_{\mathbb Z}\vdash(x,x)\in\IntRel%
  }{%
    x\in D_{\mathbb Z}\vdash(x,x)\in\IntRel%
  }[IntegerPairEquivalenceReflexive]
    \proofstepwidestar[1]{x\in D_{\mathbb Z}}{\rA}
    \proofstepwidestar[1]{\exists a\in\mathbb{N}\,
      \exists b\in\mathbb{N}\,x=(a,b)}{%
      \FormulaRefAuto{x\in A\times B\eqvdash
        \exists a\in A\,\exists b\in B\,x=(a,b)}{1}}
    \proofstepwidestar[3]{a\in\mathbb{N}\land
      b\in\mathbb{N}\land x=(a,b)}{\rA}
    \proofstepwidestar[3]{((a,b),(a,b))\in\IntRel}{%
      \FormulaRefAuto{IntRelationReflexivePairs}{3}}
    \proofstepwidestar[3]{(x,x)\in\IntRel}{\rIE{3,4}}
    \proofstepwidestar[1]{(x,x)\in\IntRel}{\rEE{2,3,5}}
  \closeproofpartwideindR

  \proofpartwideindR[Symmetrie]{%
    (x,y)\in\IntRel\vdash(y,x)\in\IntRel%
  }{%
    (x,y)\in\IntRel\vdash(y,x)\in\IntRel%
  }[IntegerPairEquivalenceSymmetric]
    \proofstepwidestar[1]{(x,y)\in\IntRel}{\rA}
    \proofstepwidestar[1]{x\in D_{\mathbb Z}\land y\in D_{\mathbb Z}}{%
      \FormulaRefAuto{A\subseteq B\dsep x\in A\vdash x\in B}{%
        \FormulaRefAuto{IntRelationCarrier},1}}
    \proofstepwidestar[1]{x\in D_{\mathbb Z}}{\rAEa{2}}
    \proofstepwidestar[1]{y\in D_{\mathbb Z}}{\rAEb{2}}
    \proofstepwidestar[1]{\exists a\in\mathbb N\,
      \exists b\in\mathbb N\,x=(a,b)}{%
      \FormulaRefAuto{x\in A\times B\eqvdash
        \exists a\in A\,\exists b\in B\,x=(a,b)}{3}}
    \proofstepwidestar[1]{\exists c\in\mathbb N\,
      \exists d\in\mathbb N\,y=(c,d)}{%
      \FormulaRefAuto{x\in A\times B\eqvdash
        \exists a\in A\,\exists b\in B\,x=(a,b)}{4}}
    \proofstepwidestar[7]{a,b,c,d\in\mathbb{N}\land
      x=(a,b)\land y=(c,d)}{\rA}
    \proofstepwidestar[1,7]{((a,b),(c,d))\in\IntRel}{\rIE{7,1}}
    \proofstepwidestar[1,7]{((c,d),(a,b))\in\IntRel}{%
      \FormulaRefAuto{IntRelationSymmetricPairs}{7,8}}
    \proofstepwidestar[1,7]{(y,x)\in\IntRel}{\rIE{7,9}}
    \proofstepwidestar[1]{(y,x)\in\IntRel}{\rEE{5,6,7,10}}
  \closeproofpartwideindR

  \proofpartwideindR[Transitivität]{%
    (x,y)\in\IntRel\dsep(y,z)\in\IntRel
    \vdash(x,z)\in\IntRel%
  }{%
    (x,y)\in\IntRel\dsep(y,z)\in\IntRel
    \vdash(x,z)\in\IntRel%
  }[IntegerPairEquivalenceTransitive]
    \proofstepwidestar[1]{(x,y)\in\IntRel}{\rA}
    \proofstepwidestar[2]{(y,z)\in\IntRel}{\rA}
    \proofstepwidestar[1]{x\in D_{\mathbb Z}\land y\in D_{\mathbb Z}}{%
      \FormulaRefAuto{A\subseteq B\dsep x\in A\vdash x\in B}{%
        \FormulaRefAuto{IntRelationCarrier},1}}
    \proofstepwidestar[2]{y\in D_{\mathbb Z}\land z\in D_{\mathbb Z}}{%
      \FormulaRefAuto{A\subseteq B\dsep x\in A\vdash x\in B}{%
        \FormulaRefAuto{IntRelationCarrier},2}}
    \proofstepwidestar[1]{x\in D_{\mathbb Z}}{\rAEa{3}}
    \proofstepwidestar[1]{y\in D_{\mathbb Z}}{\rAEb{3}}
    \proofstepwidestar[2]{z\in D_{\mathbb Z}}{\rAEb{4}}
    \proofstepwidestar[1]{\exists a\in\mathbb N\,
      \exists b\in\mathbb N\,x=(a,b)}{%
      \FormulaRefAuto{x\in A\times B\eqvdash
        \exists a\in A\,\exists b\in B\,x=(a,b)}{5}}
    \proofstepwidestar[1]{\exists c\in\mathbb N\,
      \exists d\in\mathbb N\,y=(c,d)}{%
      \FormulaRefAuto{x\in A\times B\eqvdash
        \exists a\in A\,\exists b\in B\,x=(a,b)}{6}}
    \proofstepwidestar[2]{\exists e\in\mathbb N\,
      \exists f\in\mathbb N\,z=(e,f)}{%
      \FormulaRefAuto{x\in A\times B\eqvdash
        \exists a\in A\,\exists b\in B\,x=(a,b)}{7}}
    \proofstepwidestar[11]{a,b,c,d,e,f\in\mathbb{N}\land
      x=(a,b)\land y=(c,d)\land z=(e,f)}{\rA}
    \proofstepwidestar[1,11]{((a,b),(c,d))\in\IntRel}{\rIE{11,1}}
    \proofstepwidestar[2,11]{((c,d),(e,f))\in\IntRel}{\rIE{11,2}}
    \proofstepwidestar[1,2,11]{((a,b),(e,f))\in\IntRel}{%
      \FormulaRefAuto{IntRelationTransitivePairs}{11,12,13}}
    \proofstepwidestar[1,2,11]{(x,z)\in\IntRel}{\rIE{11,14}}
    \proofstepwidestar[1,2]{(x,z)\in\IntRel}{%
      \rEE{8,9,10,11,15}}
  \closeproofpartwideindR

  \proofpartwide{Schluss}
    \proofstepwidestar[]{\IntRel\subseteq
      D_{\mathbb Z}\times D_{\mathbb Z}}{%
      \FormulaRefAuto{IntegerPairEquivalenceCarrier}}
    \proofstepwidestar[]{x\in D_{\mathbb Z}\rightarrow(x,x)\in\IntRel}{%
      \FormulaRefAuto{IntegerPairEquivalenceReflexive}}
    \proofstepwidestar[]{(x,y)\in\IntRel\rightarrow(y,x)\in\IntRel}{%
      \FormulaRefAuto{IntegerPairEquivalenceSymmetric}}
    \proofstepwidestar[]{(x,y)\in\IntRel\dsep(y,z)\in\IntRel
      \vdash(x,z)\in\IntRel}{%
      \FormulaRefAuto{IntegerPairEquivalenceTransitive}}
    \proofstepwidestar[]{\EqRel{\IntRel,D_{\mathbb Z}}}{%
      \FormulaRefAuto{EqRelDef}{1,2,3,4}}
  \closeproofpartwide
\end{tabproofsplitwide}

\section{Quotient und Klassen}

\FormulaDefDeltaK[Menge der ganzen Zahlen]{%
  \mathbb{Z}\coloneqq
  \QuotientSet{D_{\mathbb Z}}{\IntRel}%
}{IntegersAsQuotient}{%
  \DeltaRow{Mengen}{D_{\mathbb Z}\dsep\mathbb Z}%
  \DeltaRow{Relation}{\IntRel}%
  \DeltaRow{Neue Symbole}{\mathbb Z}%
}

\FormulaDefDeltaK[Klasse eines Differenzenpaares]{%
  a,b\in\mathbb{N}\vdash
  \IntClass{a}{b}\coloneqq
  \EqClass{(a,b)}{\IntRel}{D_{\mathbb Z}}%
}{IntegerClassDef}{%
  \DeltaRow{Mengen}{a\dsep b\dsep D_{\mathbb Z}}%
  \DeltaRow{Relation}{\IntRel}%
  \DeltaRow{Neue Symbole}{\IntClass{a}{b}}%
}

\FormulaThmDeltaK[Klassen sind ganze Zahlen]{%
  a,b\in\mathbb{N}\vdash\IntClass{a}{b}\in\mathbb{Z}%
}{IntegerClassInZ}{%
  \DeltaRow{Mengen}{a\dsep b}%
}
\begin{tabproof}
  \proofstep{1}{a,b\in\mathbb{N}}{\rA}
  \proofstep{1}{(a,b)\in D_{\mathbb Z}}{%
    \FormulaRefAuto{(a,b)\in A\times B\eqvdash a\in A\land b\in B}{1}}
  \proofstep{}{\EqRel{\IntRel,D_{\mathbb Z}}}{%
    \FormulaRefAuto{IntegerPairEquivalence}}
  \proofstep{1}{\EqClass{(a,b)}{\IntRel}{D_{\mathbb Z}}
    \in\QuotientSet{D_{\mathbb Z}}{\IntRel}}{%
    \FormulaRefAuto{EqClassInQuotient}{3,2}}
  \proofstep{1}{\IntClass{a}{b}=
    \EqClass{(a,b)}{\IntRel}{D_{\mathbb Z}}}{%
    \FormulaRefAuto{IntegerClassDef}{1}}
  \proofstep{}{\mathbb Z=
    \QuotientSet{D_{\mathbb Z}}{\IntRel}}{%
    \FormulaRefAuto{IntegersAsQuotient}}
  \proofstep{1}{\IntClass{a}{b}\in\mathbb Z}{\rIE{5,6,4}}
\end{tabproof}

\FormulaThmDeltaK[Gleichheit von Ganzzahlklassen]{%
  a,b,c,d\in\mathbb{N}\vdash
  \IntClass{a}{b}=\IntClass{c}{d}
  \leftrightarrow a+d=b+c%
}{IntegerClassEquality}{%
  \DeltaRow{Mengen}{a\dsep b\dsep c\dsep d}%
  \DeltaRow{Relation}{\IntRel}%
}
\begin{tabproofsplitwide}
  \proofpartwideindR[Klassenkriterium]{%
    \begin{aligned}[t]
      &a,b,c,d\in\mathbb{N}\vdash{}\\
      &\IntClass{a}{b}=\IntClass{c}{d}
      \leftrightarrow ((a,b),(c,d))\in\IntRel
    \end{aligned}
  }{%
    a,b,c,d\in\mathbb{N}\vdash
    \IntClass{a}{b}=\IntClass{c}{d}
    \leftrightarrow ((a,b),(c,d))\in\IntRel
  }[IntegerClassEqualityViaRelation]
    \proofstepwidestar[1]{a,b,c,d\in\mathbb{N}}{\rA}
    \proofstepwidestar[]{\EqRel{\IntRel,D_{\mathbb Z}}}{%
      \FormulaRefAuto{IntegerPairEquivalence}}
    \proofstepwidestar[1]{(a,b)\in D_{\mathbb Z}}{%
      \FormulaRefAuto{(a,b)\in A\times B\eqvdash a\in A\land b\in B}{1}}
    \proofstepwidestar[1]{(c,d)\in D_{\mathbb Z}}{%
      \FormulaRefAuto{(a,b)\in A\times B\eqvdash a\in A\land b\in B}{1}}
    \proofstepwidestar[1]{%
      \EqClass{(a,b)}{\IntRel}{D_{\mathbb Z}}
      =\EqClass{(c,d)}{\IntRel}{D_{\mathbb Z}}
      \leftrightarrow ((a,b),(c,d))\in\IntRel}{%
      \FormulaRefAuto{EqClassEqualIffRel}{2,3,4}}
    \proofstepwidestar[1]{\IntClass{a}{b}=
      \EqClass{(a,b)}{\IntRel}{D_{\mathbb Z}}}{%
      \FormulaRefAuto{IntegerClassDef}{1}}
    \proofstepwidestar[1]{\IntClass{c}{d}=
      \EqClass{(c,d)}{\IntRel}{D_{\mathbb Z}}}{%
      \FormulaRefAuto{IntegerClassDef}{1}}
    \proofstepwidestar[1]{%
      \IntClass{a}{b}=\IntClass{c}{d}
      \leftrightarrow ((a,b),(c,d))\in\IntRel}{\rIE{5,6,7}}
  \closeproofpartwideindR

  \proofpartwide{Schluss}
    \proofstepwidestar[1]{a,b,c,d\in\mathbb{N}}{\rA}
    \proofstepwidestar[1]{%
      \IntClass{a}{b}=\IntClass{c}{d}
      \leftrightarrow ((a,b),(c,d))\in\IntRel}{%
      \FormulaRefAuto{IntegerClassEqualityViaRelation}{1}}
    \proofstepwidestar[1]{%
      ((a,b),(c,d))\in\IntRel\leftrightarrow a+d=b+c}{%
      \FormulaRefAuto{IntRelationChar}{1}}
    \proofstepwidestar[1]{%
      \IntClass{a}{b}=\IntClass{c}{d}
      \leftrightarrow a+d=b+c}{%
      \FormulaRefAuto{P\leftrightarrow Q\dsep
        Q\leftrightarrow R\vdash P\leftrightarrow R}{2,3}}
  \closeproofpartwide
\end{tabproofsplitwide}

\FormulaThmDeltaK[Jede ganze Zahl besitzt ein Differenzenpaar]{%
  z\in\mathbb Z\vdash
  \exists a\in\mathbb N\,\exists b\in\mathbb N\,
  z=\IntClass{a}{b}%
}{IntegerPairRepresentative}{%
  \DeltaRow{Mengen}{z\dsep a\dsep b\dsep D_{\mathbb Z}}%
}
\begin{tabproofsplitwide}
  \proofpartwideindR[Quotientenrepräsentant]{%
    z\in\mathbb Z\vdash
    \exists x\in D_{\mathbb Z}\,
    z=\EqClass{x}{\IntRel}{D_{\mathbb Z}}%
  }{%
    z\in\mathbb Z\vdash
    \exists x\in D_{\mathbb Z}\,
    z=\EqClass{x}{\IntRel}{D_{\mathbb Z}}%
  }[IntegerQuotientRepresentative]
    \proofstepwidestar[1]{z\in\mathbb Z}{\rA}
    \proofstepwidestar[]{\mathbb Z=
      \QuotientSet{D_{\mathbb Z}}{\IntRel}}{%
      \FormulaRefAuto{IntegersAsQuotient}}
    \proofstepwidestar[1]{z\in
      \QuotientSet{D_{\mathbb Z}}{\IntRel}}{\rIE{2,1}}
    \proofstepwidestar[]{\EqRel{\IntRel,D_{\mathbb Z}}}{%
      \FormulaRefAuto{IntegerPairEquivalence}}
    \proofstepwidestar[1]{\exists x\in D_{\mathbb Z}\,
      z=\EqClass{x}{\IntRel}{D_{\mathbb Z}}}{%
      \FormulaRefAuto{QuotientRepresentative}{4,3}}
  \closeproofpartwideindR

  \proofpartwide{Schluss}
    \proofstepwidestar[1]{z\in\mathbb Z}{\rA}
    \proofstepwidestar[1]{\exists x\in D_{\mathbb Z}\,
      z=\EqClass{x}{\IntRel}{D_{\mathbb Z}}}{%
      \FormulaRefAuto{IntegerQuotientRepresentative}{1}}
    \proofstepwidestar[3]{x\in D_{\mathbb Z}\land
      z=\EqClass{x}{\IntRel}{D_{\mathbb Z}}}{\rA}
    \proofstepwidestar[3]{\exists a\in\mathbb N\,
      \exists b\in\mathbb N\,x=(a,b)}{%
      \FormulaRefAuto{x\in A\times B\eqvdash
        \exists a\in A\,\exists b\in B\,x=(a,b)}{\rAEa{3}}}
    \proofstepwidestar[5]{a,b\in\mathbb N\land x=(a,b)}{\rA}
    \proofstepwidestar[3,5]{z=\EqClass{(a,b)}{\IntRel}{D_{\mathbb Z}}}{%
      \rIE{5,\rAEb{3}}}
    \proofstepwidestar[5]{\IntClass{a}{b}=
      \EqClass{(a,b)}{\IntRel}{D_{\mathbb Z}}}{%
      \FormulaRefAuto{IntegerClassDef}{5}}
    \proofstepwidestar[3,5]{z=\IntClass{a}{b}}{%
      \FormulaRefAuto{a=b\dsep c=b\vdash a=c}{6,7}}
    \proofstepwidestar[3]{\exists a\in\mathbb N\,
      \exists b\in\mathbb N\,z=\IntClass{a}{b}}{\rEE{4,5,8}}
    \proofstepwidestar[1]{\exists a\in\mathbb N\,
      \exists b\in\mathbb N\,z=\IntClass{a}{b}}{\rEE{2,3,9}}
  \closeproofpartwide
\end{tabproofsplitwide}

\section{Einbettung der natürlichen Zahlen}

\FormulaDefDeltaK[Funktionsvorschrift der natürlichen Einbettung]{%
  n\in\mathbb N\vdash
  \IntEmbed(n)\coloneqq\IntClass{n}{0}%
}{NaturalIntegerEmbedding}{%
  \DeltaRow{Mengen}{n}%
  \DeltaRow{Funktionensymbol}{\IntEmbed}%
  \DeltaRow{Neue Symbole}{\IntEmbed}%
}

\FormulaThmDeltaK[Typisierung der natürlichen Einbettung]{%
  \IntEmbed\colon\mathbb N\to\mathbb Z%
}{NaturalIntegerEmbeddingFunction}{%
  \DeltaRow{Funktion}{\IntEmbed}%
}
\begin{tabproof}
  \proofstep{1}{n\in\mathbb N}{\rA}
  \proofstep{}{0\in\mathbb N}{\FormulaRefAuto{PeanoZero}[ax]}
  \proofstep{1}{\IntClass{n}{0}\in\mathbb Z}{%
    \FormulaRefAuto{IntegerClassInZ}{1,2}}
  \proofstep{1}{\IntEmbed(n)=\IntClass{n}{0}}{%
    \FormulaRefAuto{NaturalIntegerEmbedding}{1}}
  \proofstep{1}{\IntEmbed(n)\in\mathbb Z}{\rIE{4,3}}
  \proofstep{}{\IntEmbed\colon\mathbb N\to\mathbb Z}{%
    \FormulaRefAuto{x\in A\vdash t(x)\coloneq y\dsep
      x\in A\vdash t(x)\in B
      \exists! F\colon A\to B\,\forall x\in A\,F(x)=t(x)}{1,5}}
\end{tabproof}

\FormulaThmDeltaK[Injektivität der natürlichen Einbettung]{%
  \IntEmbed\colon\mathbb N\inj\mathbb Z%
}{NaturalIntegerEmbeddingInjective}{%
  \DeltaRow{Mengen}{m\dsep n}%
  \DeltaRow{Funktion}{\IntEmbed}%
}
\begin{tabproofsplitwide}
  \proofpartwideindR[Gleichheit von eingebetteten Zahlen]{%
    m,n\in\mathbb N\dsep
    \IntEmbed(m)=\IntEmbed(n)\vdash m=n%
  }{%
    m,n\in\mathbb N\dsep
    \IntEmbed(m)=\IntEmbed(n)\vdash m=n%
  }[NaturalIntegerEmbeddingReflectsEquality]
    \proofstepwidestar[1]{m,n\in\mathbb N}{\rA}
    \proofstepwidestar[2]{\IntEmbed(m)=\IntEmbed(n)}{\rA}
    \proofstepwidestar[1]{\IntEmbed(m)=\IntClass{m}{0}}{%
      \FormulaRefAuto{NaturalIntegerEmbedding}{1}}
    \proofstepwidestar[1]{\IntEmbed(n)=\IntClass{n}{0}}{%
      \FormulaRefAuto{NaturalIntegerEmbedding}{1}}
    \proofstepwidestar[1,2]{\IntClass{m}{0}=\IntClass{n}{0}}{\rIE{3,4,2}}
    \proofstepwidestar[1,2]{m+0=0+n}{%
      \FormulaRefAuto{IntegerClassEquality}{1,5}}
    \proofstepwidestar[1]{m+0=m}{\FormulaRefAuto{PeanoAddRightZero}{1}}
    \proofstepwidestar[1]{0+n=n}{\FormulaRefAuto{PeanoAddLeftZero}{1}}
    \proofstepwidestar[1,2]{m=n}{\rIE{7,8,6}}
  \closeproofpartwideindR

  \proofpartwide{Schluss}
    \proofstepwidestar[]{\IntEmbed\colon\mathbb N\to\mathbb Z}{%
      \FormulaRefAuto{NaturalIntegerEmbeddingFunction}}
    \proofstepwidestar[2]{m\in\mathbb N}{\rA}
    \proofstepwidestar[3]{n\in\mathbb N}{\rA}
    \proofstepwidestar[4]{\IntEmbed(m)=\IntEmbed(n)}{\rA}
    \proofstepwidestar[2,3,4]{m=n}{%
      \FormulaRefAuto{NaturalIntegerEmbeddingReflectsEquality}{2,3,4}}
    \proofstepwidestar[]{\IntEmbed\colon\mathbb N\inj\mathbb Z}{%
      \FormulaRefAuto{Injektive Funktion}{1,\rUI{\rRI{2,3,4,5}}}}
  \closeproofpartwide
\end{tabproofsplitwide}

\FormulaDefDeltaK[Null der ganzen Zahlen]{%
  0_{\mathbb Z}\coloneqq\IntEmbed(0)%
}{IntegerZeroDef}{%
  \DeltaRow{Mengen}{0_{\mathbb Z}}%
  \DeltaRow{Neue Symbole}{0_{\mathbb Z}}%
}

\FormulaDefDeltaK[Eins der ganzen Zahlen]{%
  1_{\mathbb Z}\coloneqq\IntEmbed(1)%
}{IntegerOneDef}{%
  \DeltaRow{Mengen}{1_{\mathbb Z}}%
  \DeltaRow{Neue Symbole}{1_{\mathbb Z}}%
}

\FormulaThmDeltaK[Nullklasse]{%
  0_{\mathbb Z}=\IntClass{0}{0}%
}{IntegerZeroClass}{%
  \DeltaRow{Mengen}{0_{\mathbb Z}}%
}
\begin{tabproof}
  \proofstep{}{0_{\mathbb Z}=\IntEmbed(0)}{\FormulaRefAuto{IntegerZeroDef}}
  \proofstep{}{0\in\mathbb N}{\FormulaRefAuto{PeanoZero}[ax]}
  \proofstep{}{\IntEmbed(0)=\IntClass{0}{0}}{%
    \FormulaRefAuto{NaturalIntegerEmbedding}{2}}
  \proofstep{}{0_{\mathbb Z}=\IntClass{0}{0}}{%
    \FormulaRefAuto{a=b\dsep b=c\vdash a=c}{1,3}}
\end{tabproof}

\FormulaThmDeltaK[Die Ganzzahlnull liegt im Quotienten]{%
  0_{\mathbb Z}\in\mathbb Z%
}{IntZeroCarrier}{%
  \DeltaRow{Mengen}{0_{\mathbb Z}}%
}
\begin{tabproof}
  \proofstep{}{0_{\mathbb Z}=\IntClass{0}{0}}{%
    \FormulaRefAuto{IntegerZeroClass}}
  \proofstep{}{0\in\mathbb N}{\FormulaRefAuto{PeanoZero}[ax]}
  \proofstep{}{\IntClass{0}{0}\in\mathbb Z}{%
    \FormulaRefAuto{IntegerClassInZ}{2,2}}
  \proofstep{}{0_{\mathbb Z}\in\mathbb Z}{\rIE{1,3}}
\end{tabproof}

\FormulaThmDeltaK[Die Ganzzahleins liegt im Quotienten]{%
  1_{\mathbb Z}\in\mathbb Z%
}{IntOneCarrier}{%
  \DeltaRow{Mengen}{1_{\mathbb Z}}%
}
\begin{tabproof}
  \proofstep{}{1_{\mathbb Z}=\IntEmbed(1)}{%
    \FormulaRefAuto{IntegerOneDef}}
  \proofstep{}{1\in\mathbb N}{\FormulaRefAuto{PeanoOneInN}}
  \proofstep{}{\IntEmbed\colon\mathbb N\to\mathbb Z}{%
    \FormulaRefAuto{NaturalIntegerEmbeddingFunction}}
  \proofstep{}{\IntEmbed(1)\in\mathbb Z}{%
    \FormulaRefAuto{F\colon A\to B\dsep x\in A\vdash F(x)\in B}{3,2}}
  \proofstep{}{1_{\mathbb Z}\in\mathbb Z}{\rIE{1,4}}
\end{tabproof}

\section{Negation und Addition}

\subsection{Negation}

\FormulaThmDeltaK[Repräsentantenunabhängigkeit der Negation]{%
  \begin{aligned}[t]
    &a,b,c,d\in\mathbb N\dsep
      \IntClass{a}{b}=\IntClass{c}{d}\vdash{}\\
    &\IntClass{b}{a}=\IntClass{d}{c}
  \end{aligned}%
}{IntegerNegationCompatible}{%
  \DeltaRow{Mengen}{a\dsep b\dsep c\dsep d}%
}
\begin{tabproofsplitwide}
  \proofpartwideindR[Kreuzsummen des Ausgangspaars]{%
    a,b,c,d\in\mathbb N\dsep
    \IntClass{a}{b}=\IntClass{c}{d}
    \vdash a+d=b+c%
  }{%
    a,b,c,d\in\mathbb N\dsep
    \IntClass{a}{b}=\IntClass{c}{d}
    \vdash a+d=b+c%
  }[IntegerNegationCompatibleCrossSum]
    \proofstepwidestar[1]{a,b,c,d\in\mathbb N}{\rA}
    \proofstepwidestar[2]{\IntClass{a}{b}=\IntClass{c}{d}}{\rA}
    \proofstepwidestar[1,2]{a+d=b+c}{%
      \FormulaRefAuto{IntegerClassEquality}{1,2}}
  \closeproofpartwideindR

  \proofpartwide{Vertauschen der Koordinaten}
    \proofstepwidestar[1]{a,b,c,d\in\mathbb N}{\rA}
    \proofstepwidestar[2]{\IntClass{a}{b}=\IntClass{c}{d}}{\rA}
    \proofstepwidestar[1,2]{a+d=b+c}{%
      \FormulaRefAuto{IntegerNegationCompatibleCrossSum}{1,2}}
    \proofstepwidestar[1,2]{b+c=a+d}{%
      \FormulaRefAuto{a=b\vdash b=a}{3}}
    \proofstepwidestar[1,2]{\IntClass{b}{a}=\IntClass{d}{c}}{%
      \FormulaRefAuto{IntegerClassEquality}{1,4}}
  \closeproofpartwide
\end{tabproofsplitwide}

\FormulaThmDeltaK[Quotientenabstieg der Negation]{%
  \begin{aligned}[t]
    z\in\mathbb Z\vdash
    \exists! u\in\mathbb Z\,
    \exists p\in\mathbb N\,\exists q\in\mathbb N\,
    \bigl(z=\IntClass{p}{q}\land u=\IntClass{q}{p}\bigr)
  \end{aligned}%
}{IntegerNegationQuotientDescent}{%
  \DeltaRow{Mengen}{z\dsep u\dsep v\dsep
    a\dsep b\dsep c\dsep d\dsep p\dsep q}%
}
\begin{tabproofsplitwide}
  \proofpartwideindR[Existenz aus einem Repräsentanten]{%
    \begin{aligned}[t]
      z\in\mathbb Z\vdash
      \exists u\in\mathbb Z\,
      \exists p\in\mathbb N\,\exists q\in\mathbb N\,
      \bigl(z=\IntClass{p}{q}\land u=\IntClass{q}{p}\bigr)
    \end{aligned}%
  }{%
    z\in\mathbb Z\vdash
    \exists u\in\mathbb Z\,
    \exists p\in\mathbb N\,\exists q\in\mathbb N\,
    \bigl(z=\IntClass{p}{q}\land u=\IntClass{q}{p}\bigr)%
  }[IntegerNegationDescentExistence]
    \proofstepwidestar[1]{z\in\mathbb Z}{\rA}
    \proofstepwidestar[1]{%
      \exists a\in\mathbb N\,\exists b\in\mathbb N\,
      z=\IntClass{a}{b}}{%
      \FormulaRefAuto{IntegerPairRepresentative}{1}}
    \proofstepwidestar[3]{%
      a,b\in\mathbb N\land z=\IntClass{a}{b}}{\rA}
    \proofstepwidestar[3]{\IntClass{b}{a}\in\mathbb Z}{%
      \FormulaRefAuto{IntegerClassInZ}{3}}
    \proofstepwidestar[]{%
      \IntClass{b}{a}=\IntClass{b}{a}}{\rII}
    \proofstepwidestar[3]{%
      \exists p\in\mathbb N\,\exists q\in\mathbb N\,
      \bigl(z=\IntClass{p}{q}\land
      \IntClass{b}{a}=\IntClass{q}{p}\bigr)}{%
      \rEI{\rAI{3,5}}}
    \proofstepwidestar[3]{%
      \IntClass{b}{a}\in\mathbb Z\land
      \exists p\in\mathbb N\,\exists q\in\mathbb N\,
      \bigl(z=\IntClass{p}{q}\land
      \IntClass{b}{a}=\IntClass{q}{p}\bigr)}{%
      \rAI{4,6}}
    \proofstepwidestar[3]{%
      \exists u\in\mathbb Z\,
      \exists p\in\mathbb N\,\exists q\in\mathbb N\,
      \bigl(z=\IntClass{p}{q}\land u=\IntClass{q}{p}\bigr)}{%
      \rEI{7}}
    \proofstepwidestar[1]{%
      \exists u\in\mathbb Z\,
      \exists p\in\mathbb N\,\exists q\in\mathbb N\,
      \bigl(z=\IntClass{p}{q}\land u=\IntClass{q}{p}\bigr)}{%
      \rEE{2,3,8}}
  \closeproofpartwideindR

  \proofpartwideindR[Eindeutigkeit für feste Repräsentanten]{%
    \begin{aligned}[t]
      &a,b,c,d\in\mathbb N\dsep
      z=\IntClass{a}{b}\dsep u=\IntClass{b}{a}\dsep{}\\
      &z=\IntClass{c}{d}\dsep v=\IntClass{d}{c}
      \vdash u=v
    \end{aligned}%
  }{%
    a,b,c,d\in\mathbb N\dsep
    z=\IntClass{a}{b}\dsep u=\IntClass{b}{a}\dsep
    z=\IntClass{c}{d}\dsep v=\IntClass{d}{c}
    \vdash u=v%
  }[IntegerNegationDescentFixedUniqueness]
    \proofstepwidestar[1]{a,b,c,d\in\mathbb N}{\rA}
    \proofstepwidestar[2]{z=\IntClass{a}{b}}{\rA}
    \proofstepwidestar[3]{u=\IntClass{b}{a}}{\rA}
    \proofstepwidestar[4]{z=\IntClass{c}{d}}{\rA}
    \proofstepwidestar[5]{v=\IntClass{d}{c}}{\rA}
    \proofstepwidestar[2,4]{%
      \IntClass{a}{b}=\IntClass{c}{d}}{\rIE{2,4}}
    \proofstepwidestar[1,2,4]{%
      \IntClass{b}{a}=\IntClass{d}{c}}{%
      \FormulaRefAuto{IntegerNegationCompatible}{1,6}}
    \proofstepwidestar[1,2,3,4,5]{u=v}{\rIE{3,5,7}}
  \closeproofpartwideindR

  \proofpartwideindR[Eindeutigkeit beliebiger Ergebniszeugen]{%
    \begin{aligned}[t]
      &\bigl(u\in\mathbb Z\land
      \exists a\in\mathbb N\,\exists b\in\mathbb N\,
        (z=\IntClass{a}{b}\land u=\IntClass{b}{a})\bigr)\dsep{}\\
      &\bigl(v\in\mathbb Z\land
      \exists c\in\mathbb N\,\exists d\in\mathbb N\,
        (z=\IntClass{c}{d}\land v=\IntClass{d}{c})\bigr)
      \vdash u=v
    \end{aligned}%
  }{%
    \bigl(u\in\mathbb Z\land
    \exists a\in\mathbb N\,\exists b\in\mathbb N\,
      (z=\IntClass{a}{b}\land u=\IntClass{b}{a})\bigr)\dsep
    \bigl(v\in\mathbb Z\land
    \exists c\in\mathbb N\,\exists d\in\mathbb N\,
      (z=\IntClass{c}{d}\land v=\IntClass{d}{c})\bigr)
    \vdash u=v%
  }[IntegerNegationDescentUniqueness]
    \proofstepwidestar[1]{%
      u\in\mathbb Z\land
      \exists a\in\mathbb N\,\exists b\in\mathbb N\,
      (z=\IntClass{a}{b}\land u=\IntClass{b}{a})}{\rA}
    \proofstepwidestar[2]{%
      v\in\mathbb Z\land
      \exists c\in\mathbb N\,\exists d\in\mathbb N\,
      (z=\IntClass{c}{d}\land v=\IntClass{d}{c})}{\rA}
    \proofstepwidestar[1]{%
      \exists a\in\mathbb N\,\exists b\in\mathbb N\,
      (z=\IntClass{a}{b}\land u=\IntClass{b}{a})}{\rAEb{1}}
    \proofstepwidestar[2]{%
      \exists c\in\mathbb N\,\exists d\in\mathbb N\,
      (z=\IntClass{c}{d}\land v=\IntClass{d}{c})}{\rAEb{2}}
    \proofstepwidestar[5]{%
      a,b\in\mathbb N\land
      z=\IntClass{a}{b}\land u=\IntClass{b}{a}}{\rA}
    \proofstepwidestar[6]{%
      c,d\in\mathbb N\land
      z=\IntClass{c}{d}\land v=\IntClass{d}{c}}{\rA}
    \proofstepwidestar[5]{a,b\in\mathbb N}{\rAEa{5}}
    \proofstepwidestar[5]{z=\IntClass{a}{b}}{\rAEa{\rAEb{5}}}
    \proofstepwidestar[5]{u=\IntClass{b}{a}}{\rAEb{\rAEb{5}}}
    \proofstepwidestar[6]{c,d\in\mathbb N}{\rAEa{6}}
    \proofstepwidestar[6]{z=\IntClass{c}{d}}{\rAEa{\rAEb{6}}}
    \proofstepwidestar[6]{v=\IntClass{d}{c}}{\rAEb{\rAEb{6}}}
    \proofstepwidestar[5,6]{a,b,c,d\in\mathbb N}{\rAI{7,10}}
    \proofstepwidestar[5,6]{u=v}{%
      \FormulaRefAuto{IntegerNegationDescentFixedUniqueness}{%
        13,8,9,11,12}}
    \proofstepwidestar[5]{u=v}{\rEE{4,6,14}}
    \proofstepwidestar[1,2]{u=v}{\rEE{3,5,15}}
  \closeproofpartwideindR

  \proofpartwideindR[Abschluss des Abstiegs]{%
    \begin{aligned}[t]
      z\in\mathbb Z\vdash
      \exists! u\in\mathbb Z\,
      \exists p\in\mathbb N\,\exists q\in\mathbb N\,
      \bigl(z=\IntClass{p}{q}\land u=\IntClass{q}{p}\bigr)
    \end{aligned}%
  }{%
    z\in\mathbb Z\vdash
    \exists! u\in\mathbb Z\,
    \exists p\in\mathbb N\,\exists q\in\mathbb N\,
    \bigl(z=\IntClass{p}{q}\land u=\IntClass{q}{p}\bigr)%
  }[IntegerNegationDescentConclusion]
    \proofstepwidestar[1]{z\in\mathbb Z}{\rA}
    \proofstepwidestar[1]{%
      \exists u\in\mathbb Z\,
      \exists p\in\mathbb N\,\exists q\in\mathbb N\,
      (z=\IntClass{p}{q}\land u=\IntClass{q}{p})}{%
      \FormulaRefAuto{IntegerNegationDescentExistence}{1}}
    \proofstepwidestar[3]{%
      u\in\mathbb Z\land
      \exists p\in\mathbb N\,\exists q\in\mathbb N\,
      (z=\IntClass{p}{q}\land u=\IntClass{q}{p})}{\rA}
    \proofstepwidestar[4]{%
      v\in\mathbb Z\land
      \exists p\in\mathbb N\,\exists q\in\mathbb N\,
      (z=\IntClass{p}{q}\land v=\IntClass{q}{p})}{\rA}
    \proofstepwidestar[3,4]{u=v}{%
      \FormulaRefAuto{IntegerNegationDescentUniqueness}{3,4}}
    \proofstepwidestar[1]{%
      \exists! u\in\mathbb Z\,
      \exists p\in\mathbb N\,\exists q\in\mathbb N\,
      (z=\IntClass{p}{q}\land u=\IntClass{q}{p})}{%
      \UEI{2,3,4,5}}
  \closeproofpartwideindR
\end{tabproofsplitwide}

Der vorangehende Satz zeigt, dass die folgende Klassenregel unabhängig
vom gewählten Differenzenpaar genau ein Ergebnis in \(\mathbb Z\) festlegt.

\FormulaDefDeltaK[Negation ganzer Zahlen]{%
  a,b\in\mathbb N\vdash
  -\IntClass{a}{b}\coloneqq\IntClass{b}{a}%
}{IntegerNegationDef}{%
  \DeltaRow{Mengen}{a\dsep b}%
  \DeltaRow{Neue Symbole}{-}%
}

\FormulaThmDeltaK[Abgeschlossenheit der Negation]{%
  z\in\mathbb Z\vdash -z\in\mathbb Z%
}{IntNegClosure}{%
  \DeltaRow{Mengen}{z\dsep a\dsep b}%
}
\begin{tabproof}
  \proofstep{1}{z\in\mathbb Z}{\rA}
  \proofstep{1}{\exists a\in\mathbb N\,\exists b\in\mathbb N\,
    z=\IntClass{a}{b}}{\FormulaRefAuto{IntegerPairRepresentative}{1}}
  \proofstep{3}{a,b\in\mathbb N\land z=\IntClass{a}{b}}{\rA}
  \proofstep{3}{-\IntClass{a}{b}=\IntClass{b}{a}}{%
    \FormulaRefAuto{IntegerNegationDef}{3}}
  \proofstep{3}{\IntClass{b}{a}\in\mathbb Z}{%
    \FormulaRefAuto{IntegerClassInZ}{3}}
  \proofstep{3}{-z\in\mathbb Z}{\rIE{3,4,5}}
  \proofstep{1}{-z\in\mathbb Z}{\rEE{2,3,6}}
\end{tabproof}

\FormulaThmDeltaK[Doppelte Negation]{%
  z\in\mathbb Z\vdash -(-z)=z%
}{IntegerDoubleNegation}{%
  \DeltaRow{Mengen}{z\dsep a\dsep b}%
}
\begin{tabproof}
  \proofstep{1}{z\in\mathbb Z}{\rA}
  \proofstep{1}{\exists a,b\in\mathbb N\,
    z=\IntClass{a}{b}}{\FormulaRefAuto{IntegerPairRepresentative}{1}}
  \proofstep{3}{a,b\in\mathbb N\land z=\IntClass{a}{b}}{\rA}
  \proofstep{3}{-\IntClass{a}{b}=\IntClass{b}{a}}{%
    \FormulaRefAuto{IntegerNegationDef}{3}}
  \proofstep{3}{-\IntClass{b}{a}=\IntClass{a}{b}}{%
    \FormulaRefAuto{IntegerNegationDef}{3}}
  \proofstep{3}{-(-z)=z}{\rIE{3,4,5}}
  \proofstep{1}{-(-z)=z}{\rEE{2,3,6}}
\end{tabproof}

\subsection{Addition}

\FormulaThmDeltaK[Repräsentantenunabhängigkeit der Addition]{%
  \begin{aligned}[t]
    &a,b,a',b',c,d,c',d'\in\mathbb N\dsep
      \IntClass{a}{b}=\IntClass{a'}{b'}\dsep
      \IntClass{c}{d}=\IntClass{c'}{d'}\vdash{}\\
    &\IntClass{a+c}{b+d}=\IntClass{a'+c'}{b'+d'}
  \end{aligned}%
}{IntegerAdditionCompatible}{%
  \DeltaRow{Mengen}{a\dsep b\dsep a'\dsep b'\dsep
    c\dsep d\dsep c'\dsep d'}%
}
\begin{tabproofsplitwide}
  \proofpartwideindR[Kreuzsummen der beiden Eingaben]{%
    \begin{aligned}[t]
      &a,b,a',b',c,d,c',d'\in\mathbb N\dsep
      \IntClass{a}{b}=\IntClass{a'}{b'}\dsep
      \IntClass{c}{d}=\IntClass{c'}{d'}\vdash{}\\
      &a+b'=b+a'\quad\land\quad c+d'=d+c'
    \end{aligned}
  }{%
    a,b,a',b',c,d,c',d'\in\mathbb N\dsep
    \IntClass{a}{b}=\IntClass{a'}{b'}\dsep
    \IntClass{c}{d}=\IntClass{c'}{d'}
    \vdash a+b'=b+a'\land c+d'=d+c'
  }[IntegerAdditionInputCrossSums]
    \proofstepwidestar[1]{a,b,a',b',c,d,c',d'\in\mathbb N}{\rA}
    \proofstepwidestar[2]{\IntClass{a}{b}=\IntClass{a'}{b'}}{\rA}
    \proofstepwidestar[3]{\IntClass{c}{d}=\IntClass{c'}{d'}}{\rA}
    \proofstepwidestar[1,2]{a+b'=b+a'}{%
      \FormulaRefAuto{IntegerClassEquality}{1,2}}
    \proofstepwidestar[1,3]{c+d'=d+c'}{%
      \FormulaRefAuto{IntegerClassEquality}{1,3}}
    \proofstepwidestar[1,2,3]{a+b'=b+a'\land c+d'=d+c'}{\rAI{4,5}}
  \closeproofpartwideindR

  \proofpartwideindR[Zusammenführen und Umsortieren]{%
    \begin{aligned}[t]
      &a,b,a',b',c,d,c',d'\in\mathbb N\dsep
      a+b'=b+a'\dsep c+d'=d+c'\vdash{}\\
      &(a+c)+(b'+d')=(b+d)+(a'+c')
    \end{aligned}
  }{%
    a,b,a',b',c,d,c',d'\in\mathbb N\dsep
    a+b'=b+a'\dsep c+d'=d+c'
    \vdash (a+c)+(b'+d')=(b+d)+(a'+c')
  }[IntegerAdditionOutputCrossSum]
    \proofstepwidestar[1]{a,b,a',b',c,d,c',d'\in\mathbb N}{\rA}
    \proofstepwidestar[2]{a+b'=b+a'}{\rA}
    \proofstepwidestar[3]{c+d'=d+c'}{\rA}
    \proofstepwide[1,2,3]{(a+b')+(c+d')}{=}{(b+a')+(d+c')}{\rIE{2,3}}
    \proofstepwide[1]{(a+b')+(c+d')}{=}{(a+c)+(b'+d')}{%
      \FormulaRefAuto{PeanoAddInterchange}{1}}
    \proofstepwide[1]{(b+a')+(d+c')}{=}{(b+d)+(a'+c')}{%
      \FormulaRefAuto{PeanoAddInterchange}{1}}
    \proofstepwide[1,2,3]{(a+c)+(b'+d')}{=}{(b+d)+(a'+c')}{%
      \rIE{4,5,6}}
  \closeproofpartwideindR

  \proofpartwide{Schluss}
    \proofstepwidestar[1]{a,b,a',b',c,d,c',d'\in\mathbb N}{\rA}
    \proofstepwidestar[2]{\IntClass{a}{b}=\IntClass{a'}{b'}}{\rA}
    \proofstepwidestar[3]{\IntClass{c}{d}=\IntClass{c'}{d'}}{\rA}
    \proofstepwidestar[1,2,3]{a+b'=b+a'\land c+d'=d+c'}{%
      \FormulaRefAuto{IntegerAdditionInputCrossSums}{1,2,3}}
    \proofstepwidestar[1,2,3]{(a+c)+(b'+d')=(b+d)+(a'+c')}{%
      \FormulaRefAuto{IntegerAdditionOutputCrossSum}{1,\rAEa{4},\rAEb{4}}}
    \proofstepwidestar[1,2,3]{%
      \IntClass{a+c}{b+d}=\IntClass{a'+c'}{b'+d'}}{%
      \FormulaRefAuto{IntegerClassEquality}{1,5}}
  \closeproofpartwide
\end{tabproofsplitwide}

\FormulaThmDeltaK[Quotientenabstieg der Addition]{%
  \begin{aligned}[t]
    &z,w\in\mathbb Z\vdash
    \exists! u\in\mathbb Z\,
    \exists p\in\mathbb N\,\exists q\in\mathbb N\,
    \exists r\in\mathbb N\,\exists s\in\mathbb N\,{}\\
    &\qquad
    \bigl(z=\IntClass{p}{q}\land
      w=\IntClass{r}{s}\land
      u=\IntClass{p+r}{q+s}\bigr)
  \end{aligned}%
}{IntegerAdditionQuotientDescent}{%
  \DeltaRow{Mengen}{z\dsep w\dsep u\dsep v\dsep
    a\dsep b\dsep c\dsep d\dsep
    a'\dsep b'\dsep c'\dsep d'\dsep
    p\dsep q\dsep r\dsep s}%
}
\begin{tabproofsplitwide}
  \proofpartwideindR[Ergebniskandidat fester Repräsentanten]{%
    \begin{aligned}[t]
      &a,b,c,d\in\mathbb N\dsep
      z=\IntClass{a}{b}\dsep w=\IntClass{c}{d}\vdash{}\\
      &\IntClass{a+c}{b+d}\in\mathbb Z\land
      \exists p\in\mathbb N\,\exists q\in\mathbb N\,
      \exists r\in\mathbb N\,\exists s\in\mathbb N\,{}\\
      &\quad
      \bigl(z=\IntClass{p}{q}\land
        w=\IntClass{r}{s}\land
        \IntClass{a+c}{b+d}=\IntClass{p+r}{q+s}\bigr)
    \end{aligned}%
  }{%
    a,b,c,d\in\mathbb N\dsep
    z=\IntClass{a}{b}\dsep w=\IntClass{c}{d}\vdash
    \IntClass{a+c}{b+d}\in\mathbb Z\land
    \exists p\in\mathbb N\,\exists q\in\mathbb N\,
    \exists r\in\mathbb N\,\exists s\in\mathbb N\,
    \bigl(z=\IntClass{p}{q}\land
      w=\IntClass{r}{s}\land
      \IntClass{a+c}{b+d}=\IntClass{p+r}{q+s}\bigr)%
  }[IntegerAdditionDescentCandidate]
    \proofstepwidestar[1]{a,b,c,d\in\mathbb N}{\rA}
    \proofstepwidestar[2]{z=\IntClass{a}{b}}{\rA}
    \proofstepwidestar[3]{w=\IntClass{c}{d}}{\rA}
    \proofstepwidestar[1]{a+c\in\mathbb N}{%
      \FormulaRefAuto{PeanoAddClosure}{1}}
    \proofstepwidestar[1]{b+d\in\mathbb N}{%
      \FormulaRefAuto{PeanoAddClosure}{1}}
    \proofstepwidestar[1]{\IntClass{a+c}{b+d}\in\mathbb Z}{%
      \FormulaRefAuto{IntegerClassInZ}{4,5}}
    \proofstepwidestar[]{%
      \IntClass{a+c}{b+d}=\IntClass{a+c}{b+d}}{\rII}
    \proofstepwidestar[1,2,3]{%
      \exists p\in\mathbb N\,\exists q\in\mathbb N\,
      \exists r\in\mathbb N\,\exists s\in\mathbb N\,
      \bigl(z=\IntClass{p}{q}\land
        w=\IntClass{r}{s}\land
        \IntClass{a+c}{b+d}=\IntClass{p+r}{q+s}\bigr)}{%
      \rEI{\rAI{1,\rAI{2,\rAI{3,7}}}}}
    \proofstepwidestar[1,2,3]{%
      \IntClass{a+c}{b+d}\in\mathbb Z\land
      \exists p\in\mathbb N\,\exists q\in\mathbb N\,
      \exists r\in\mathbb N\,\exists s\in\mathbb N\,
      \bigl(z=\IntClass{p}{q}\land
        w=\IntClass{r}{s}\land
        \IntClass{a+c}{b+d}=\IntClass{p+r}{q+s}\bigr)}{%
      \rAI{6,8}}
  \closeproofpartwideindR

  \proofpartwideindR[Existenz durch getrennte Repräsentantenwahl]{%
    \begin{aligned}[t]
      &z,w\in\mathbb Z\vdash
      \exists u\in\mathbb Z\,
      \exists p\in\mathbb N\,\exists q\in\mathbb N\,
      \exists r\in\mathbb N\,\exists s\in\mathbb N\,{}\\
      &\qquad
      \bigl(z=\IntClass{p}{q}\land
        w=\IntClass{r}{s}\land
        u=\IntClass{p+r}{q+s}\bigr)
    \end{aligned}%
  }{%
    z,w\in\mathbb Z\vdash
    \exists u\in\mathbb Z\,
    \exists p\in\mathbb N\,\exists q\in\mathbb N\,
    \exists r\in\mathbb N\,\exists s\in\mathbb N\,
    \bigl(z=\IntClass{p}{q}\land
      w=\IntClass{r}{s}\land
      u=\IntClass{p+r}{q+s}\bigr)%
  }[IntegerAdditionDescentExistence]
    \proofstepwidestar[1]{z\in\mathbb Z}{\rA}
    \proofstepwidestar[2]{w\in\mathbb Z}{\rA}
    \proofstepwidestar[1]{%
      \exists a\in\mathbb N\,\exists b\in\mathbb N\,
      z=\IntClass{a}{b}}{%
      \FormulaRefAuto{IntegerPairRepresentative}{1}}
    \proofstepwidestar[2]{%
      \exists c\in\mathbb N\,\exists d\in\mathbb N\,
      w=\IntClass{c}{d}}{%
      \FormulaRefAuto{IntegerPairRepresentative}{2}}
    \proofstepwidestar[5]{%
      a,b\in\mathbb N\land z=\IntClass{a}{b}}{\rA}
    \proofstepwidestar[6]{%
      c,d\in\mathbb N\land w=\IntClass{c}{d}}{\rA}
    \proofstepwidestar[5]{a,b\in\mathbb N}{\rAEa{5}}
    \proofstepwidestar[5]{z=\IntClass{a}{b}}{\rAEb{5}}
    \proofstepwidestar[6]{c,d\in\mathbb N}{\rAEa{6}}
    \proofstepwidestar[6]{w=\IntClass{c}{d}}{\rAEb{6}}
    \proofstepwidestar[5,6]{a,b,c,d\in\mathbb N}{\rAI{7,9}}
    \proofstepwidestar[5,6]{%
      \IntClass{a+c}{b+d}\in\mathbb Z\land
      \exists p\in\mathbb N\,\exists q\in\mathbb N\,
      \exists r\in\mathbb N\,\exists s\in\mathbb N\,
      \bigl(z=\IntClass{p}{q}\land
        w=\IntClass{r}{s}\land
        \IntClass{a+c}{b+d}=\IntClass{p+r}{q+s}\bigr)}{%
      \FormulaRefAuto{IntegerAdditionDescentCandidate}{11,8,10}}
    \proofstepwidestar[5,6]{%
      \exists u\in\mathbb Z\,
      \exists p\in\mathbb N\,\exists q\in\mathbb N\,
      \exists r\in\mathbb N\,\exists s\in\mathbb N\,
      \bigl(z=\IntClass{p}{q}\land
        w=\IntClass{r}{s}\land
        u=\IntClass{p+r}{q+s}\bigr)}{\rEI{12}}
    \proofstepwidestar[5]{%
      \exists u\in\mathbb Z\,
      \exists p\in\mathbb N\,\exists q\in\mathbb N\,
      \exists r\in\mathbb N\,\exists s\in\mathbb N\,
      \bigl(z=\IntClass{p}{q}\land
        w=\IntClass{r}{s}\land
        u=\IntClass{p+r}{q+s}\bigr)}{\rEE{4,6,13}}
    \proofstepwidestar[1,2]{%
      \exists u\in\mathbb Z\,
      \exists p\in\mathbb N\,\exists q\in\mathbb N\,
      \exists r\in\mathbb N\,\exists s\in\mathbb N\,
      \bigl(z=\IntClass{p}{q}\land
        w=\IntClass{r}{s}\land
        u=\IntClass{p+r}{q+s}\bigr)}{\rEE{3,5,14}}
  \closeproofpartwideindR

  \proofpartwideindR[Eindeutigkeit für feste Repräsentanten]{%
    \begin{aligned}[t]
      &a,b,c,d,a',b',c',d'\in\mathbb N\dsep{}\\
      &z=\IntClass{a}{b}\land w=\IntClass{c}{d}
        \land u=\IntClass{a+c}{b+d}\dsep{}\\
      &z=\IntClass{a'}{b'}\land w=\IntClass{c'}{d'}
        \land v=\IntClass{a'+c'}{b'+d'}
      \vdash u=v
    \end{aligned}%
  }{%
    a,b,c,d,a',b',c',d'\in\mathbb N\dsep
    z=\IntClass{a}{b}\land w=\IntClass{c}{d}
      \land u=\IntClass{a+c}{b+d}\dsep
    z=\IntClass{a'}{b'}\land w=\IntClass{c'}{d'}
      \land v=\IntClass{a'+c'}{b'+d'}
    \vdash u=v%
  }[IntegerAdditionDescentFixedUniqueness]
    \proofstepwidestar[1]{%
      a,b,c,d,a',b',c',d'\in\mathbb N}{\rA}
    \proofstepwidestar[2]{%
      z=\IntClass{a}{b}\land w=\IntClass{c}{d}
      \land u=\IntClass{a+c}{b+d}}{\rA}
    \proofstepwidestar[3]{%
      z=\IntClass{a'}{b'}\land w=\IntClass{c'}{d'}
      \land v=\IntClass{a'+c'}{b'+d'}}{\rA}
    \proofstepwidestar[2]{z=\IntClass{a}{b}}{\rAEa{2}}
    \proofstepwidestar[2]{w=\IntClass{c}{d}}{\rAEa{\rAEb{2}}}
    \proofstepwidestar[2]{u=\IntClass{a+c}{b+d}}{\rAEb{\rAEb{2}}}
    \proofstepwidestar[3]{z=\IntClass{a'}{b'}}{\rAEa{3}}
    \proofstepwidestar[3]{w=\IntClass{c'}{d'}}{\rAEa{\rAEb{3}}}
    \proofstepwidestar[3]{v=\IntClass{a'+c'}{b'+d'}}{%
      \rAEb{\rAEb{3}}}
    \proofstepwidestar[2,3]{%
      \IntClass{a}{b}=\IntClass{a'}{b'}}{\rIE{4,7}}
    \proofstepwidestar[2,3]{%
      \IntClass{c}{d}=\IntClass{c'}{d'}}{\rIE{5,8}}
    \proofstepwidestar[1,2,3]{%
      \IntClass{a+c}{b+d}=\IntClass{a'+c'}{b'+d'}}{%
      \FormulaRefAuto{IntegerAdditionCompatible}{1,10,11}}
    \proofstepwidestar[1,2,3]{u=v}{\rIE{6,9,12}}
  \closeproofpartwideindR

  \proofpartwideindR[Eindeutigkeit beliebiger Ergebniszeugen]{%
    \begin{aligned}[t]
      &\bigl(u\in\mathbb Z\land
      \exists a,b,c,d\in\mathbb N\,
      (z=\IntClass{a}{b}\land w=\IntClass{c}{d}
       \land u=\IntClass{a+c}{b+d})\bigr)\dsep{}\\
      &\bigl(v\in\mathbb Z\land
      \exists a',b',c',d'\in\mathbb N\,
      (z=\IntClass{a'}{b'}\land w=\IntClass{c'}{d'}
       \land v=\IntClass{a'+c'}{b'+d'})\bigr)
      \vdash u=v
    \end{aligned}%
  }{%
    \bigl(u\in\mathbb Z\land
    \exists a,b,c,d\in\mathbb N\,
    (z=\IntClass{a}{b}\land w=\IntClass{c}{d}
     \land u=\IntClass{a+c}{b+d})\bigr)\dsep
    \bigl(v\in\mathbb Z\land
    \exists a',b',c',d'\in\mathbb N\,
    (z=\IntClass{a'}{b'}\land w=\IntClass{c'}{d'}
     \land v=\IntClass{a'+c'}{b'+d'})\bigr)
    \vdash u=v%
  }[IntegerAdditionDescentUniqueness]
    \proofstepwidestar[1]{%
      u\in\mathbb Z\land
      \exists a,b,c,d\in\mathbb N\,
      (z=\IntClass{a}{b}\land w=\IntClass{c}{d}
       \land u=\IntClass{a+c}{b+d})}{\rA}
    \proofstepwidestar[2]{%
      v\in\mathbb Z\land
      \exists a',b',c',d'\in\mathbb N\,
      (z=\IntClass{a'}{b'}\land w=\IntClass{c'}{d'}
       \land v=\IntClass{a'+c'}{b'+d'})}{\rA}
    \proofstepwidestar[1]{%
      \exists a,b,c,d\in\mathbb N\,
      (z=\IntClass{a}{b}\land w=\IntClass{c}{d}
       \land u=\IntClass{a+c}{b+d})}{\rAEb{1}}
    \proofstepwidestar[2]{%
      \exists a',b',c',d'\in\mathbb N\,
      (z=\IntClass{a'}{b'}\land w=\IntClass{c'}{d'}
       \land v=\IntClass{a'+c'}{b'+d'})}{\rAEb{2}}
    \proofstepwidestar[5]{%
      a,b,c,d\in\mathbb N\land
      z=\IntClass{a}{b}\land w=\IntClass{c}{d}
      \land u=\IntClass{a+c}{b+d}}{\rA}
    \proofstepwidestar[6]{%
      a',b',c',d'\in\mathbb N\land
      z=\IntClass{a'}{b'}\land w=\IntClass{c'}{d'}
      \land v=\IntClass{a'+c'}{b'+d'}}{\rA}
    \proofstepwidestar[5]{a,b,c,d\in\mathbb N}{\rAEa{5}}
    \proofstepwidestar[5]{%
      z=\IntClass{a}{b}\land w=\IntClass{c}{d}
      \land u=\IntClass{a+c}{b+d}}{\rAEb{5}}
    \proofstepwidestar[6]{a',b',c',d'\in\mathbb N}{\rAEa{6}}
    \proofstepwidestar[6]{%
      z=\IntClass{a'}{b'}\land w=\IntClass{c'}{d'}
      \land v=\IntClass{a'+c'}{b'+d'}}{\rAEb{6}}
    \proofstepwidestar[5,6]{%
      a,b,c,d,a',b',c',d'\in\mathbb N}{\rAI{7,9}}
    \proofstepwidestar[5,6]{u=v}{%
      \FormulaRefAuto{IntegerAdditionDescentFixedUniqueness}{11,8,10}}
    \proofstepwidestar[5]{u=v}{\rEE{4,6,12}}
    \proofstepwidestar[1,2]{u=v}{\rEE{3,5,13}}
  \closeproofpartwideindR

  \proofpartwideindR[Abschluss des Abstiegs]{%
    \begin{aligned}[t]
      &z,w\in\mathbb Z\vdash
      \exists! u\in\mathbb Z\,
      \exists p\in\mathbb N\,\exists q\in\mathbb N\,
      \exists r\in\mathbb N\,\exists s\in\mathbb N\,{}\\
      &\qquad
      \bigl(z=\IntClass{p}{q}\land w=\IntClass{r}{s}
      \land u=\IntClass{p+r}{q+s}\bigr)
    \end{aligned}%
  }{%
    z,w\in\mathbb Z\vdash
    \exists! u\in\mathbb Z\,
    \exists p\in\mathbb N\,\exists q\in\mathbb N\,
    \exists r\in\mathbb N\,\exists s\in\mathbb N\,
    \bigl(z=\IntClass{p}{q}\land w=\IntClass{r}{s}
    \land u=\IntClass{p+r}{q+s}\bigr)%
  }[IntegerAdditionDescentConclusion]
    \proofstepwidestar[1]{z,w\in\mathbb Z}{\rA}
    \proofstepwidestar[1]{%
      \exists u\in\mathbb Z\,
      \exists p\in\mathbb N\,\exists q\in\mathbb N\,
      \exists r\in\mathbb N\,\exists s\in\mathbb N\,
      (z=\IntClass{p}{q}\land w=\IntClass{r}{s}
       \land u=\IntClass{p+r}{q+s})}{%
      \FormulaRefAuto{IntegerAdditionDescentExistence}{1}}
    \proofstepwidestar[3]{%
      u\in\mathbb Z\land
      \exists p\in\mathbb N\,\exists q\in\mathbb N\,
      \exists r\in\mathbb N\,\exists s\in\mathbb N\,
      (z=\IntClass{p}{q}\land w=\IntClass{r}{s}
       \land u=\IntClass{p+r}{q+s})}{\rA}
    \proofstepwidestar[4]{%
      v\in\mathbb Z\land
      \exists p\in\mathbb N\,\exists q\in\mathbb N\,
      \exists r\in\mathbb N\,\exists s\in\mathbb N\,
      (z=\IntClass{p}{q}\land w=\IntClass{r}{s}
       \land v=\IntClass{p+r}{q+s})}{\rA}
    \proofstepwidestar[3,4]{u=v}{%
      \FormulaRefAuto{IntegerAdditionDescentUniqueness}{3,4}}
    \proofstepwidestar[1]{%
      \exists! u\in\mathbb Z\,
      \exists p\in\mathbb N\,\exists q\in\mathbb N\,
      \exists r\in\mathbb N\,\exists s\in\mathbb N\,
      (z=\IntClass{p}{q}\land w=\IntClass{r}{s}
       \land u=\IntClass{p+r}{q+s})}{\UEI{2,3,4,5}}
  \closeproofpartwideindR
\end{tabproofsplitwide}

Der eindeutige Existenzsatz ist der Quotientenabstieg der Addition; die
folgende Klassenregel bezeichnet dieses eindeutig bestimmte Ergebnis.

\FormulaDefDeltaK[Addition ganzer Zahlen]{%
  a,b,c,d\in\mathbb N\vdash
  \IntClass{a}{b}+\IntClass{c}{d}
  \coloneqq\IntClass{a+c}{b+d}%
}{IntegerAdditionDef}{%
  \DeltaRow{Mengen}{a\dsep b\dsep c\dsep d}%
  \DeltaRow{Neue Symbole}{+}%
}

\FormulaThmDeltaK[Abgeschlossenheit der Addition ganzer Zahlen]{%
  z\in\mathbb Z\dsep w\in\mathbb Z\vdash z+w\in\mathbb Z%
}{IntAddClosure}{%
  \DeltaRow{Mengen}{z\dsep w\dsep a\dsep b\dsep c\dsep d}%
}
\begin{tabproofsplitwide}
  \proofpartwideindR[Abgeschlossenheit für feste Repräsentanten]{%
    \begin{aligned}[t]
      &a,b,c,d\in\mathbb N\land
        z=\IntClass{a}{b}\land w=\IntClass{c}{d}\\
      &\vdash z+w\in\mathbb Z
    \end{aligned}
  }{%
    a,b,c,d\in\mathbb N\land
    z=\IntClass{a}{b}\land w=\IntClass{c}{d}
    \vdash z+w\in\mathbb Z%
  }[IntAddClosureForRepresentatives]
    \proofstepwidestar[1]{a,b,c,d\in\mathbb N\land
      z=\IntClass{a}{b}\land w=\IntClass{c}{d}}{\rA}
    \proofstepwidestar[1]{z+w=\IntClass{a+c}{b+d}}{%
      \FormulaRefAuto{IntegerAdditionDef}{1}}
    \proofstepwidestar[1]{a+c\in\mathbb N}{%
      \FormulaRefAuto{PeanoAddClosure}{1}}
    \proofstepwidestar[1]{b+d\in\mathbb N}{%
      \FormulaRefAuto{PeanoAddClosure}{1}}
    \proofstepwidestar[1]{\IntClass{a+c}{b+d}\in\mathbb Z}{%
      \FormulaRefAuto{IntegerClassInZ}{3,4}}
    \proofstepwidestar[1]{z+w\in\mathbb Z}{\rIE{2,5}}
  \closeproofpartwideindR

  \proofpartwideindR[Elimination der Repräsentanten]{%
    z\in\mathbb Z\dsep w\in\mathbb Z\vdash z+w\in\mathbb Z%
  }{%
    z\in\mathbb Z\dsep w\in\mathbb Z\vdash z+w\in\mathbb Z%
  }[IntAddClosureRepresentativeElimination]
    \proofstepwidestar[1]{z\in\mathbb Z}{\rA}
    \proofstepwidestar[2]{w\in\mathbb Z}{\rA}
    \proofstepwidestar[1]{\exists a,b\in\mathbb N\,
      z=\IntClass{a}{b}}{\FormulaRefAuto{IntegerPairRepresentative}{1}}
    \proofstepwidestar[2]{\exists c,d\in\mathbb N\,
      w=\IntClass{c}{d}}{\FormulaRefAuto{IntegerPairRepresentative}{2}}
    \proofstepwidestar[5]{a,b,c,d\in\mathbb N\land
      z=\IntClass{a}{b}\land w=\IntClass{c}{d}}{\rA}
    \proofstepwidestar[5]{z+w\in\mathbb Z}{%
      \FormulaRefAuto{IntAddClosureForRepresentatives}{5}}
    \proofstepwidestar[1,2]{z+w\in\mathbb Z}{\rEE{3,4,5,6}}
  \closeproofpartwideindR
\end{tabproofsplitwide}

\FormulaThmDeltaK[Assoziativität der Addition ganzer Zahlen]{%
  z,w,u\in\mathbb Z\vdash(z+w)+u=z+(w+u)%
}{IntAddAssociative}{%
  \DeltaRow{Mengen}{z\dsep w\dsep u\dsep
    a\dsep b\dsep c\dsep d\dsep e\dsep f}%
}
\begin{tabproofsplitwide}
  \proofpartwideindR[Rechnung mit Repräsentanten]{%
    \begin{aligned}[t]
      &a,b,c,d,e,f\in\mathbb N\vdash{}\\
      &(\IntClass{a}{b}+\IntClass{c}{d})+\IntClass{e}{f}
       =\IntClass{a}{b}+(\IntClass{c}{d}+\IntClass{e}{f})
    \end{aligned}
  }{%
    a,b,c,d,e,f\in\mathbb N\vdash
    (\IntClass{a}{b}+\IntClass{c}{d})+\IntClass{e}{f}
    =\IntClass{a}{b}+(\IntClass{c}{d}+\IntClass{e}{f})
  }[IntAddAssociativeOnClasses]
    \proofstepwidestar[1]{a,b,c,d,e,f\in\mathbb N}{\rA}
    \proofstepwidestar[1]{%
      (\IntClass{a}{b}+\IntClass{c}{d})+\IntClass{e}{f}
      =\IntClass{(a+c)+e}{(b+d)+f}}{%
      \FormulaRefAuto{IntegerAdditionDef}{1}}
    \proofstepwidestar[1]{%
      \IntClass{a}{b}+(\IntClass{c}{d}+\IntClass{e}{f})
      =\IntClass{a+(c+e)}{b+(d+f)}}{%
      \FormulaRefAuto{IntegerAdditionDef}{1}}
    \proofstepwidestar[1]{(a+c)+e=a+(c+e)}{%
      \FormulaRefAuto{PeanoAddAssociative}{1}}
    \proofstepwidestar[1]{(b+d)+f=b+(d+f)}{%
      \FormulaRefAuto{PeanoAddAssociative}{1}}
    \proofstepwidestar[1]{%
      \IntClass{(a+c)+e}{(b+d)+f}
      =\IntClass{a+(c+e)}{b+(d+f)}}{\rIE{4,5}}
    \proofstepwidestar[1]{%
      (\IntClass{a}{b}+\IntClass{c}{d})+\IntClass{e}{f}
      =\IntClass{a}{b}+(\IntClass{c}{d}+\IntClass{e}{f})}{%
      \rIE{2,3,6}}
  \closeproofpartwideindR

  \proofpartwide{Schluss}
    \proofstepwidestar[1]{z,w,u\in\mathbb Z}{\rA}
    \proofstepwidestar[1]{z\in\mathbb Z}{\rAEa{1}}
    \proofstepwidestar[1]{w\in\mathbb Z}{\rAEa{\rAEb{1}}}
    \proofstepwidestar[1]{u\in\mathbb Z}{\rAEb{\rAEb{1}}}
    \proofstepwidestar[1]{\exists a,b\in\mathbb N\,
      z=\IntClass{a}{b}}{\FormulaRefAuto{IntegerPairRepresentative}{2}}
    \proofstepwidestar[1]{\exists c,d\in\mathbb N\,
      w=\IntClass{c}{d}}{\FormulaRefAuto{IntegerPairRepresentative}{3}}
    \proofstepwidestar[1]{\exists e,f\in\mathbb N\,
      u=\IntClass{e}{f}}{\FormulaRefAuto{IntegerPairRepresentative}{4}}
    \proofstepwidestar[8]{a,b,c,d,e,f\in\mathbb N\land
      z=\IntClass{a}{b}\land w=\IntClass{c}{d}\land
      u=\IntClass{e}{f}}{\rA}
    \proofstepwidestar[8]{%
      (\IntClass{a}{b}+\IntClass{c}{d})+\IntClass{e}{f}
      =\IntClass{a}{b}+(\IntClass{c}{d}+\IntClass{e}{f})}{%
      \FormulaRefAuto{IntAddAssociativeOnClasses}{8}}
    \proofstepwidestar[8]{(z+w)+u=z+(w+u)}{\rIE{8,9}}
    \proofstepwidestar[1]{(z+w)+u=z+(w+u)}{%
      \rEE{5,6,7,8,10}}
  \closeproofpartwide
\end{tabproofsplitwide}

\FormulaThmDeltaK[Kommutativität der Addition ganzer Zahlen]{%
  z,w\in\mathbb Z\vdash z+w=w+z%
}{IntAddCommutative}{%
  \DeltaRow{Mengen}{z\dsep w\dsep a\dsep b\dsep c\dsep d}%
}
\begin{tabproofsplitwide}
  \proofpartwideindR[Vertauschung fester Repräsentanten]{%
    \begin{aligned}[t]
      &a,b,c,d\in\mathbb N\land
        z=\IntClass{a}{b}\land w=\IntClass{c}{d}\\
      &\vdash z+w=w+z
    \end{aligned}
  }{%
    a,b,c,d\in\mathbb N\land
    z=\IntClass{a}{b}\land w=\IntClass{c}{d}
    \vdash z+w=w+z%
  }[IntAddCommutativeOnRepresentatives]
    \proofstepwidestar[1]{a,b,c,d\in\mathbb N\land
      z=\IntClass{a}{b}\land w=\IntClass{c}{d}}{\rA}
    \proofstepwidestar[1]{z+w=\IntClass{a+c}{b+d}}{%
      \FormulaRefAuto{IntegerAdditionDef}{1}}
    \proofstepwidestar[1]{w+z=\IntClass{c+a}{d+b}}{%
      \FormulaRefAuto{IntegerAdditionDef}{1}}
    \proofstepwidestar[1]{a+c=c+a}{%
      \FormulaRefAuto{PeanoAddCommutative}{1}}
    \proofstepwidestar[1]{b+d=d+b}{%
      \FormulaRefAuto{PeanoAddCommutative}{1}}
    \proofstepwidestar[1]{%
      \IntClass{a+c}{b+d}=\IntClass{c+a}{d+b}}{\rIE{4,5}}
    \proofstepwidestar[1]{z+w=w+z}{\rIE{2,3,6}}
  \closeproofpartwideindR

  \proofpartwideindR[Getrennte Repräsentantenwahl]{%
    z,w\in\mathbb Z\vdash z+w=w+z%
  }{%
    z,w\in\mathbb Z\vdash z+w=w+z%
  }[IntAddCommutativeRepresentativeElimination]
    \proofstepwidestar[1]{z,w\in\mathbb Z}{\rA}
    \proofstepwidestar[1]{z\in\mathbb Z}{\rAEa{1}}
    \proofstepwidestar[1]{w\in\mathbb Z}{\rAEb{1}}
    \proofstepwidestar[1]{\exists a,b\in\mathbb N\,
      z=\IntClass{a}{b}}{\FormulaRefAuto{IntegerPairRepresentative}{2}}
    \proofstepwidestar[1]{\exists c,d\in\mathbb N\,
      w=\IntClass{c}{d}}{\FormulaRefAuto{IntegerPairRepresentative}{3}}
    \proofstepwidestar[6]{a,b,c,d\in\mathbb N\land
      z=\IntClass{a}{b}\land w=\IntClass{c}{d}}{\rA}
    \proofstepwidestar[6]{z+w=w+z}{%
      \FormulaRefAuto{IntAddCommutativeOnRepresentatives}{6}}
    \proofstepwidestar[1]{z+w=w+z}{\rEE{4,5,6,7}}
  \closeproofpartwideindR
\end{tabproofsplitwide}

\FormulaThmDeltaK[Nullgesetze der Addition ganzer Zahlen]{%
  z\in\mathbb Z\vdash
  z+0_{\mathbb Z}=z\land0_{\mathbb Z}+z=z%
}{IntAddZero}{%
  \DeltaRow{Mengen}{z\dsep a\dsep b}%
}
\begin{tabproofsplitwide}
  \proofpartwideindR[Rechtsneutrales Element für einen Repräsentanten]{%
    \begin{aligned}[t]
      &a,b\in\mathbb N\land z=\IntClass{a}{b}\\
      &\vdash z+0_{\mathbb Z}=z
    \end{aligned}
  }{%
    a,b\in\mathbb N\land z=\IntClass{a}{b}
    \vdash z+0_{\mathbb Z}=z%
  }[IntAddRightZeroForRepresentative]
    \proofstepwidestar[1]{a,b\in\mathbb N\land
      z=\IntClass{a}{b}}{\rA}
    \proofstepwidestar[]{0_{\mathbb Z}=\IntClass{0}{0}}{%
      \FormulaRefAuto{IntegerZeroClass}}
    \proofstepwidestar[1]{z+0_{\mathbb Z}=\IntClass{a+0}{b+0}}{%
      \FormulaRefAuto{IntegerAdditionDef}{1,2}}
    \proofstepwidestar[1]{a+0=a}{%
      \FormulaRefAuto{PeanoAddRightZero}{1}}
    \proofstepwidestar[1]{b+0=b}{%
      \FormulaRefAuto{PeanoAddRightZero}{1}}
    \proofstepwidestar[1]{z+0_{\mathbb Z}=z}{\rIE{1,3,4,5}}
  \closeproofpartwideindR

  \proofpartwideindR[Elimination des Repräsentanten]{%
    z\in\mathbb Z\vdash
    z+0_{\mathbb Z}=z\land0_{\mathbb Z}+z=z%
  }{%
    z\in\mathbb Z\vdash
    z+0_{\mathbb Z}=z\land0_{\mathbb Z}+z=z%
  }[IntAddZeroRepresentativeElimination]
    \proofstepwidestar[1]{z\in\mathbb Z}{\rA}
    \proofstepwidestar[1]{\exists a,b\in\mathbb N\,
      z=\IntClass{a}{b}}{\FormulaRefAuto{IntegerPairRepresentative}{1}}
    \proofstepwidestar[3]{a,b\in\mathbb N\land
      z=\IntClass{a}{b}}{\rA}
    \proofstepwidestar[3]{z+0_{\mathbb Z}=z}{%
      \FormulaRefAuto{IntAddRightZeroForRepresentative}{3}}
    \proofstepwidestar[3]{0_{\mathbb Z}+z=z}{%
      \FormulaRefAuto{IntAddCommutative}{%
        \FormulaRefAuto{IntZeroCarrier},1,
        \FormulaRefAuto{a=b\vdash b=a}{4}}}
    \proofstepwidestar[3]{%
      z+0_{\mathbb Z}=z\land0_{\mathbb Z}+z=z}{\rAI{4,5}}
    \proofstepwidestar[1]{%
      z+0_{\mathbb Z}=z\land0_{\mathbb Z}+z=z}{\rEE{2,3,6}}
  \closeproofpartwideindR
\end{tabproofsplitwide}

\FormulaThmDeltaK[Additive Inversenregeln]{%
  z\in\mathbb Z\vdash
  z+(-z)=0_{\mathbb Z}\land(-z)+z=0_{\mathbb Z}%
}{IntAddInverse}{%
  \DeltaRow{Mengen}{z\dsep a\dsep b}%
}
\begin{tabproofsplitwide}
  \proofpartwideindR[Inverse Rechnung auf einer Klasse]{%
    a,b\in\mathbb N\vdash
    \IntClass{a}{b}+(-\IntClass{a}{b})=0_{\mathbb Z}%
  }{%
    a,b\in\mathbb N\vdash
    \IntClass{a}{b}+(-\IntClass{a}{b})=0_{\mathbb Z}%
  }[IntAddInverseOnClass]
    \proofstepwidestar[1]{a,b\in\mathbb N}{\rA}
    \proofstepwidestar[1]{-\IntClass{a}{b}=\IntClass{b}{a}}{%
      \FormulaRefAuto{IntegerNegationDef}{1}}
    \proofstepwidestar[1]{%
      \IntClass{a}{b}+(-\IntClass{a}{b})=\IntClass{a+b}{b+a}}{%
      \FormulaRefAuto{IntegerAdditionDef}{1,2}}
    \proofstepwidestar[1]{a+b=b+a}{\FormulaRefAuto{PeanoAddCommutative}{1}}
    \proofstepwidestar[1]{\IntClass{a+b}{b+a}=\IntClass{0}{0}}{%
      \FormulaRefAuto{IntegerClassEquality}{%
        \FormulaRefAuto{PeanoAddClosure}{1},
        \FormulaRefAuto{PeanoAddClosure}{1},
        \FormulaRefAuto{PeanoZero},
        \FormulaRefAuto{PeanoZero},
        \rIE{\FormulaRefAuto{PeanoAddRightZero},%
             \FormulaRefAuto{PeanoAddLeftZero},4}}}
    \proofstepwidestar[]{0_{\mathbb Z}=\IntClass{0}{0}}{%
      \FormulaRefAuto{IntegerZeroClass}}
    \proofstepwidestar[1]{%
      \IntClass{a}{b}+(-\IntClass{a}{b})=0_{\mathbb Z}}{\rIE{3,5,6}}
  \closeproofpartwideindR

  \proofpartwide{Schluss}
    \proofstepwidestar[1]{z\in\mathbb Z}{\rA}
    \proofstepwidestar[1]{\exists a,b\in\mathbb N\,
      z=\IntClass{a}{b}}{\FormulaRefAuto{IntegerPairRepresentative}{1}}
    \proofstepwidestar[3]{a,b\in\mathbb N\land z=\IntClass{a}{b}}{\rA}
    \proofstepwidestar[3]{z+(-z)=0_{\mathbb Z}}{%
      \FormulaRefAuto{IntAddInverseOnClass}{3}}
    \proofstepwidestar[1]{-z\in\mathbb Z}{\FormulaRefAuto{IntNegClosure}{1}}
    \proofstepwidestar[1,3]{(-z)+z=z+(-z)}{%
      \FormulaRefAuto{IntAddCommutative}{5,1}}
    \proofstepwidestar[1,3]{(-z)+z=0_{\mathbb Z}}{\rIE{6,4}}
    \proofstepwidestar[1,3]{%
      z+(-z)=0_{\mathbb Z}\land(-z)+z=0_{\mathbb Z}}{\rAI{4,7}}
    \proofstepwidestar[1]{%
      z+(-z)=0_{\mathbb Z}\land(-z)+z=0_{\mathbb Z}}{\rEE{2,3,8}}
  \closeproofpartwide
\end{tabproofsplitwide}

\FormulaThmDeltaK[Negation der Ganzzahlnull]{%
  -0_{\mathbb Z}=0_{\mathbb Z}%
}{IntegerNegativeZero}{%
  \DeltaRow{Mengen}{0_{\mathbb Z}}%
}
\begin{tabproof}
  \proofstep{}{0_{\mathbb Z}\in\mathbb Z}{%
    \FormulaRefAuto{IntZeroCarrier}}
  \proofstep{}{-0_{\mathbb Z}\in\mathbb Z}{%
    \FormulaRefAuto{IntNegClosure}{1}}
  \proofstep{}{(-0_{\mathbb Z})+0_{\mathbb Z}=-0_{\mathbb Z}}{%
    \rAEa{\FormulaRefAuto{IntAddZero}{2}}}
  \proofstep{}{(-0_{\mathbb Z})+0_{\mathbb Z}=0_{\mathbb Z}}{%
    \rAEb{\FormulaRefAuto{IntAddInverse}{1}}}
  \proofstep{}{-0_{\mathbb Z}=(-0_{\mathbb Z})+0_{\mathbb Z}}{%
    \FormulaRefAuto{a=b\vdash b=a}{3}}
  \proofstep{}{-0_{\mathbb Z}=0_{\mathbb Z}}{%
    \rIE{5,4}}
\end{tabproof}

\section{Multiplikation}

\FormulaThmDeltaK[Repräsentantenunabhängigkeit der Multiplikation]{%
  \begin{aligned}[t]
    &a,b,a',b',c,d,c',d'\in\mathbb N\dsep
      \IntClass{a}{b}=\IntClass{a'}{b'}\dsep
      \IntClass{c}{d}=\IntClass{c'}{d'}\vdash{}\\
    &\IntClass{a\cdot c+b\cdot d}{a\cdot d+b\cdot c}
      =
      \IntClass{a'\cdot c'+b'\cdot d'}{a'\cdot d'+b'\cdot c'}
  \end{aligned}%
}{IntegerMultiplicationCompatible}{%
  \DeltaRow{Mengen}{a\dsep b\dsep a'\dsep b'\dsep
    c\dsep d\dsep c'\dsep d'}%
}
\begin{tabproofsplitwide}
  \proofpartwideindR[Wechsel des ersten Repräsentanten]{%
    \begin{aligned}[t]
      &a,b,a',b',c,d\in\mathbb N\dsep
      \IntClass{a}{b}=\IntClass{a'}{b'}\vdash{}\\
      &\IntClass{a c+b d}{a d+b c}
       =\IntClass{a'c+b'd}{a'd+b'c}
    \end{aligned}
  }{%
    a,b,a',b',c,d\in\mathbb N\dsep
    \IntClass{a}{b}=\IntClass{a'}{b'}
    \vdash
    \IntClass{a c+b d}{a d+b c}
    =\IntClass{a'c+b'd}{a'd+b'c}
  }[IntegerMultiplicationCompatibleLeft]
    \proofstepwidestar[1]{a,b,a',b',c,d\in\mathbb N}{\rA}
    \proofstepwidestar[2]{\IntClass{a}{b}=\IntClass{a'}{b'}}{\rA}
    \proofstepwidestar[1,2]{a+b'=b+a'}{%
      \FormulaRefAuto{IntegerClassEquality}{1,2}}
    \proofstepwide[1,2]{(a+b')c}{=}{(b+a')c}{\rIE{3}}
    \proofstepwide[1]{(a+b')c}{=}{ac+b'c}{%
      \FormulaRefAuto{PeanoMulRightDistributive}{1}}
    \proofstepwide[1]{(b+a')c}{=}{bc+a'c}{%
      \FormulaRefAuto{PeanoMulRightDistributive}{1}}
    \proofstepwidestar[1,2]{ac+b'c=bc+a'c}{\rIE{4,5,6}}
    \proofstepwide[1,2]{(a+b')d}{=}{(b+a')d}{\rIE{3}}
    \proofstepwide[1]{(a+b')d}{=}{ad+b'd}{%
      \FormulaRefAuto{PeanoMulRightDistributive}{1}}
    \proofstepwide[1]{(b+a')d}{=}{bd+a'd}{%
      \FormulaRefAuto{PeanoMulRightDistributive}{1}}
    \proofstepwidestar[1,2]{ad+b'd=bd+a'd}{\rIE{8,9,10}}
    \proofstepwidestar[1]{a'd+b'c=b'c+a'd}{%
      \FormulaRefAuto{PeanoAddCommutative}{1}}
    \proofstepwidestar[1,2]{%
      (ac+bd)+(a'd+b'c)=(ac+bd)+(b'c+a'd)}{\rIE{12}}
    \proofstepwidestar[1]{%
      (ac+bd)+(b'c+a'd)=(ac+b'c)+(bd+a'd)}{%
      \FormulaRefAuto{PeanoAddInterchange}{1}}
    \proofstepwidestar[1,2]{bd+a'd=ad+b'd}{%
      \FormulaRefAuto{a=b\vdash b=a}{11}}
    \proofstepwidestar[1,2]{%
      (ac+b'c)+(bd+a'd)=(bc+a'c)+(ad+b'd)}{\rIE{7,15}}
    \proofstepwidestar[1]{%
      (bc+a'c)+(ad+b'd)=(bc+ad)+(a'c+b'd)}{%
      \FormulaRefAuto{PeanoAddInterchange}{1}}
    \proofstepwidestar[1]{bc+ad=ad+bc}{%
      \FormulaRefAuto{PeanoAddCommutative}{1}}
    \proofstepwidestar[1]{%
      (bc+ad)+(a'c+b'd)=(ad+bc)+(a'c+b'd)}{\rIE{18}}
    \proofstepwidestar[1,2]{%
      (ac+bd)+(a'd+b'c)=(ad+bc)+(a'c+b'd)}{%
      \rIE{13,14,16,17,19}}
    \proofstepwidestar[1,2]{%
      \IntClass{ac+bd}{ad+bc}=\IntClass{a'c+b'd}{a'd+b'c}}{%
      \FormulaRefAuto{IntegerClassEquality}{1,20}}
  \closeproofpartwideindR

  \proofpartwideindR[Wechsel des zweiten Repräsentanten]{%
    \begin{aligned}[t]
      &a,b,c,d,c',d'\in\mathbb N\dsep
      \IntClass{c}{d}=\IntClass{c'}{d'}\vdash{}\\
      &\IntClass{a c+b d}{a d+b c}
       =\IntClass{a c'+b d'}{a d'+b c'}
    \end{aligned}
  }{%
    a,b,c,d,c',d'\in\mathbb N\dsep
    \IntClass{c}{d}=\IntClass{c'}{d'}
    \vdash
    \IntClass{a c+b d}{a d+b c}
    =\IntClass{a c'+b d'}{a d'+b c'}
  }[IntegerMultiplicationCompatibleRight]
    \proofstepwidestar[1]{a,b,c,d,c',d'\in\mathbb N}{\rA}
    \proofstepwidestar[2]{\IntClass{c}{d}=\IntClass{c'}{d'}}{\rA}
    \proofstepwidestar[1,2]{c+d'=d+c'}{%
      \FormulaRefAuto{IntegerClassEquality}{1,2}}
    \proofstepwide[1,2]{a(c+d')}{=}{a(d+c')}{\rIE{3}}
    \proofstepwide[1]{a(c+d')}{=}{ac+ad'}{%
      \FormulaRefAuto{PeanoMulLeftDistributive}{1}}
    \proofstepwide[1]{a(d+c')}{=}{ad+ac'}{%
      \FormulaRefAuto{PeanoMulLeftDistributive}{1}}
    \proofstepwidestar[1,2]{ac+ad'=ad+ac'}{\rIE{4,5,6}}
    \proofstepwide[1,2]{b(c+d')}{=}{b(d+c')}{\rIE{3}}
    \proofstepwide[1]{b(c+d')}{=}{bc+bd'}{%
      \FormulaRefAuto{PeanoMulLeftDistributive}{1}}
    \proofstepwide[1]{b(d+c')}{=}{bd+bc'}{%
      \FormulaRefAuto{PeanoMulLeftDistributive}{1}}
    \proofstepwidestar[1,2]{bc+bd'=bd+bc'}{\rIE{8,9,10}}
    \proofstepwidestar[1]{%
      (ac+bd)+(ad'+bc')=(ac+ad')+(bd+bc')}{%
      \FormulaRefAuto{PeanoAddInterchange}{1}}
    \proofstepwidestar[1,2]{bd+bc'=bc+bd'}{%
      \FormulaRefAuto{a=b\vdash b=a}{11}}
    \proofstepwidestar[1,2]{%
      (ac+ad')+(bd+bc')=(ad+ac')+(bc+bd')}{\rIE{7,13}}
    \proofstepwidestar[1]{%
      (ad+ac')+(bc+bd')=(ad+bc)+(ac'+bd')}{%
      \FormulaRefAuto{PeanoAddInterchange}{1}}
    \proofstepwidestar[1,2]{%
      (ac+bd)+(ad'+bc')=(ad+bc)+(ac'+bd')}{%
      \rIE{12,14,15}}
    \proofstepwidestar[1,2]{%
      \IntClass{ac+bd}{ad+bc}=\IntClass{ac'+bd'}{ad'+bc'}}{%
      \FormulaRefAuto{IntegerClassEquality}{1,16}}
  \closeproofpartwideindR

  \proofpartwide{Zusammenführung}
    \proofstepwidestar[1]{a,b,a',b',c,d,c',d'\in\mathbb N}{\rA}
    \proofstepwidestar[2]{\IntClass{a}{b}=\IntClass{a'}{b'}}{\rA}
    \proofstepwidestar[3]{\IntClass{c}{d}=\IntClass{c'}{d'}}{\rA}
    \proofstepwidestar[1,2]{%
      \IntClass{ac+bd}{ad+bc}
      =\IntClass{a'c+b'd}{a'd+b'c}}{%
      \FormulaRefAuto{IntegerMultiplicationCompatibleLeft}{1,2}}
    \proofstepwidestar[1,3]{%
      \IntClass{a'c+b'd}{a'd+b'c}
      =\IntClass{a'c'+b'd'}{a'd'+b'c'}}{%
      \FormulaRefAuto{IntegerMultiplicationCompatibleRight}{1,3}}
    \proofstepwidestar[1,2,3]{%
      \IntClass{ac+bd}{ad+bc}
      =\IntClass{a'c'+b'd'}{a'd'+b'c'}}{%
      \FormulaRefAuto{a=b\dsep b=c\vdash a=c}{4,5}}
  \closeproofpartwide
\end{tabproofsplitwide}

\FormulaThmDeltaK[Quotientenabstieg der Multiplikation]{%
  \begin{aligned}[t]
    &z,w\in\mathbb Z\vdash
    \exists! u\in\mathbb Z\,
    \exists p\in\mathbb N\,\exists q\in\mathbb N\,
    \exists r\in\mathbb N\,\exists s\in\mathbb N\,{}\\
    &\qquad
    \bigl(z=\IntClass{p}{q}\land
      w=\IntClass{r}{s}\land
      u=\IntClass{p r+q s}{p s+q r}\bigr)
  \end{aligned}%
}{IntegerMultiplicationQuotientDescent}{%
  \DeltaRow{Mengen}{z\dsep w\dsep u\dsep v\dsep
    a\dsep b\dsep c\dsep d\dsep
    a'\dsep b'\dsep c'\dsep d'\dsep
    p\dsep q\dsep r\dsep s}%
}
\begin{tabproofsplitwide}
  \proofpartwideindR[Abgeschlossenheit der Ergebniskoordinaten]{%
    a,b,c,d\in\mathbb N\vdash
    ac+bd\in\mathbb N\land ad+bc\in\mathbb N%
  }{%
    a,b,c,d\in\mathbb N\vdash
    ac+bd\in\mathbb N\land ad+bc\in\mathbb N%
  }[IntegerMultiplicationDescentCoordinates]
    \proofstepwidestar[1]{a,b,c,d\in\mathbb N}{\rA}
    \proofstepwidestar[1]{ac\in\mathbb N}{%
      \FormulaRefAuto{PeanoMulClosure}{1}}
    \proofstepwidestar[1]{bd\in\mathbb N}{%
      \FormulaRefAuto{PeanoMulClosure}{1}}
    \proofstepwidestar[1]{ad\in\mathbb N}{%
      \FormulaRefAuto{PeanoMulClosure}{1}}
    \proofstepwidestar[1]{bc\in\mathbb N}{%
      \FormulaRefAuto{PeanoMulClosure}{1}}
    \proofstepwidestar[1]{ac+bd\in\mathbb N}{%
      \FormulaRefAuto{PeanoAddClosure}{2,3}}
    \proofstepwidestar[1]{ad+bc\in\mathbb N}{%
      \FormulaRefAuto{PeanoAddClosure}{4,5}}
    \proofstepwidestar[1]{%
      ac+bd\in\mathbb N\land ad+bc\in\mathbb N}{\rAI{6,7}}
  \closeproofpartwideindR

  \proofpartwideindR[Ergebniskandidat fester Repräsentanten]{%
    \begin{aligned}[t]
      &a,b,c,d\in\mathbb N\dsep
      z=\IntClass{a}{b}\dsep w=\IntClass{c}{d}\vdash{}\\
      &\IntClass{ac+bd}{ad+bc}\in\mathbb Z\land
      \exists p\in\mathbb N\,\exists q\in\mathbb N\,
      \exists r\in\mathbb N\,\exists s\in\mathbb N\,{}\\
      &\quad
      \bigl(z=\IntClass{p}{q}\land
        w=\IntClass{r}{s}\land
        \IntClass{ac+bd}{ad+bc}
          =\IntClass{p r+q s}{p s+q r}\bigr)
    \end{aligned}%
  }{%
    a,b,c,d\in\mathbb N\dsep
    z=\IntClass{a}{b}\dsep w=\IntClass{c}{d}\vdash
    \IntClass{ac+bd}{ad+bc}\in\mathbb Z\land
    \exists p\in\mathbb N\,\exists q\in\mathbb N\,
    \exists r\in\mathbb N\,\exists s\in\mathbb N\,
    \bigl(z=\IntClass{p}{q}\land
      w=\IntClass{r}{s}\land
      \IntClass{ac+bd}{ad+bc}
        =\IntClass{p r+q s}{p s+q r}\bigr)%
  }[IntegerMultiplicationDescentCandidate]
    \proofstepwidestar[1]{a,b,c,d\in\mathbb N}{\rA}
    \proofstepwidestar[2]{z=\IntClass{a}{b}}{\rA}
    \proofstepwidestar[3]{w=\IntClass{c}{d}}{\rA}
    \proofstepwidestar[1]{%
      ac+bd\in\mathbb N\land ad+bc\in\mathbb N}{%
      \FormulaRefAuto{IntegerMultiplicationDescentCoordinates}{1}}
    \proofstepwidestar[1]{\IntClass{ac+bd}{ad+bc}\in\mathbb Z}{%
      \FormulaRefAuto{IntegerClassInZ}{\rAEa{4},\rAEb{4}}}
    \proofstepwidestar[]{%
      \IntClass{ac+bd}{ad+bc}=\IntClass{ac+bd}{ad+bc}}{\rII}
    \proofstepwidestar[1,2,3]{%
      \exists p\in\mathbb N\,\exists q\in\mathbb N\,
      \exists r\in\mathbb N\,\exists s\in\mathbb N\,
      \bigl(z=\IntClass{p}{q}\land
        w=\IntClass{r}{s}\land
        \IntClass{ac+bd}{ad+bc}
          =\IntClass{p r+q s}{p s+q r}\bigr)}{%
      \rEI{\rAI{1,\rAI{2,\rAI{3,6}}}}}
    \proofstepwidestar[1,2,3]{%
      \IntClass{ac+bd}{ad+bc}\in\mathbb Z\land
      \exists p\in\mathbb N\,\exists q\in\mathbb N\,
      \exists r\in\mathbb N\,\exists s\in\mathbb N\,
      \bigl(z=\IntClass{p}{q}\land
        w=\IntClass{r}{s}\land
        \IntClass{ac+bd}{ad+bc}
          =\IntClass{p r+q s}{p s+q r}\bigr)}{\rAI{5,7}}
  \closeproofpartwideindR

  \proofpartwideindR[Existenz durch getrennte Repräsentantenwahl]{%
    \begin{aligned}[t]
      &z,w\in\mathbb Z\vdash
      \exists u\in\mathbb Z\,
      \exists p\in\mathbb N\,\exists q\in\mathbb N\,
      \exists r\in\mathbb N\,\exists s\in\mathbb N\,{}\\
      &\qquad
      \bigl(z=\IntClass{p}{q}\land
        w=\IntClass{r}{s}\land
        u=\IntClass{p r+q s}{p s+q r}\bigr)
    \end{aligned}%
  }{%
    z,w\in\mathbb Z\vdash
    \exists u\in\mathbb Z\,
    \exists p\in\mathbb N\,\exists q\in\mathbb N\,
    \exists r\in\mathbb N\,\exists s\in\mathbb N\,
    \bigl(z=\IntClass{p}{q}\land
      w=\IntClass{r}{s}\land
      u=\IntClass{p r+q s}{p s+q r}\bigr)%
  }[IntegerMultiplicationDescentExistence]
    \proofstepwidestar[1]{z,w\in\mathbb Z}{\rA}
    \proofstepwidestar[1]{z\in\mathbb Z}{\rAEa{1}}
    \proofstepwidestar[1]{w\in\mathbb Z}{\rAEb{1}}
    \proofstepwidestar[1]{%
      \exists a\in\mathbb N\,\exists b\in\mathbb N\,
      z=\IntClass{a}{b}}{%
      \FormulaRefAuto{IntegerPairRepresentative}{2}}
    \proofstepwidestar[1]{%
      \exists c\in\mathbb N\,\exists d\in\mathbb N\,
      w=\IntClass{c}{d}}{%
      \FormulaRefAuto{IntegerPairRepresentative}{3}}
    \proofstepwidestar[6]{%
      a,b\in\mathbb N\land z=\IntClass{a}{b}}{\rA}
    \proofstepwidestar[7]{%
      c,d\in\mathbb N\land w=\IntClass{c}{d}}{\rA}
    \proofstepwidestar[6]{a,b\in\mathbb N}{\rAEa{6}}
    \proofstepwidestar[6]{z=\IntClass{a}{b}}{\rAEb{6}}
    \proofstepwidestar[7]{c,d\in\mathbb N}{\rAEa{7}}
    \proofstepwidestar[7]{w=\IntClass{c}{d}}{\rAEb{7}}
    \proofstepwidestar[6,7]{a,b,c,d\in\mathbb N}{\rAI{8,10}}
    \proofstepwidestar[6,7]{%
      \IntClass{ac+bd}{ad+bc}\in\mathbb Z\land
      \exists p\in\mathbb N\,\exists q\in\mathbb N\,
      \exists r\in\mathbb N\,\exists s\in\mathbb N\,
      \bigl(z=\IntClass{p}{q}\land
        w=\IntClass{r}{s}\land
        \IntClass{ac+bd}{ad+bc}
          =\IntClass{p r+q s}{p s+q r}\bigr)}{%
      \FormulaRefAuto{IntegerMultiplicationDescentCandidate}{12,9,11}}
    \proofstepwidestar[6,7]{%
      \exists u\in\mathbb Z\,
      \exists p\in\mathbb N\,\exists q\in\mathbb N\,
      \exists r\in\mathbb N\,\exists s\in\mathbb N\,
      \bigl(z=\IntClass{p}{q}\land
        w=\IntClass{r}{s}\land
        u=\IntClass{p r+q s}{p s+q r}\bigr)}{\rEI{13}}
    \proofstepwidestar[6]{%
      \exists u\in\mathbb Z\,
      \exists p\in\mathbb N\,\exists q\in\mathbb N\,
      \exists r\in\mathbb N\,\exists s\in\mathbb N\,
      \bigl(z=\IntClass{p}{q}\land
        w=\IntClass{r}{s}\land
        u=\IntClass{p r+q s}{p s+q r}\bigr)}{\rEE{5,7,14}}
    \proofstepwidestar[1]{%
      \exists u\in\mathbb Z\,
      \exists p\in\mathbb N\,\exists q\in\mathbb N\,
      \exists r\in\mathbb N\,\exists s\in\mathbb N\,
      \bigl(z=\IntClass{p}{q}\land
        w=\IntClass{r}{s}\land
        u=\IntClass{p r+q s}{p s+q r}\bigr)}{\rEE{4,6,15}}
  \closeproofpartwideindR

  \proofpartwideindR[Eindeutigkeit für feste Repräsentanten]{%
    \begin{aligned}[t]
      &a,b,c,d,a',b',c',d'\in\mathbb N\dsep{}\\
      &z=\IntClass{a}{b}\land w=\IntClass{c}{d}
        \land u=\IntClass{ac+bd}{ad+bc}\dsep{}\\
      &z=\IntClass{a'}{b'}\land w=\IntClass{c'}{d'}
        \land v=\IntClass{a'c'+b'd'}{a'd'+b'c'}
      \vdash u=v
    \end{aligned}%
  }{%
    a,b,c,d,a',b',c',d'\in\mathbb N\dsep
    z=\IntClass{a}{b}\land w=\IntClass{c}{d}
      \land u=\IntClass{ac+bd}{ad+bc}\dsep
    z=\IntClass{a'}{b'}\land w=\IntClass{c'}{d'}
      \land v=\IntClass{a'c'+b'd'}{a'd'+b'c'}
    \vdash u=v%
  }[IntegerMultiplicationDescentFixedUniqueness]
    \proofstepwidestar[1]{%
      a,b,c,d,a',b',c',d'\in\mathbb N}{\rA}
    \proofstepwidestar[2]{%
      z=\IntClass{a}{b}\land w=\IntClass{c}{d}
      \land u=\IntClass{ac+bd}{ad+bc}}{\rA}
    \proofstepwidestar[3]{%
      z=\IntClass{a'}{b'}\land w=\IntClass{c'}{d'}
      \land v=\IntClass{a'c'+b'd'}{a'd'+b'c'}}{\rA}
    \proofstepwidestar[2]{z=\IntClass{a}{b}}{\rAEa{2}}
    \proofstepwidestar[2]{w=\IntClass{c}{d}}{\rAEa{\rAEb{2}}}
    \proofstepwidestar[2]{u=\IntClass{ac+bd}{ad+bc}}{\rAEb{\rAEb{2}}}
    \proofstepwidestar[3]{z=\IntClass{a'}{b'}}{\rAEa{3}}
    \proofstepwidestar[3]{w=\IntClass{c'}{d'}}{\rAEa{\rAEb{3}}}
    \proofstepwidestar[3]{%
      v=\IntClass{a'c'+b'd'}{a'd'+b'c'}}{\rAEb{\rAEb{3}}}
    \proofstepwidestar[2,3]{%
      \IntClass{a}{b}=\IntClass{a'}{b'}}{\rIE{4,7}}
    \proofstepwidestar[2,3]{%
      \IntClass{c}{d}=\IntClass{c'}{d'}}{\rIE{5,8}}
    \proofstepwidestar[1,2,3]{%
      \IntClass{ac+bd}{ad+bc}
        =\IntClass{a'c'+b'd'}{a'd'+b'c'}}{%
      \FormulaRefAuto{IntegerMultiplicationCompatible}{1,10,11}}
    \proofstepwidestar[1,2,3]{u=v}{\rIE{6,9,12}}
  \closeproofpartwideindR

  \proofpartwideindR[Eindeutigkeit beliebiger Ergebniszeugen]{%
    \begin{aligned}[t]
      &\bigl(u\in\mathbb Z\land
      \exists a,b,c,d\in\mathbb N\,
      (z=\IntClass{a}{b}\land w=\IntClass{c}{d}
       \land u=\IntClass{ac+bd}{ad+bc})\bigr)\dsep{}\\
      &\bigl(v\in\mathbb Z\land
      \exists a',b',c',d'\in\mathbb N\,
      (z=\IntClass{a'}{b'}\land w=\IntClass{c'}{d'}
       \land v=\IntClass{a'c'+b'd'}{a'd'+b'c'})\bigr)
      \vdash u=v
    \end{aligned}%
  }{%
    \bigl(u\in\mathbb Z\land
    \exists a,b,c,d\in\mathbb N\,
    (z=\IntClass{a}{b}\land w=\IntClass{c}{d}
     \land u=\IntClass{ac+bd}{ad+bc})\bigr)\dsep
    \bigl(v\in\mathbb Z\land
    \exists a',b',c',d'\in\mathbb N\,
    (z=\IntClass{a'}{b'}\land w=\IntClass{c'}{d'}
     \land v=\IntClass{a'c'+b'd'}{a'd'+b'c'})\bigr)
    \vdash u=v%
  }[IntegerMultiplicationDescentUniqueness]
    \proofstepwidestar[1]{%
      u\in\mathbb Z\land
      \exists a,b,c,d\in\mathbb N\,
      (z=\IntClass{a}{b}\land w=\IntClass{c}{d}
       \land u=\IntClass{ac+bd}{ad+bc})}{\rA}
    \proofstepwidestar[2]{%
      v\in\mathbb Z\land
      \exists a',b',c',d'\in\mathbb N\,
      (z=\IntClass{a'}{b'}\land w=\IntClass{c'}{d'}
       \land v=\IntClass{a'c'+b'd'}{a'd'+b'c'})}{\rA}
    \proofstepwidestar[1]{%
      \exists a,b,c,d\in\mathbb N\,
      (z=\IntClass{a}{b}\land w=\IntClass{c}{d}
       \land u=\IntClass{ac+bd}{ad+bc})}{\rAEb{1}}
    \proofstepwidestar[2]{%
      \exists a',b',c',d'\in\mathbb N\,
      (z=\IntClass{a'}{b'}\land w=\IntClass{c'}{d'}
       \land v=\IntClass{a'c'+b'd'}{a'd'+b'c'})}{\rAEb{2}}
    \proofstepwidestar[5]{%
      a,b,c,d\in\mathbb N\land
      z=\IntClass{a}{b}\land w=\IntClass{c}{d}
      \land u=\IntClass{ac+bd}{ad+bc}}{\rA}
    \proofstepwidestar[6]{%
      a',b',c',d'\in\mathbb N\land
      z=\IntClass{a'}{b'}\land w=\IntClass{c'}{d'}
      \land v=\IntClass{a'c'+b'd'}{a'd'+b'c'}}{\rA}
    \proofstepwidestar[5]{a,b,c,d\in\mathbb N}{\rAEa{5}}
    \proofstepwidestar[5]{%
      z=\IntClass{a}{b}\land w=\IntClass{c}{d}
      \land u=\IntClass{ac+bd}{ad+bc}}{\rAEb{5}}
    \proofstepwidestar[6]{a',b',c',d'\in\mathbb N}{\rAEa{6}}
    \proofstepwidestar[6]{%
      z=\IntClass{a'}{b'}\land w=\IntClass{c'}{d'}
      \land v=\IntClass{a'c'+b'd'}{a'd'+b'c'}}{\rAEb{6}}
    \proofstepwidestar[5,6]{%
      a,b,c,d,a',b',c',d'\in\mathbb N}{\rAI{7,9}}
    \proofstepwidestar[5,6]{u=v}{%
      \FormulaRefAuto{IntegerMultiplicationDescentFixedUniqueness}{%
        11,8,10}}
    \proofstepwidestar[5]{u=v}{\rEE{4,6,12}}
    \proofstepwidestar[1,2]{u=v}{\rEE{3,5,13}}
  \closeproofpartwideindR

  \proofpartwideindR[Abschluss des Abstiegs]{%
    \begin{aligned}[t]
      &z,w\in\mathbb Z\vdash
      \exists! u\in\mathbb Z\,
      \exists p\in\mathbb N\,\exists q\in\mathbb N\,
      \exists r\in\mathbb N\,\exists s\in\mathbb N\,{}\\
      &\qquad
      \bigl(z=\IntClass{p}{q}\land w=\IntClass{r}{s}
      \land u=\IntClass{p r+q s}{p s+q r}\bigr)
    \end{aligned}%
  }{%
    z,w\in\mathbb Z\vdash
    \exists! u\in\mathbb Z\,
    \exists p\in\mathbb N\,\exists q\in\mathbb N\,
    \exists r\in\mathbb N\,\exists s\in\mathbb N\,
    \bigl(z=\IntClass{p}{q}\land w=\IntClass{r}{s}
    \land u=\IntClass{p r+q s}{p s+q r}\bigr)%
  }[IntegerMultiplicationDescentConclusion]
    \proofstepwidestar[1]{z,w\in\mathbb Z}{\rA}
    \proofstepwidestar[1]{%
      \exists u\in\mathbb Z\,
      \exists p\in\mathbb N\,\exists q\in\mathbb N\,
      \exists r\in\mathbb N\,\exists s\in\mathbb N\,
      (z=\IntClass{p}{q}\land w=\IntClass{r}{s}
       \land u=\IntClass{p r+q s}{p s+q r})}{%
      \FormulaRefAuto{IntegerMultiplicationDescentExistence}{1}}
    \proofstepwidestar[3]{%
      u\in\mathbb Z\land
      \exists p\in\mathbb N\,\exists q\in\mathbb N\,
      \exists r\in\mathbb N\,\exists s\in\mathbb N\,
      (z=\IntClass{p}{q}\land w=\IntClass{r}{s}
       \land u=\IntClass{p r+q s}{p s+q r})}{\rA}
    \proofstepwidestar[4]{%
      v\in\mathbb Z\land
      \exists p\in\mathbb N\,\exists q\in\mathbb N\,
      \exists r\in\mathbb N\,\exists s\in\mathbb N\,
      (z=\IntClass{p}{q}\land w=\IntClass{r}{s}
       \land v=\IntClass{p r+q s}{p s+q r})}{\rA}
    \proofstepwidestar[3,4]{u=v}{%
      \FormulaRefAuto{IntegerMultiplicationDescentUniqueness}{3,4}}
    \proofstepwidestar[1]{%
      \exists! u\in\mathbb Z\,
      \exists p\in\mathbb N\,\exists q\in\mathbb N\,
      \exists r\in\mathbb N\,\exists s\in\mathbb N\,
      (z=\IntClass{p}{q}\land w=\IntClass{r}{s}
       \land u=\IntClass{p r+q s}{p s+q r})}{\UEI{2,3,4,5}}
  \closeproofpartwideindR
\end{tabproofsplitwide}

Der Quotientenabstieg liefert damit auch für die Multiplikation unabhängig
von beiden gewählten Differenzenpaaren genau ein Ergebnis in \(\mathbb Z\).

\FormulaDefDeltaK[Multiplikation ganzer Zahlen]{%
  a,b,c,d\in\mathbb N\vdash
  \IntClass{a}{b}\cdot\IntClass{c}{d}
  \coloneqq
  \IntClass{a\cdot c+b\cdot d}{a\cdot d+b\cdot c}%
}{IntegerMultiplicationDef}{%
  \DeltaRow{Mengen}{a\dsep b\dsep c\dsep d}%
  \DeltaRow{Neue Symbole}{\cdot}%
}

\FormulaThmDeltaK[Abgeschlossenheit der Multiplikation ganzer Zahlen]{%
  z\in\mathbb Z\dsep w\in\mathbb Z
  \vdash z\cdot w\in\mathbb Z%
}{IntMulClosure}{%
  \DeltaRow{Mengen}{z\dsep w\dsep a\dsep b\dsep c\dsep d}%
}
\begin{tabproofsplitwide}
  \proofpartwideindR[Abgeschlossenheit für feste Repräsentanten]{%
    \begin{aligned}[t]
      &a,b,c,d\in\mathbb N\land
        z=\IntClass{a}{b}\land w=\IntClass{c}{d}\\
      &\vdash z\cdot w\in\mathbb Z
    \end{aligned}
  }{%
    a,b,c,d\in\mathbb N\land
    z=\IntClass{a}{b}\land w=\IntClass{c}{d}
    \vdash z\cdot w\in\mathbb Z%
  }[IntMulClosureForRepresentatives]
    \proofstepwidestar[1]{a,b,c,d\in\mathbb N\land
      z=\IntClass{a}{b}\land w=\IntClass{c}{d}}{\rA}
    \proofstepwidestar[1]{%
      z\cdot w=\IntClass{ac+bd}{ad+bc}}{%
      \FormulaRefAuto{IntegerMultiplicationDef}{1}}
    \proofstepwidestar[1]{ac+bd\in\mathbb N}{%
      \FormulaRefAuto{PeanoAddClosure}{%
        \FormulaRefAuto{PeanoMulClosure}{1},
        \FormulaRefAuto{PeanoMulClosure}{1}}}
    \proofstepwidestar[1]{ad+bc\in\mathbb N}{%
      \FormulaRefAuto{PeanoAddClosure}{%
        \FormulaRefAuto{PeanoMulClosure}{1},
        \FormulaRefAuto{PeanoMulClosure}{1}}}
    \proofstepwidestar[1]{\IntClass{ac+bd}{ad+bc}\in\mathbb Z}{%
      \FormulaRefAuto{IntegerClassInZ}{3,4}}
    \proofstepwidestar[1]{z\cdot w\in\mathbb Z}{\rIE{2,5}}
  \closeproofpartwideindR

  \proofpartwideindR[Elimination der Repräsentanten]{%
    z\in\mathbb Z\dsep w\in\mathbb Z
    \vdash z\cdot w\in\mathbb Z%
  }{%
    z\in\mathbb Z\dsep w\in\mathbb Z
    \vdash z\cdot w\in\mathbb Z%
  }[IntMulClosureRepresentativeElimination]
    \proofstepwidestar[1]{z\in\mathbb Z}{\rA}
    \proofstepwidestar[2]{w\in\mathbb Z}{\rA}
    \proofstepwidestar[1]{\exists a,b\in\mathbb N\,
      z=\IntClass{a}{b}}{\FormulaRefAuto{IntegerPairRepresentative}{1}}
    \proofstepwidestar[2]{\exists c,d\in\mathbb N\,
      w=\IntClass{c}{d}}{\FormulaRefAuto{IntegerPairRepresentative}{2}}
    \proofstepwidestar[5]{a,b,c,d\in\mathbb N\land
      z=\IntClass{a}{b}\land w=\IntClass{c}{d}}{\rA}
    \proofstepwidestar[5]{z\cdot w\in\mathbb Z}{%
      \FormulaRefAuto{IntMulClosureForRepresentatives}{5}}
    \proofstepwidestar[1,2]{z\cdot w\in\mathbb Z}{\rEE{3,4,5,6}}
  \closeproofpartwideindR
\end{tabproofsplitwide}

\FormulaThmDeltaK[Kommutativität der Multiplikation ganzer Zahlen]{%
  z,w\in\mathbb Z\vdash z\cdot w=w\cdot z%
}{IntMulCommutative}{%
  \DeltaRow{Mengen}{z\dsep w\dsep a\dsep b\dsep c\dsep d}%
}
\begin{tabproofsplitwide}
  \proofpartwideindR[Produkte gewählter Repräsentanten]{%
    \begin{aligned}[t]
      &a,b,c,d\in\mathbb N\land
        z=\IntClass{a}{b}\land w=\IntClass{c}{d}\\
      &\vdash
        z\cdot w=\IntClass{ac+bd}{ad+bc}\land
        w\cdot z=\IntClass{ca+db}{cb+da}
    \end{aligned}
  }{%
    a,b,c,d\in\mathbb N\land
    z=\IntClass{a}{b}\land w=\IntClass{c}{d}
    \vdash
    z\cdot w=\IntClass{ac+bd}{ad+bc}\land
    w\cdot z=\IntClass{ca+db}{cb+da}
  }[IntMulCommutativeRepresentativeProducts]
    \proofstepwidestar[1]{a,b,c,d\in\mathbb N\land
      z=\IntClass{a}{b}\land w=\IntClass{c}{d}}{\rA}
    \proofstepwidestar[1]{z\cdot w=\IntClass{ac+bd}{ad+bc}}{%
      \FormulaRefAuto{IntegerMultiplicationDef}{1}}
    \proofstepwidestar[1]{w\cdot z=\IntClass{ca+db}{cb+da}}{%
      \FormulaRefAuto{IntegerMultiplicationDef}{1}}
    \proofstepwidestar[1]{%
      z\cdot w=\IntClass{ac+bd}{ad+bc}\land
      w\cdot z=\IntClass{ca+db}{cb+da}}{\rAI{2,3}}
  \closeproofpartwideindR

  \proofpartwideindR[Vertauschung der Klassenkoordinaten]{%
    a,b,c,d\in\mathbb N\vdash
    \IntClass{ac+bd}{ad+bc}=\IntClass{ca+db}{cb+da}%
  }{%
    a,b,c,d\in\mathbb N\vdash
    \IntClass{ac+bd}{ad+bc}=\IntClass{ca+db}{cb+da}%
  }[IntMulCommutativeClassCoordinates]
    \proofstepwidestar[1]{a,b,c,d\in\mathbb N}{\rA}
    \proofstepwidestar[1]{ac=ca}{%
      \FormulaRefAuto{PeanoMulCommutative}{1}}
    \proofstepwidestar[1]{bd=db}{%
      \FormulaRefAuto{PeanoMulCommutative}{1}}
    \proofstepwidestar[1]{ad=da}{%
      \FormulaRefAuto{PeanoMulCommutative}{1}}
    \proofstepwidestar[1]{bc=cb}{%
      \FormulaRefAuto{PeanoMulCommutative}{1}}
    \proofstepwidestar[1]{ac+bd=ca+db}{\rIE{2,3}}
    \proofstepwidestar[1]{ad+bc=cb+da}{%
      \FormulaRefAuto{PeanoAddCommutative}{\rIE{4,5}}}
    \proofstepwidestar[1]{\IntClass{ac+bd}{ad+bc}=
      \IntClass{ca+db}{cb+da}}{\rIE{6,7}}
  \closeproofpartwideindR

  \proofpartwideindR[Elimination der Repräsentanten]{%
    z,w\in\mathbb Z\vdash z\cdot w=w\cdot z%
  }{%
    z,w\in\mathbb Z\vdash z\cdot w=w\cdot z%
  }[IntMulCommutativeRepresentativeElimination]
    \proofstepwidestar[1]{z,w\in\mathbb Z}{\rA}
    \proofstepwidestar[1]{z\in\mathbb Z}{\rAEa{1}}
    \proofstepwidestar[1]{w\in\mathbb Z}{\rAEb{1}}
    \proofstepwidestar[1]{\exists a,b\in\mathbb N\,
      z=\IntClass{a}{b}}{\FormulaRefAuto{IntegerPairRepresentative}{2}}
    \proofstepwidestar[1]{\exists c,d\in\mathbb N\,
      w=\IntClass{c}{d}}{\FormulaRefAuto{IntegerPairRepresentative}{3}}
    \proofstepwidestar[6]{a,b,c,d\in\mathbb N\land
      z=\IntClass{a}{b}\land w=\IntClass{c}{d}}{\rA}
    \proofstepwidestar[6]{%
      z\cdot w=\IntClass{ac+bd}{ad+bc}\land
      w\cdot z=\IntClass{ca+db}{cb+da}}{%
      \FormulaRefAuto{IntMulCommutativeRepresentativeProducts}{6}}
    \proofstepwidestar[6]{%
      z\cdot w=\IntClass{ac+bd}{ad+bc}}{\rAEa{7}}
    \proofstepwidestar[6]{%
      w\cdot z=\IntClass{ca+db}{cb+da}}{\rAEb{7}}
    \proofstepwidestar[6]{\IntClass{ac+bd}{ad+bc}=
      \IntClass{ca+db}{cb+da}}{%
      \FormulaRefAuto{IntMulCommutativeClassCoordinates}{6}}
    \proofstepwidestar[6]{z\cdot w=w\cdot z}{\rIE{8,9,10}}
    \proofstepwidestar[1]{z\cdot w=w\cdot z}{\rEE{4,5,6,11}}
  \closeproofpartwideindR
\end{tabproofsplitwide}

\FormulaThmDeltaK[Einheitsgesetze der Multiplikation ganzer Zahlen]{%
  z\in\mathbb Z\vdash
  z\cdot1_{\mathbb Z}=z\land1_{\mathbb Z}\cdot z=z%
}{IntMulOne}{%
  \DeltaRow{Mengen}{z\dsep a\dsep b}%
}
\begin{tabproofsplitwide}
  \proofpartwideindR[Rechte Einheit am Repräsentanten]{%
    a,b\in\mathbb N\land z=\IntClass{a}{b}
    \vdash z\cdot1_{\mathbb Z}=z%
  }{%
    a,b\in\mathbb N\land z=\IntClass{a}{b}
    \vdash z\cdot1_{\mathbb Z}=z%
  }[IntMulRightIdentityRepresentative]
    \proofstepwidestar[1]{%
      a,b\in\mathbb N\land z=\IntClass{a}{b}}{\rA}
    \proofstepwidestar[]{1_{\mathbb Z}=\IntEmbed(1)}{%
      \FormulaRefAuto{IntegerOneDef}}
    \proofstepwidestar[]{\IntEmbed(1)=\IntClass{1}{0}}{%
      \FormulaRefAuto{NaturalIntegerEmbedding}{%
        \FormulaRefAuto{PeanoOneInN}}}
    \proofstepwidestar[1]{z\cdot1_{\mathbb Z}=
      \IntClass{a\cdot1+b\cdot0}{a\cdot0+b\cdot1}}{%
      \FormulaRefAuto{IntegerMultiplicationDef}{1,2,3}}
    \proofstepwidestar[1]{a\cdot1=a}{%
      \FormulaRefAuto{PeanoMulRightOne}{1}}
    \proofstepwidestar[1]{b\cdot1=b}{%
      \FormulaRefAuto{PeanoMulRightOne}{1}}
    \proofstepwidestar[1]{a\cdot0=0}{%
      \FormulaRefAuto{PeanoMulRightZero}{1}}
    \proofstepwidestar[1]{b\cdot0=0}{%
      \FormulaRefAuto{PeanoMulRightZero}{1}}
    \proofstepwidestar[1]{a\cdot1+b\cdot0=a}{%
      \rIE{5,8,\FormulaRefAuto{PeanoAddRightZero}{1}}}
    \proofstepwidestar[1]{a\cdot0+b\cdot1=b}{%
      \rIE{7,6,\FormulaRefAuto{PeanoAddLeftZero}{1}}}
    \proofstepwidestar[1]{z\cdot1_{\mathbb Z}=z}{%
      \rIE{1,4,9,10}}
  \closeproofpartwideindR

  \proofpartwideindR[Beidseitige Einheit]{%
    z\in\mathbb Z\vdash
    z\cdot1_{\mathbb Z}=z\land1_{\mathbb Z}\cdot z=z%
  }{%
    z\in\mathbb Z\vdash
    z\cdot1_{\mathbb Z}=z\land1_{\mathbb Z}\cdot z=z%
  }[IntMulTwoSidedIdentityConclusion]
    \proofstepwidestar[1]{z\in\mathbb Z}{\rA}
    \proofstepwidestar[1]{\exists a,b\in\mathbb N\,
      z=\IntClass{a}{b}}{%
      \FormulaRefAuto{IntegerPairRepresentative}{1}}
    \proofstepwidestar[3]{%
      a,b\in\mathbb N\land z=\IntClass{a}{b}}{\rA}
    \proofstepwidestar[3]{z\cdot1_{\mathbb Z}=z}{%
      \FormulaRefAuto{IntMulRightIdentityRepresentative}{3}}
    \proofstepwidestar[3]{1_{\mathbb Z}\cdot z=z}{%
      \FormulaRefAuto{IntMulCommutative}{%
        \FormulaRefAuto{IntOneCarrier},1,
        \FormulaRefAuto{a=b\vdash b=a}{4}}}
    \proofstepwidestar[3]{z\cdot1_{\mathbb Z}=z\land
      1_{\mathbb Z}\cdot z=z}{\rAI{4,5}}
    \proofstepwidestar[1]{z\cdot1_{\mathbb Z}=z\land
      1_{\mathbb Z}\cdot z=z}{\rEE{2,3,6}}
  \closeproofpartwideindR
\end{tabproofsplitwide}

\FormulaThmDeltaK[Assoziativität der Multiplikation ganzer Zahlen]{%
  z,w,u\in\mathbb Z\vdash
  (z\cdot w)\cdot u=z\cdot(w\cdot u)%
}{IntMulAssociative}{%
  \DeltaRow{Mengen}{z\dsep w\dsep u\dsep
    a\dsep b\dsep c\dsep d\dsep e\dsep f}%
}
\begin{tabproofsplitwide}
  \proofpartwideindR[Positive Koordinate: linke Seite ausmultiplizieren]{%
    \begin{aligned}[t]
      &a,b,c,d,e,f\in\mathbb N\vdash{}\\
      &(ac+bd)e+(ad+bc)f
      =\bigl((ac)e+(bd)e\bigr)+\bigl((ad)f+(bc)f\bigr)
    \end{aligned}
  }{%
    a,b,c,d,e,f\in\mathbb N\vdash
    (ac+bd)e+(ad+bc)f
    =\bigl((ac)e+(bd)e\bigr)+\bigl((ad)f+(bc)f\bigr)
  }[IntMulAssociativePositiveLeftExpansion]
    \proofstepwidestar[1]{a,b,c,d,e,f\in\mathbb N}{\rA}
    \proofstepwide[1]{(ac+bd)e}{=}{(ac)e+(bd)e}{%
      \FormulaRefAuto{PeanoMulRightDistributive}{1}}
    \proofstepwide[1]{(ad+bc)f}{=}{(ad)f+(bc)f}{%
      \FormulaRefAuto{PeanoMulRightDistributive}{1}}
    \proofstepwidestar[1]{%
      (ac+bd)e+(ad+bc)f
      =\bigl((ac)e+(bd)e\bigr)+\bigl((ad)f+(bc)f\bigr)}{\rIE{2,3}}
  \closeproofpartwideindR

  \proofpartwideindR[Positive Koordinate: rechte Seite ausmultiplizieren]{%
    \begin{aligned}[t]
      &a,b,c,d,e,f\in\mathbb N\vdash{}\\
      &a(ce+df)+b(cf+de)
      =\bigl(a(ce)+a(df)\bigr)+\bigl(b(cf)+b(de)\bigr)
    \end{aligned}
  }{%
    a,b,c,d,e,f\in\mathbb N\vdash
    a(ce+df)+b(cf+de)
    =\bigl(a(ce)+a(df)\bigr)+\bigl(b(cf)+b(de)\bigr)
  }[IntMulAssociativePositiveRightExpansion]
    \proofstepwidestar[1]{a,b,c,d,e,f\in\mathbb N}{\rA}
    \proofstepwide[1]{a(ce+df)}{=}{a(ce)+a(df)}{%
      \FormulaRefAuto{PeanoMulLeftDistributive}{1}}
    \proofstepwide[1]{b(cf+de)}{=}{b(cf)+b(de)}{%
      \FormulaRefAuto{PeanoMulLeftDistributive}{1}}
    \proofstepwidestar[1]{%
      a(ce+df)+b(cf+de)
      =\bigl(a(ce)+a(df)\bigr)+\bigl(b(cf)+b(de)\bigr)}{\rIE{2,3}}
  \closeproofpartwideindR

  \proofpartwideindR[Positive Koordinate: ausmultiplizierte Terme ordnen]{%
    \begin{aligned}[t]
      &a,b,c,d,e,f\in\mathbb N\vdash{}\\
      &\bigl((ac)e+(bd)e\bigr)+\bigl((ad)f+(bc)f\bigr)\\
      &\qquad=
      \bigl(a(ce)+a(df)\bigr)+\bigl(b(cf)+b(de)\bigr)
    \end{aligned}
  }{%
    \begin{aligned}[t]
      &a,b,c,d,e,f\in\mathbb N\vdash{}\\
      &\bigl((ac)e+(bd)e\bigr)+\bigl((ad)f+(bc)f\bigr)
      =
      \bigl(a(ce)+a(df)\bigr)+\bigl(b(cf)+b(de)\bigr)
    \end{aligned}
  }[IntMulAssociativePositiveExpandedComparison]
    \proofstepwidestar[1]{a,b,c,d,e,f\in\mathbb N}{\rA}
    \proofstepwidestar[1]{(ac)e=a(ce)}{%
      \FormulaRefAuto{PeanoMulAssociative}{1}}
    \proofstepwidestar[1]{(bd)e=b(de)}{%
      \FormulaRefAuto{PeanoMulAssociative}{1}}
    \proofstepwidestar[1]{(ad)f=a(df)}{%
      \FormulaRefAuto{PeanoMulAssociative}{1}}
    \proofstepwidestar[1]{(bc)f=b(cf)}{%
      \FormulaRefAuto{PeanoMulAssociative}{1}}
    \proofstepwidestar[1]{%
      \bigl((ac)e+(bd)e\bigr)+\bigl((ad)f+(bc)f\bigr)
      =
      \bigl(a(ce)+b(de)\bigr)+\bigl(a(df)+b(cf)\bigr)}{%
      \rIE{2,3,4,5}}
    \proofstepwidestar[1]{%
      \bigl(a(ce)+b(de)\bigr)+\bigl(a(df)+b(cf)\bigr)
      =
      \bigl(a(ce)+a(df)\bigr)+\bigl(b(de)+b(cf)\bigr)}{%
      \FormulaRefAuto{PeanoAddInterchange}{1}}
    \proofstepwidestar[1]{b(de)+b(cf)=b(cf)+b(de)}{%
      \FormulaRefAuto{PeanoAddCommutative}{1}}
    \proofstepwidestar[1]{%
      \bigl(a(ce)+a(df)\bigr)+\bigl(b(de)+b(cf)\bigr)
      =
      \bigl(a(ce)+a(df)\bigr)+\bigl(b(cf)+b(de)\bigr)}{\rIE{8}}
    \proofstepwidestar[1]{%
      \bigl((ac)e+(bd)e\bigr)+\bigl((ad)f+(bc)f\bigr)
      =
      \bigl(a(ce)+a(df)\bigr)+\bigl(b(cf)+b(de)\bigr)}{%
      \rIE{6,7,9}}
  \closeproofpartwideindR

  \proofpartwideindR[Gleichheit der positiven Koordinaten]{%
    \begin{aligned}[t]
      &a,b,c,d,e,f\in\mathbb N\vdash{}\\
      &(ac+bd)e+(ad+bc)f
      =a(ce+df)+b(cf+de)
    \end{aligned}
  }{%
    a,b,c,d,e,f\in\mathbb N\vdash
    (ac+bd)e+(ad+bc)f
    =a(ce+df)+b(cf+de)
  }[IntMulAssociativePositiveCoordinate]
    \proofstepwidestar[1]{a,b,c,d,e,f\in\mathbb N}{\rA}
    \proofstepwidestar[1]{%
      (ac+bd)e+(ad+bc)f
      =\bigl((ac)e+(bd)e\bigr)+\bigl((ad)f+(bc)f\bigr)}{%
      \FormulaRefAuto{IntMulAssociativePositiveLeftExpansion}{1}}
    \proofstepwidestar[1]{%
      \bigl((ac)e+(bd)e\bigr)+\bigl((ad)f+(bc)f\bigr)
      =
      \bigl(a(ce)+a(df)\bigr)+\bigl(b(cf)+b(de)\bigr)}{%
      \FormulaRefAuto{IntMulAssociativePositiveExpandedComparison}{1}}
    \proofstepwidestar[1]{%
      a(ce+df)+b(cf+de)
      =\bigl(a(ce)+a(df)\bigr)+\bigl(b(cf)+b(de)\bigr)}{%
      \FormulaRefAuto{IntMulAssociativePositiveRightExpansion}{1}}
    \proofstepwidestar[1]{%
      \bigl(a(ce)+a(df)\bigr)+\bigl(b(cf)+b(de)\bigr)
      =a(ce+df)+b(cf+de)}{%
      \FormulaRefAuto{a=b\vdash b=a}{4}}
    \proofstepwidestar[1]{%
      (ac+bd)e+(ad+bc)f
      =a(ce+df)+b(cf+de)}{%
      \rIE{2,3,5}}
  \closeproofpartwideindR

  \proofpartwideindR[Negative Koordinate: linke Seite ausmultiplizieren]{%
    \begin{aligned}[t]
      &a,b,c,d,e,f\in\mathbb N\vdash{}\\
      &(ac+bd)f+(ad+bc)e
      =\bigl((ac)f+(bd)f\bigr)+\bigl((ad)e+(bc)e\bigr)
    \end{aligned}
  }{%
    a,b,c,d,e,f\in\mathbb N\vdash
    (ac+bd)f+(ad+bc)e
    =\bigl((ac)f+(bd)f\bigr)+\bigl((ad)e+(bc)e\bigr)
  }[IntMulAssociativeNegativeLeftExpansion]
    \proofstepwidestar[1]{a,b,c,d,e,f\in\mathbb N}{\rA}
    \proofstepwide[1]{(ac+bd)f}{=}{(ac)f+(bd)f}{%
      \FormulaRefAuto{PeanoMulRightDistributive}{1}}
    \proofstepwide[1]{(ad+bc)e}{=}{(ad)e+(bc)e}{%
      \FormulaRefAuto{PeanoMulRightDistributive}{1}}
    \proofstepwidestar[1]{%
      (ac+bd)f+(ad+bc)e
      =\bigl((ac)f+(bd)f\bigr)+\bigl((ad)e+(bc)e\bigr)}{\rIE{2,3}}
  \closeproofpartwideindR

  \proofpartwideindR[Negative Koordinate: rechte Seite ausmultiplizieren]{%
    \begin{aligned}[t]
      &a,b,c,d,e,f\in\mathbb N\vdash{}\\
      &a(cf+de)+b(ce+df)
      =\bigl(a(cf)+a(de)\bigr)+\bigl(b(ce)+b(df)\bigr)
    \end{aligned}
  }{%
    a,b,c,d,e,f\in\mathbb N\vdash
    a(cf+de)+b(ce+df)
    =\bigl(a(cf)+a(de)\bigr)+\bigl(b(ce)+b(df)\bigr)
  }[IntMulAssociativeNegativeRightExpansion]
    \proofstepwidestar[1]{a,b,c,d,e,f\in\mathbb N}{\rA}
    \proofstepwide[1]{a(cf+de)}{=}{a(cf)+a(de)}{%
      \FormulaRefAuto{PeanoMulLeftDistributive}{1}}
    \proofstepwide[1]{b(ce+df)}{=}{b(ce)+b(df)}{%
      \FormulaRefAuto{PeanoMulLeftDistributive}{1}}
    \proofstepwidestar[1]{%
      a(cf+de)+b(ce+df)
      =\bigl(a(cf)+a(de)\bigr)+\bigl(b(ce)+b(df)\bigr)}{\rIE{2,3}}
  \closeproofpartwideindR

  \proofpartwideindR[Negative Koordinate: ausmultiplizierte Terme ordnen]{%
    \begin{aligned}[t]
      &a,b,c,d,e,f\in\mathbb N\vdash{}\\
      &\bigl((ac)f+(bd)f\bigr)+\bigl((ad)e+(bc)e\bigr)\\
      &\qquad=
      \bigl(a(cf)+a(de)\bigr)+\bigl(b(ce)+b(df)\bigr)
    \end{aligned}
  }{%
    \begin{aligned}[t]
      &a,b,c,d,e,f\in\mathbb N\vdash{}\\
      &\bigl((ac)f+(bd)f\bigr)+\bigl((ad)e+(bc)e\bigr)
      =
      \bigl(a(cf)+a(de)\bigr)+\bigl(b(ce)+b(df)\bigr)
    \end{aligned}
  }[IntMulAssociativeNegativeExpandedComparison]
    \proofstepwidestar[1]{a,b,c,d,e,f\in\mathbb N}{\rA}
    \proofstepwidestar[1]{(ac)f=a(cf)}{%
      \FormulaRefAuto{PeanoMulAssociative}{1}}
    \proofstepwidestar[1]{(bd)f=b(df)}{%
      \FormulaRefAuto{PeanoMulAssociative}{1}}
    \proofstepwidestar[1]{(ad)e=a(de)}{%
      \FormulaRefAuto{PeanoMulAssociative}{1}}
    \proofstepwidestar[1]{(bc)e=b(ce)}{%
      \FormulaRefAuto{PeanoMulAssociative}{1}}
    \proofstepwidestar[1]{%
      \bigl((ac)f+(bd)f\bigr)+\bigl((ad)e+(bc)e\bigr)
      =
      \bigl(a(cf)+b(df)\bigr)+\bigl(a(de)+b(ce)\bigr)}{%
      \rIE{2,3,4,5}}
    \proofstepwidestar[1]{%
      \bigl(a(cf)+b(df)\bigr)+\bigl(a(de)+b(ce)\bigr)
      =
      \bigl(a(cf)+a(de)\bigr)+\bigl(b(df)+b(ce)\bigr)}{%
      \FormulaRefAuto{PeanoAddInterchange}{1}}
    \proofstepwidestar[1]{b(df)+b(ce)=b(ce)+b(df)}{%
      \FormulaRefAuto{PeanoAddCommutative}{1}}
    \proofstepwidestar[1]{%
      \bigl(a(cf)+a(de)\bigr)+\bigl(b(df)+b(ce)\bigr)
      =
      \bigl(a(cf)+a(de)\bigr)+\bigl(b(ce)+b(df)\bigr)}{\rIE{8}}
    \proofstepwidestar[1]{%
      \bigl((ac)f+(bd)f\bigr)+\bigl((ad)e+(bc)e\bigr)
      =
      \bigl(a(cf)+a(de)\bigr)+\bigl(b(ce)+b(df)\bigr)}{%
      \rIE{6,7,9}}
  \closeproofpartwideindR

  \proofpartwideindR[Gleichheit der negativen Koordinaten]{%
    \begin{aligned}[t]
      &a,b,c,d,e,f\in\mathbb N\vdash{}\\
      &(ac+bd)f+(ad+bc)e
      =a(cf+de)+b(ce+df)
    \end{aligned}
  }{%
    a,b,c,d,e,f\in\mathbb N\vdash
    (ac+bd)f+(ad+bc)e
    =a(cf+de)+b(ce+df)
  }[IntMulAssociativeNegativeCoordinate]
    \proofstepwidestar[1]{a,b,c,d,e,f\in\mathbb N}{\rA}
    \proofstepwidestar[1]{%
      (ac+bd)f+(ad+bc)e
      =\bigl((ac)f+(bd)f\bigr)+\bigl((ad)e+(bc)e\bigr)}{%
      \FormulaRefAuto{IntMulAssociativeNegativeLeftExpansion}{1}}
    \proofstepwidestar[1]{%
      \bigl((ac)f+(bd)f\bigr)+\bigl((ad)e+(bc)e\bigr)
      =
      \bigl(a(cf)+a(de)\bigr)+\bigl(b(ce)+b(df)\bigr)}{%
      \FormulaRefAuto{IntMulAssociativeNegativeExpandedComparison}{1}}
    \proofstepwidestar[1]{%
      a(cf+de)+b(ce+df)
      =\bigl(a(cf)+a(de)\bigr)+\bigl(b(ce)+b(df)\bigr)}{%
      \FormulaRefAuto{IntMulAssociativeNegativeRightExpansion}{1}}
    \proofstepwidestar[1]{%
      \bigl(a(cf)+a(de)\bigr)+\bigl(b(ce)+b(df)\bigr)
      =a(cf+de)+b(ce+df)}{%
      \FormulaRefAuto{a=b\vdash b=a}{4}}
    \proofstepwidestar[1]{%
      (ac+bd)f+(ad+bc)e
      =a(cf+de)+b(ce+df)}{%
      \rIE{2,3,5}}
  \closeproofpartwideindR

  \proofpartwideindR[Rechnung mit Klassen]{%
    \begin{aligned}[t]
      &a,b,c,d,e,f\in\mathbb N\vdash{}\\
      &(\IntClass{a}{b}\cdot\IntClass{c}{d})\cdot\IntClass{e}{f}
      =\IntClass{a}{b}\cdot
       (\IntClass{c}{d}\cdot\IntClass{e}{f})
    \end{aligned}
  }{%
    a,b,c,d,e,f\in\mathbb N\vdash
    (\IntClass{a}{b}\cdot\IntClass{c}{d})\cdot\IntClass{e}{f}
    =\IntClass{a}{b}\cdot
     (\IntClass{c}{d}\cdot\IntClass{e}{f})
  }[IntMulAssociativeOnClasses]
    \proofstepwidestar[1]{a,b,c,d,e,f\in\mathbb N}{\rA}
    \proofstepwidestar[1]{%
      (\IntClass{a}{b}\cdot\IntClass{c}{d})\cdot\IntClass{e}{f}
      =
      \IntClass{(ac+bd)e+(ad+bc)f}
               {(ac+bd)f+(ad+bc)e}}{%
      \FormulaRefAuto{IntegerMultiplicationDef}{1}}
    \proofstepwidestar[1]{%
      \IntClass{a}{b}\cdot
      (\IntClass{c}{d}\cdot\IntClass{e}{f})
      =
      \IntClass{a(ce+df)+b(cf+de)}
               {a(cf+de)+b(ce+df)}}{%
      \FormulaRefAuto{IntegerMultiplicationDef}{1}}
    \proofstepwidestar[1]{%
      (ac+bd)e+(ad+bc)f=a(ce+df)+b(cf+de)}{%
      \FormulaRefAuto{IntMulAssociativePositiveCoordinate}{1}}
    \proofstepwidestar[1]{%
      (ac+bd)f+(ad+bc)e=a(cf+de)+b(ce+df)}{%
      \FormulaRefAuto{IntMulAssociativeNegativeCoordinate}{1}}
    \proofstepwidestar[1]{%
      (\IntClass{a}{b}\cdot\IntClass{c}{d})\cdot\IntClass{e}{f}
      =\IntClass{a}{b}\cdot
       (\IntClass{c}{d}\cdot\IntClass{e}{f})}{\rIE{2,3,4,5}}
  \closeproofpartwideindR

  \proofpartwide{Schluss}
    \proofstepwidestar[1]{z,w,u\in\mathbb Z}{\rA}
    \proofstepwidestar[1]{z\in\mathbb Z}{\rAEa{1}}
    \proofstepwidestar[1]{w\in\mathbb Z}{\rAEa{\rAEb{1}}}
    \proofstepwidestar[1]{u\in\mathbb Z}{\rAEb{\rAEb{1}}}
    \proofstepwidestar[1]{\exists a,b\in\mathbb N\,
      z=\IntClass{a}{b}}{\FormulaRefAuto{IntegerPairRepresentative}{2}}
    \proofstepwidestar[1]{\exists c,d\in\mathbb N\,
      w=\IntClass{c}{d}}{\FormulaRefAuto{IntegerPairRepresentative}{3}}
    \proofstepwidestar[1]{\exists e,f\in\mathbb N\,
      u=\IntClass{e}{f}}{\FormulaRefAuto{IntegerPairRepresentative}{4}}
    \proofstepwidestar[8]{a,b,c,d,e,f\in\mathbb N\land
      z=\IntClass{a}{b}\land w=\IntClass{c}{d}\land
      u=\IntClass{e}{f}}{\rA}
    \proofstepwidestar[8]{%
      (\IntClass{a}{b}\cdot\IntClass{c}{d})\cdot\IntClass{e}{f}
      =\IntClass{a}{b}\cdot
       (\IntClass{c}{d}\cdot\IntClass{e}{f})}{%
      \FormulaRefAuto{IntMulAssociativeOnClasses}{8}}
    \proofstepwidestar[8]{%
      (z\cdot w)\cdot u=z\cdot(w\cdot u)}{\rIE{8,9}}
    \proofstepwidestar[1]{(z\cdot w)\cdot u=z\cdot(w\cdot u)}{%
      \rEE{5,6,7,8,10}}
  \closeproofpartwide
\end{tabproofsplitwide}

\FormulaThmDeltaK[Linksdistributivität ganzer Zahlen]{%
  z,w,u\in\mathbb Z\vdash
  z\cdot(w+u)=z\cdot w+z\cdot u%
}{IntMulLeftDistributive}{%
  \DeltaRow{Mengen}{z\dsep w\dsep u\dsep
    a\dsep b\dsep c\dsep d\dsep e\dsep f}%
}
\begin{tabproofsplitwide}
  \proofpartwideindR[Positive Koordinate distributiv ausmultiplizieren]{%
    \begin{aligned}[t]
      &a,b,c,d,e,f\in\mathbb N\vdash{}\\
      &a(c+e)+b(d+f)=(ac+bd)+(ae+bf)
    \end{aligned}
  }{%
    a,b,c,d,e,f\in\mathbb N\vdash
    a(c+e)+b(d+f)=(ac+bd)+(ae+bf)
  }[IntMulLeftDistributivePositiveCoordinate]
    \proofstepwidestar[1]{a,b,c,d,e,f\in\mathbb N}{\rA}
    \proofstepwidestar[1]{a(c+e)=ac+ae}{%
      \FormulaRefAuto{PeanoMulLeftDistributive}{1}}
    \proofstepwidestar[1]{b(d+f)=bd+bf}{%
      \FormulaRefAuto{PeanoMulLeftDistributive}{1}}
    \proofstepwidestar[1]{%
      a(c+e)+b(d+f)=(ac+ae)+(bd+bf)}{\rIE{2,3}}
    \proofstepwidestar[1]{%
      (ac+ae)+(bd+bf)=(ac+bd)+(ae+bf)}{%
      \FormulaRefAuto{PeanoAddInterchange}{1}}
    \proofstepwidestar[1]{%
      a(c+e)+b(d+f)=(ac+bd)+(ae+bf)}{%
      \rIE{4,5}}
  \closeproofpartwideindR

  \proofpartwideindR[Negative Koordinate distributiv ausmultiplizieren]{%
    \begin{aligned}[t]
      &a,b,c,d,e,f\in\mathbb N\vdash{}\\
      &a(d+f)+b(c+e)=(ad+bc)+(af+be)
    \end{aligned}
  }{%
    a,b,c,d,e,f\in\mathbb N\vdash
    a(d+f)+b(c+e)=(ad+bc)+(af+be)
  }[IntMulLeftDistributiveNegativeCoordinate]
    \proofstepwidestar[1]{a,b,c,d,e,f\in\mathbb N}{\rA}
    \proofstepwidestar[1]{a(d+f)=ad+af}{%
      \FormulaRefAuto{PeanoMulLeftDistributive}{1}}
    \proofstepwidestar[1]{b(c+e)=bc+be}{%
      \FormulaRefAuto{PeanoMulLeftDistributive}{1}}
    \proofstepwidestar[1]{%
      a(d+f)+b(c+e)=(ad+af)+(bc+be)}{\rIE{2,3}}
    \proofstepwidestar[1]{%
      (ad+af)+(bc+be)=(ad+bc)+(af+be)}{%
      \FormulaRefAuto{PeanoAddInterchange}{1}}
    \proofstepwidestar[1]{%
      a(d+f)+b(c+e)=(ad+bc)+(af+be)}{%
      \rIE{4,5}}
  \closeproofpartwideindR

  \proofpartwideindR[Distributivrechnung auf Klassen]{%
    \begin{aligned}[t]
      &a,b,c,d,e,f\in\mathbb N\vdash{}\\
      &\IntClass{a}{b}\cdot
       (\IntClass{c}{d}+\IntClass{e}{f})
       =
       \IntClass{a}{b}\cdot\IntClass{c}{d}
       +\IntClass{a}{b}\cdot\IntClass{e}{f}
    \end{aligned}
  }{%
    a,b,c,d,e,f\in\mathbb N\vdash
    \IntClass{a}{b}\cdot
    (\IntClass{c}{d}+\IntClass{e}{f})
    =
    \IntClass{a}{b}\cdot\IntClass{c}{d}
    +\IntClass{a}{b}\cdot\IntClass{e}{f}
  }[IntMulLeftDistributiveOnClasses]
    \proofstepwidestar[1]{a,b,c,d,e,f\in\mathbb N}{\rA}
    \proofstepwidestar[1]{%
      \IntClass{a}{b}\cdot
      (\IntClass{c}{d}+\IntClass{e}{f})
      =
      \IntClass{a(c+e)+b(d+f)}
               {a(d+f)+b(c+e)}}{%
      \FormulaRefAuto{IntegerAdditionDef}{%
        \FormulaRefAuto{IntegerMultiplicationDef}{1}}}
    \proofstepwidestar[1]{%
      \IntClass{a}{b}\cdot\IntClass{c}{d}
      +\IntClass{a}{b}\cdot\IntClass{e}{f}
      =
      \IntClass{(ac+bd)+(ae+bf)}
               {(ad+bc)+(af+be)}}{%
      \FormulaRefAuto{IntegerMultiplicationDef}{%
        \FormulaRefAuto{IntegerAdditionDef}{1}}}
    \proofstepwidestar[1]{%
      a(c+e)+b(d+f)=(ac+bd)+(ae+bf)}{%
      \FormulaRefAuto{IntMulLeftDistributivePositiveCoordinate}{1}}
    \proofstepwidestar[1]{%
      a(d+f)+b(c+e)=(ad+bc)+(af+be)}{%
      \FormulaRefAuto{IntMulLeftDistributiveNegativeCoordinate}{1}}
    \proofstepwidestar[1]{%
      \IntClass{a}{b}\cdot
      (\IntClass{c}{d}+\IntClass{e}{f})
      =
      \IntClass{a}{b}\cdot\IntClass{c}{d}
      +\IntClass{a}{b}\cdot\IntClass{e}{f}}{\rIE{2,3,4,5}}
  \closeproofpartwideindR

  \proofpartwide{Schluss}
    \proofstepwidestar[1]{z,w,u\in\mathbb Z}{\rA}
    \proofstepwidestar[1]{z\in\mathbb Z}{\rAEa{1}}
    \proofstepwidestar[1]{w\in\mathbb Z}{\rAEa{\rAEb{1}}}
    \proofstepwidestar[1]{u\in\mathbb Z}{\rAEb{\rAEb{1}}}
    \proofstepwidestar[1]{\exists a,b\in\mathbb N\,
      z=\IntClass{a}{b}}{\FormulaRefAuto{IntegerPairRepresentative}{2}}
    \proofstepwidestar[1]{\exists c,d\in\mathbb N\,
      w=\IntClass{c}{d}}{\FormulaRefAuto{IntegerPairRepresentative}{3}}
    \proofstepwidestar[1]{\exists e,f\in\mathbb N\,
      u=\IntClass{e}{f}}{\FormulaRefAuto{IntegerPairRepresentative}{4}}
    \proofstepwidestar[8]{a,b,c,d,e,f\in\mathbb N\land
      z=\IntClass{a}{b}\land w=\IntClass{c}{d}\land
      u=\IntClass{e}{f}}{\rA}
    \proofstepwidestar[8]{%
      \IntClass{a}{b}\cdot
      (\IntClass{c}{d}+\IntClass{e}{f})
      =
      \IntClass{a}{b}\cdot\IntClass{c}{d}
      +\IntClass{a}{b}\cdot\IntClass{e}{f}}{%
      \FormulaRefAuto{IntMulLeftDistributiveOnClasses}{8}}
    \proofstepwidestar[8]{%
      z\cdot(w+u)=z\cdot w+z\cdot u}{\rIE{8,9}}
    \proofstepwidestar[1]{z\cdot(w+u)=z\cdot w+z\cdot u}{%
      \rEE{5,6,7,8,10}}
  \closeproofpartwide
\end{tabproofsplitwide}

\FormulaThmDeltaK[Rechtsdistributivität ganzer Zahlen]{%
  z,w,u\in\mathbb Z\vdash
  (z+w)\cdot u=z\cdot u+w\cdot u%
}{IntMulRightDistributive}{%
  \DeltaRow{Mengen}{z\dsep w\dsep u}%
}
\begin{tabproof}
  \proofstep{1}{z,w,u\in\mathbb Z}{\rA}
  \proofstep{1}{(z+w)\cdot u=u\cdot(z+w)}{%
    \FormulaRefAuto{IntMulCommutative}{%
      \FormulaRefAuto{IntAddClosure}{1},1}}
  \proofstep{1}{u\cdot(z+w)=u\cdot z+u\cdot w}{%
    \FormulaRefAuto{IntMulLeftDistributive}{1}}
  \proofstep{1}{u\cdot z=z\cdot u}{\FormulaRefAuto{IntMulCommutative}{1}}
  \proofstep{1}{u\cdot w=w\cdot u}{\FormulaRefAuto{IntMulCommutative}{1}}
  \proofstep{1}{u\cdot z+u\cdot w=z\cdot u+w\cdot u}{\rIE{4,5}}
  \proofstep{1}{(z+w)\cdot u=z\cdot u+w\cdot u}{%
    \rIE{2,3,6}}
\end{tabproof}

\section{Verträglichkeit der Einbettung mit der Arithmetik}

\FormulaThmDeltaK[Additivität der natürlichen Einbettung]{%
  m,n\in\mathbb N\vdash
  \IntEmbed(m+n)=\IntEmbed(m)+\IntEmbed(n)%
}{NaturalIntegerEmbeddingAdditive}{%
  \DeltaRow{Mengen}{m\dsep n}%
  \DeltaRow{Funktion}{\IntEmbed}%
}
\begin{tabproof}
  \proofstep{1}{m,n\in\mathbb N}{\rA}
  \proofstep{1}{m+n\in\mathbb N}{\FormulaRefAuto{PeanoAddClosure}{1}}
  \proofstep{1}{\IntEmbed(m+n)=\IntClass{m+n}{0}}{%
    \FormulaRefAuto{NaturalIntegerEmbedding}{2}}
  \proofstep{1}{\IntEmbed(m)=\IntClass{m}{0}}{%
    \FormulaRefAuto{NaturalIntegerEmbedding}{1}}
  \proofstep{1}{\IntEmbed(n)=\IntClass{n}{0}}{%
    \FormulaRefAuto{NaturalIntegerEmbedding}{1}}
  \proofstep{1}{\IntEmbed(m)+\IntEmbed(n)=
    \IntClass{m+n}{0+0}}{%
    \FormulaRefAuto{IntegerAdditionDef}{1,4,5}}
  \proofstep{}{0+0=0}{%
    \FormulaRefAuto{PeanoAddRightZero}{\FormulaRefAuto{PeanoZero}}}
  \proofstep{1}{\IntEmbed(m)+\IntEmbed(n)=
    \IntClass{m+n}{0}}{\rIE{7,6}}
  \proofstep{1}{\IntEmbed(m+n)=
    \IntEmbed(m)+\IntEmbed(n)}{%
    \FormulaRefAuto{a=b\dsep c=b\vdash a=c}{3,8}}
\end{tabproof}

\FormulaThmDeltaK[Multiplikativität der natürlichen Einbettung]{%
  m,n\in\mathbb N\vdash
  \IntEmbed(m\cdot n)=\IntEmbed(m)\cdot\IntEmbed(n)%
}{NaturalIntegerEmbeddingMultiplicative}{%
  \DeltaRow{Mengen}{m\dsep n}%
  \DeltaRow{Funktion}{\IntEmbed}%
}
\begin{tabproofsplitwide}
  \proofpartwideindR[Produkt der eingebetteten Faktoren]{%
    m,n\in\mathbb N\vdash
    \IntEmbed(m)\cdot\IntEmbed(n)=\IntClass{mn}{0}%
  }{%
    m,n\in\mathbb N\vdash
    \IntEmbed(m)\cdot\IntEmbed(n)=\IntClass{mn}{0}%
  }[NaturalIntegerEmbeddingProductClass]
    \proofstepwidestar[1]{m,n\in\mathbb N}{\rA}
    \proofstepwidestar[1]{mn\in\mathbb N}{%
      \FormulaRefAuto{PeanoMulClosure}{1}}
    \proofstepwidestar[1]{\IntEmbed(m)=\IntClass{m}{0}}{%
      \FormulaRefAuto{NaturalIntegerEmbedding}{1}}
    \proofstepwidestar[1]{\IntEmbed(n)=\IntClass{n}{0}}{%
      \FormulaRefAuto{NaturalIntegerEmbedding}{1}}
    \proofstepwidestar[1]{\IntEmbed(m)\cdot\IntEmbed(n)=
      \IntClass{mn+0\cdot0}{m\cdot0+0\cdot n}}{%
      \FormulaRefAuto{IntegerMultiplicationDef}{1,3,4}}
    \proofstepwidestar[]{0\cdot0=0}{%
      \FormulaRefAuto{PeanoMulLeftZero}{\FormulaRefAuto{PeanoZero}}}
    \proofstepwidestar[1]{m\cdot0=0}{%
      \FormulaRefAuto{PeanoMulRightZero}{1}}
    \proofstepwidestar[1]{0\cdot n=0}{%
      \FormulaRefAuto{PeanoMulLeftZero}{1}}
    \proofstepwidestar[1]{mn+0\cdot0=mn}{%
      \rIE{6,\FormulaRefAuto{PeanoAddRightZero}{2}}}
    \proofstepwidestar[1]{m\cdot0+0\cdot n=0}{%
      \rIE{7,8,
        \FormulaRefAuto{PeanoAddRightZero}{\FormulaRefAuto{PeanoZero}}}}
    \proofstepwidestar[1]{%
      \IntEmbed(m)\cdot\IntEmbed(n)=\IntClass{mn}{0}}{%
      \rIE{5,9,10}}
  \closeproofpartwideindR

  \proofpartwideindR[Vergleich mit dem Produktbild]{%
    m,n\in\mathbb N\vdash
    \IntEmbed(m\cdot n)=\IntEmbed(m)\cdot\IntEmbed(n)%
  }{%
    m,n\in\mathbb N\vdash
    \IntEmbed(m\cdot n)=\IntEmbed(m)\cdot\IntEmbed(n)%
  }[NaturalIntegerEmbeddingMultiplicativeConclusion]
    \proofstepwidestar[1]{m,n\in\mathbb N}{\rA}
    \proofstepwidestar[1]{mn\in\mathbb N}{%
      \FormulaRefAuto{PeanoMulClosure}{1}}
    \proofstepwidestar[1]{\IntEmbed(mn)=\IntClass{mn}{0}}{%
      \FormulaRefAuto{NaturalIntegerEmbedding}{2}}
    \proofstepwidestar[1]{%
      \IntEmbed(m)\cdot\IntEmbed(n)=\IntClass{mn}{0}}{%
      \FormulaRefAuto{NaturalIntegerEmbeddingProductClass}{1}}
    \proofstepwidestar[1]{\IntEmbed(mn)=
      \IntEmbed(m)\cdot\IntEmbed(n)}{%
      \FormulaRefAuto{a=b\dsep c=b\vdash a=c}{3,4}}
  \closeproofpartwideindR
\end{tabproofsplitwide}

\chapter{Nachfolger, Vorgänger und Normalformen}

\section{Nachfolger und Vorgänger}

\FormulaDefDeltaK[Nachfolger einer ganzen Zahl]{%
  z\in\mathbb Z\vdash
  \IntSucc(z)\coloneqq z+1_{\mathbb Z}%
}{IntegerSuccessorDef}{%
  \DeltaRow{Mengen}{z}%
  \DeltaRow{Funktionensymbol}{\IntSucc}%
  \DeltaRow{Neue Symbole}{\IntSucc}%
}

\FormulaDefDeltaK[Vorgänger einer ganzen Zahl]{%
  z\in\mathbb Z\vdash
  \IntPred(z)\coloneqq z+(-1_{\mathbb Z})%
}{IntegerPredecessorDef}{%
  \DeltaRow{Mengen}{z}%
  \DeltaRow{Funktionensymbol}{\IntPred}%
  \DeltaRow{Neue Symbole}{\IntPred}%
}

\FormulaThmDeltaK[Nachfolgerfunktion auf den ganzen Zahlen]{%
  \IntSucc\colon\mathbb Z\to\mathbb Z%
}{IntegerSuccessorFunction}{%
  \DeltaRow{Funktion}{\IntSucc}%
}
\begin{tabproof}
  \proofstep{1}{z\in\mathbb Z}{\rA}
  \proofstep{}{1_{\mathbb Z}\in\mathbb Z}{\FormulaRefAuto{IntOneCarrier}}
  \proofstep{1}{z+1_{\mathbb Z}\in\mathbb Z}{%
    \FormulaRefAuto{IntAddClosure}{1,2}}
  \proofstep{1}{\IntSucc(z)=z+1_{\mathbb Z}}{%
    \FormulaRefAuto{IntegerSuccessorDef}{1}}
  \proofstep{1}{\IntSucc(z)\in\mathbb Z}{\rIE{4,3}}
  \proofstep{}{\IntSucc\colon\mathbb Z\to\mathbb Z}{%
    \FormulaRefAuto{x\in A\vdash t(x)\coloneq y\dsep
      x\in A\vdash t(x)\in B
      \exists! F\colon A\to B\,\forall x\in A\,F(x)=t(x)}{1,5}}
\end{tabproof}

\FormulaThmDeltaK[Vorgängerfunktion auf den ganzen Zahlen]{%
  \IntPred\colon\mathbb Z\to\mathbb Z%
}{IntegerPredecessorFunction}{%
  \DeltaRow{Funktion}{\IntPred}%
}
\begin{tabproof}
  \proofstep{1}{z\in\mathbb Z}{\rA}
  \proofstep{}{1_{\mathbb Z}\in\mathbb Z}{\FormulaRefAuto{IntOneCarrier}}
  \proofstep{}{-1_{\mathbb Z}\in\mathbb Z}{%
    \FormulaRefAuto{IntNegClosure}{2}}
  \proofstep{1}{z+(-1_{\mathbb Z})\in\mathbb Z}{%
    \FormulaRefAuto{IntAddClosure}{1,3}}
  \proofstep{1}{\IntPred(z)=z+(-1_{\mathbb Z})}{%
    \FormulaRefAuto{IntegerPredecessorDef}{1}}
  \proofstep{1}{\IntPred(z)\in\mathbb Z}{\rIE{5,4}}
  \proofstep{}{\IntPred\colon\mathbb Z\to\mathbb Z}{%
    \FormulaRefAuto{x\in A\vdash t(x)\coloneq y\dsep
      x\in A\vdash t(x)\in B
      \exists! F\colon A\to B\,\forall x\in A\,F(x)=t(x)}{1,6}}
\end{tabproof}

\FormulaThmDeltaK[Nachfolger auf Klassen]{%
  a,b\in\mathbb N\vdash
  \IntSucc(\IntClass{a}{b})=\IntClass{\Succ(a)}{b}%
}{IntegerSuccessorClassValue}{%
  \DeltaRow{Mengen}{a\dsep b}%
}
\begin{tabproof}
  \proofstep{1}{a,b\in\mathbb N}{\rA}
  \proofstep{1}{\IntClass{a}{b}\in\mathbb Z}{%
    \FormulaRefAuto{IntegerClassInZ}{1}}
  \proofstep{1}{\IntSucc(\IntClass{a}{b})=
    \IntClass{a}{b}+1_{\mathbb Z}}{%
    \FormulaRefAuto{IntegerSuccessorDef}{2}}
  \proofstep{}{1_{\mathbb Z}=\IntClass{1}{0}}{%
    \FormulaRefAuto{IntegerOneDef}{%
      \FormulaRefAuto{NaturalIntegerEmbedding}{%
        \FormulaRefAuto{PeanoOneInN}}}}
  \proofstep{1}{\IntClass{a}{b}+1_{\mathbb Z}
    =\IntClass{a+1}{b+0}}{%
    \FormulaRefAuto{IntegerAdditionDef}{1,4}}
  \proofstep{1}{a+1=\Succ(a)}{\FormulaRefAuto{PeanoAddOneSucc}{1}}
  \proofstep{1}{b+0=b}{\FormulaRefAuto{PeanoAddRightZero}{1}}
  \proofstep{1}{\IntSucc(\IntClass{a}{b})=
    \IntClass{\Succ(a)}{b}}{\rIE{3,5,6,7}}
\end{tabproof}

\FormulaThmDeltaK[Vorgänger auf Klassen]{%
  a,b\in\mathbb N\vdash
  \IntPred(\IntClass{a}{b})=\IntClass{a}{\Succ(b)}%
}{IntegerPredecessorClassValue}{%
  \DeltaRow{Mengen}{a\dsep b}%
}
\begin{tabproof}
  \proofstep{1}{a,b\in\mathbb N}{\rA}
  \proofstep{1}{\IntClass{a}{b}\in\mathbb Z}{%
    \FormulaRefAuto{IntegerClassInZ}{1}}
  \proofstep{1}{\IntPred(\IntClass{a}{b})=
    \IntClass{a}{b}+(-1_{\mathbb Z})}{%
    \FormulaRefAuto{IntegerPredecessorDef}{2}}
  \proofstep{}{-1_{\mathbb Z}=\IntClass{0}{1}}{%
    \FormulaRefAuto{IntegerNegationDef}{%
      \FormulaRefAuto{IntegerOneDef},%
      \FormulaRefAuto{NaturalIntegerEmbedding}}}
  \proofstep{1}{\IntClass{a}{b}+(-1_{\mathbb Z})
    =\IntClass{a+0}{b+1}}{%
    \FormulaRefAuto{IntegerAdditionDef}{1,4}}
  \proofstep{1}{a+0=a}{\FormulaRefAuto{PeanoAddRightZero}{1}}
  \proofstep{1}{b+1=\Succ(b)}{\FormulaRefAuto{PeanoAddOneSucc}{1}}
  \proofstep{1}{\IntPred(\IntClass{a}{b})=
    \IntClass{a}{\Succ(b)}}{\rIE{3,5,6,7}}
\end{tabproof}

\FormulaThmDeltaK[Vorgänger nach Nachfolger]{%
  z\in\mathbb Z\vdash
  \IntPred(\IntSucc(z))=z%
}{IntegerPredAfterSucc}{%
  \DeltaRow{Mengen}{z}%
}
\begin{tabproofsplitwide}
  \proofpartwideindR[Auflösen der Definitionen]{%
    z\in\mathbb Z\vdash
    \IntPred(\IntSucc(z))=(z+1_{\mathbb Z})+(-1_{\mathbb Z})%
  }{%
    z\in\mathbb Z\vdash
    \IntPred(\IntSucc(z))=(z+1_{\mathbb Z})+(-1_{\mathbb Z})%
  }[IntegerPredAfterSuccExpansion]
    \proofstepwidestar[1]{z\in\mathbb Z}{\rA}
    \proofstepwidestar[1]{\IntSucc(z)=z+1_{\mathbb Z}}{%
      \FormulaRefAuto{IntegerSuccessorDef}{1}}
    \proofstepwidestar[1]{\IntSucc(z)\in\mathbb Z}{%
      \FormulaRefAuto{F\colon A\to B\dsep x\in A\vdash F(x)\in B}{%
        \FormulaRefAuto{IntegerSuccessorFunction},1}}
    \proofstepwidestar[1]{\IntPred(\IntSucc(z))=
      \IntSucc(z)+(-1_{\mathbb Z})}{%
      \FormulaRefAuto{IntegerPredecessorDef}{3}}
    \proofstepwidestar[1]{\IntPred(\IntSucc(z))=
      (z+1_{\mathbb Z})+(-1_{\mathbb Z})}{\rIE{2,4}}
  \closeproofpartwideindR

  \proofpartwide{Schluss}
    \proofstepwidestar[1]{z\in\mathbb Z}{\rA}
    \proofstepwidestar[1]{\IntPred(\IntSucc(z))=
      (z+1_{\mathbb Z})+(-1_{\mathbb Z})}{%
      \FormulaRefAuto{IntegerPredAfterSuccExpansion}{1}}
    \proofstepwidestar[]{1_{\mathbb Z}\in\mathbb Z}{%
      \FormulaRefAuto{IntOneCarrier}}
    \proofstepwidestar[]{1_{\mathbb Z}+(-1_{\mathbb Z})=0_{\mathbb Z}}{%
      \rAEa{\FormulaRefAuto{IntAddInverse}{3}}}
    \proofstepwidestar[1]{(z+1_{\mathbb Z})+(-1_{\mathbb Z})
      =z+(1_{\mathbb Z}+(-1_{\mathbb Z}))}{%
      \FormulaRefAuto{IntAddAssociative}{1,3,%
        \FormulaRefAuto{IntNegClosure}{3}}}
    \proofstepwidestar[1]{z+(1_{\mathbb Z}+(-1_{\mathbb Z}))
      =z+0_{\mathbb Z}}{\rIE{4}}
    \proofstepwidestar[1]{z+0_{\mathbb Z}=z}{%
      \rAEa{\FormulaRefAuto{IntAddZero}{1}}}
    \proofstepwidestar[1]{\IntPred(\IntSucc(z))=z}{%
      \rIE{2,5,6,7}}
  \closeproofpartwide
\end{tabproofsplitwide}

\FormulaThmDeltaK[Nachfolger nach Vorgänger]{%
  z\in\mathbb Z\vdash
  \IntSucc(\IntPred(z))=z%
}{IntegerSuccAfterPred}{%
  \DeltaRow{Mengen}{z}%
}
\begin{tabproofsplitwide}
  \proofpartwideindR[Auflösen der Definitionen]{%
    z\in\mathbb Z\vdash
    \IntSucc(\IntPred(z))=(z+(-1_{\mathbb Z}))+1_{\mathbb Z}%
  }{%
    z\in\mathbb Z\vdash
    \IntSucc(\IntPred(z))=(z+(-1_{\mathbb Z}))+1_{\mathbb Z}%
  }[IntegerSuccAfterPredExpansion]
    \proofstepwidestar[1]{z\in\mathbb Z}{\rA}
    \proofstepwidestar[1]{\IntPred(z)=z+(-1_{\mathbb Z})}{%
      \FormulaRefAuto{IntegerPredecessorDef}{1}}
    \proofstepwidestar[1]{\IntPred(z)\in\mathbb Z}{%
      \FormulaRefAuto{F\colon A\to B\dsep x\in A\vdash F(x)\in B}{%
        \FormulaRefAuto{IntegerPredecessorFunction},1}}
    \proofstepwidestar[1]{%
      \IntSucc(\IntPred(z))=\IntPred(z)+1_{\mathbb Z}}{%
      \FormulaRefAuto{IntegerSuccessorDef}{3}}
    \proofstepwidestar[1]{%
      \IntSucc(\IntPred(z))=(z+(-1_{\mathbb Z}))+1_{\mathbb Z}}{%
      \rIE{2,4}}
  \closeproofpartwideindR

  \proofpartwideindR[Assoziativität und Inverses]{%
    z\in\mathbb Z\vdash\IntSucc(\IntPred(z))=z%
  }{%
    z\in\mathbb Z\vdash\IntSucc(\IntPred(z))=z%
  }[IntegerSuccAfterPredConclusion]
    \proofstepwidestar[1]{z\in\mathbb Z}{\rA}
    \proofstepwidestar[1]{%
      \IntSucc(\IntPred(z))=(z+(-1_{\mathbb Z}))+1_{\mathbb Z}}{%
      \FormulaRefAuto{IntegerSuccAfterPredExpansion}{1}}
    \proofstepwidestar[]{1_{\mathbb Z}\in\mathbb Z}{%
      \FormulaRefAuto{IntOneCarrier}}
    \proofstepwidestar[]{-1_{\mathbb Z}\in\mathbb Z}{%
      \FormulaRefAuto{IntNegClosure}{3}}
    \proofstepwidestar[]{%
      (-1_{\mathbb Z})+1_{\mathbb Z}=0_{\mathbb Z}}{%
      \rAEb{\FormulaRefAuto{IntAddInverse}{3}}}
    \proofstepwidestar[1]{%
      (z+(-1_{\mathbb Z}))+1_{\mathbb Z}
      =z+((-1_{\mathbb Z})+1_{\mathbb Z})}{%
      \FormulaRefAuto{IntAddAssociative}{1,4,3}}
    \proofstepwidestar[1]{%
      z+((-1_{\mathbb Z})+1_{\mathbb Z})=z+0_{\mathbb Z}}{\rIE{5}}
    \proofstepwidestar[1]{z+0_{\mathbb Z}=z}{%
      \rAEa{\FormulaRefAuto{IntAddZero}{1}}}
    \proofstepwidestar[1]{\IntSucc(\IntPred(z))=z}{%
      \rIE{2,6,7,8}}
  \closeproofpartwideindR
\end{tabproofsplitwide}

\FormulaThmDeltaK[Bijektivität des Ganzzahlnachfolgers]{%
  \IntSucc\colon\mathbb Z\bij\mathbb Z%
}{IntegerSuccessorBijective}{%
  \DeltaRow{Funktionen}{\IntSucc\dsep\IntPred}%
}
\begin{tabproofsplitwide}
  \proofpartwideindR[Injektivität]{%
    x,y\in\mathbb Z\dsep
    \IntSucc(x)=\IntSucc(y)\vdash x=y%
  }{%
    x,y\in\mathbb Z\dsep
    \IntSucc(x)=\IntSucc(y)\vdash x=y%
  }[IntegerSuccessorInjective]
    \proofstepwidestar[1]{x,y\in\mathbb Z}{\rA}
    \proofstepwidestar[2]{\IntSucc(x)=\IntSucc(y)}{\rA}
    \proofstepwidestar[1]{\IntPred(\IntSucc(x))=x}{%
      \FormulaRefAuto{IntegerPredAfterSucc}{1}}
    \proofstepwidestar[1]{\IntPred(\IntSucc(y))=y}{%
      \FormulaRefAuto{IntegerPredAfterSucc}{1}}
    \proofstepwidestar[2]{\IntPred(\IntSucc(x))=
      \IntPred(\IntSucc(y))}{\rIE{2}}
    \proofstepwidestar[1,2]{x=y}{\rIE{3,4,5}}
  \closeproofpartwideindR

  \proofpartwideindR[Surjektivität]{%
    y\in\mathbb Z\vdash
    \exists x\in\mathbb Z\,\IntSucc(x)=y%
  }{%
    y\in\mathbb Z\vdash
    \exists x\in\mathbb Z\,\IntSucc(x)=y%
  }[IntegerSuccessorSurjective]
    \proofstepwidestar[1]{y\in\mathbb Z}{\rA}
    \proofstepwidestar[1]{\IntPred(y)\in\mathbb Z}{%
      \FormulaRefAuto{F\colon A\to B\dsep x\in A\vdash F(x)\in B}{%
        \FormulaRefAuto{IntegerPredecessorFunction},1}}
    \proofstepwidestar[1]{\IntSucc(\IntPred(y))=y}{%
      \FormulaRefAuto{IntegerSuccAfterPred}{1}}
    \proofstepwidestar[1]{\exists x\in\mathbb Z\,\IntSucc(x)=y}{%
      \rEI{\rAI{2,3}}}
  \closeproofpartwideindR

  \proofpartwide{Schluss}
    \proofstepwidestar[]{\IntSucc\colon\mathbb Z\to\mathbb Z}{%
      \FormulaRefAuto{IntegerSuccessorFunction}}
    \proofstepwidestar[]{%
      \forall x,y\in\mathbb Z\,
      (\IntSucc(x)=\IntSucc(y)\rightarrow x=y)}{%
      \FormulaRefAuto{IntegerSuccessorInjective}}
    \proofstepwidestar[]{\IntSucc\colon\mathbb Z\inj\mathbb Z}{%
      \FormulaRefAuto{Injektive Funktion}{1,2}}
    \proofstepwidestar[]{%
      \forall y\in\mathbb Z\,\exists x\in\mathbb Z\,
      \IntSucc(x)=y}{\FormulaRefAuto{IntegerSuccessorSurjective}}
    \proofstepwidestar[]{\IntSucc\colon\mathbb Z\sur\mathbb Z}{%
      \FormulaRefAuto{Surjektive Funktion}{1,4}}
    \proofstepwidestar[]{\IntSucc\colon\mathbb Z\bij\mathbb Z}{%
      \FormulaRefAuto{F\colon A\inj B\dsep F\colon A\sur B
        \vdash F\colon A\bij B}{3,5}}
  \closeproofpartwide
\end{tabproofsplitwide}

\section{Normalformen}

\FormulaThmDeltaK[Positiver Strahl und Nachfolger]{%
  n\in\mathbb N\vdash
  \IntSucc(\IntEmbed(n))=\IntEmbed(\Succ(n))%
}{IntegerSuccessorOnEmbedding}{%
  \DeltaRow{Mengen}{n}%
}
\begin{tabproof}
  \proofstep{1}{n\in\mathbb N}{\rA}
  \proofstep{1}{\IntSucc(\IntEmbed(n))=
    \IntEmbed(n)+1_{\mathbb Z}}{%
    \FormulaRefAuto{IntegerSuccessorDef}{%
      \FormulaRefAuto{F\colon A\to B\dsep x\in A\vdash F(x)\in B}{%
        \FormulaRefAuto{NaturalIntegerEmbeddingFunction},1}}}
  \proofstep{}{1_{\mathbb Z}=\IntEmbed(1)}{\FormulaRefAuto{IntegerOneDef}}
  \proofstep{1}{\IntEmbed(n)+\IntEmbed(1)=\IntEmbed(n+1)}{%
    \FormulaRefAuto{a=b\vdash b=a}{%
      \FormulaRefAuto{NaturalIntegerEmbeddingAdditive}{1,%
        \FormulaRefAuto{PeanoOneInN}}}}
  \proofstep{1}{n+1=\Succ(n)}{\FormulaRefAuto{PeanoAddOneSucc}{1}}
  \proofstep{1}{\IntSucc(\IntEmbed(n))=\IntEmbed(\Succ(n))}{\rIE{2,3,4,5}}
\end{tabproof}

\FormulaThmDeltaK[Negativer Strahl und Vorgänger]{%
  n\in\mathbb N\vdash
  \IntPred(-\IntEmbed(n))=-\IntEmbed(\Succ(n))%
}{IntegerPredecessorOnNegativeEmbedding}{%
  \DeltaRow{Mengen}{n}%
}
\begin{tabproof}
  \proofstep{1}{n\in\mathbb N}{\rA}
  \proofstep{1}{\IntEmbed(n)=\IntClass{n}{0}}{%
    \FormulaRefAuto{NaturalIntegerEmbedding}{1}}
  \proofstep{1}{-\IntEmbed(n)=\IntClass{0}{n}}{%
    \FormulaRefAuto{IntegerNegationDef}{1,2}}
  \proofstep{1}{\IntPred(-\IntEmbed(n))=\IntClass{0}{\Succ(n)}}{%
    \FormulaRefAuto{IntegerPredecessorClassValue}{1,3}}
  \proofstep{1}{\IntEmbed(\Succ(n))=\IntClass{\Succ(n)}{0}}{%
    \FormulaRefAuto{NaturalIntegerEmbedding}{%
      \FormulaRefAuto{PeanoSuccClosure}{1}}}
  \proofstep{1}{-\IntEmbed(\Succ(n))=\IntClass{0}{\Succ(n)}}{%
    \FormulaRefAuto{IntegerNegationDef}{1,5}}
  \proofstep{1}{\IntPred(-\IntEmbed(n))=-\IntEmbed(\Succ(n))}{%
    \FormulaRefAuto{a=b\dsep c=b\vdash a=c}{4,6}}
\end{tabproof}

\FormulaThmDeltaK[Normalform ganzer Zahlen]{%
  z\in\mathbb Z\vdash
  \exists n\in\mathbb N\,
  \bigl(z=\IntEmbed(n)\lor z=-\IntEmbed(n)\bigr)%
}{IntegerSignedNormalForm}{%
  \DeltaRow{Mengen}{z\dsep a\dsep b\dsep n}%
}
\begin{tabproofsplitwide}
  \proofpartwideindR[Fall \(b\leq a\)]{%
    \begin{aligned}[t]
      &a,b\in\mathbb N\dsep b\leq a\vdash{}\\
      &\exists n\in\mathbb N\,
      \IntClass{a}{b}=\IntEmbed(n)
    \end{aligned}
  }{%
    a,b\in\mathbb N\dsep b\leq a
    \vdash\exists n\in\mathbb N\,
    \IntClass{a}{b}=\IntEmbed(n)
  }[IntegerNormalFormNonnegative]
    \proofstepwidestar[1]{a,b\in\mathbb N}{\rA}
    \proofstepwidestar[2]{b\leq a}{\rA}
    \proofstepwidestar[1,2]{\exists n\in\mathbb N\,(b+n=a)}{%
      \FormulaRefAuto{m\in\mathbb{N}\dsep n\in\mathbb{N}\vdash
        m\leq n\leftrightarrow
        \exists k\in\mathbb{N}\,(m+k=n)}{1,2}}
    \proofstepwidestar[4]{n\in\mathbb N\land b+n=a}{\rA}
    \proofstepwidestar[1,4]{a+0=b+n}{%
      \rIE{\FormulaRefAuto{PeanoAddRightZero}{1},\rAEb{4}}}
    \proofstepwidestar[1,4]{\IntClass{a}{b}=\IntClass{n}{0}}{%
      \FormulaRefAuto{IntegerClassEquality}{1,4,5}}
    \proofstepwidestar[4]{\IntEmbed(n)=\IntClass{n}{0}}{%
      \FormulaRefAuto{NaturalIntegerEmbedding}{4}}
    \proofstepwidestar[1,4]{\IntClass{a}{b}=\IntEmbed(n)}{%
      \FormulaRefAuto{a=b\dsep c=b\vdash a=c}{6,7}}
    \proofstepwidestar[1,2]{\exists n\in\mathbb N\,
      \IntClass{a}{b}=\IntEmbed(n)}{\rEE{3,4,8}}
  \closeproofpartwideindR

  \proofpartwideindR[Fall \(a\leq b\)]{%
    \begin{aligned}[t]
      &a,b\in\mathbb N\dsep a\leq b\vdash{}\\
      &\exists n\in\mathbb N\,
      \IntClass{a}{b}=-\IntEmbed(n)
    \end{aligned}
  }{%
    a,b\in\mathbb N\dsep a\leq b
    \vdash\exists n\in\mathbb N\,
    \IntClass{a}{b}=-\IntEmbed(n)
  }[IntegerNormalFormNonpositive]
    \proofstepwidestar[1]{a,b\in\mathbb N}{\rA}
    \proofstepwidestar[2]{a\leq b}{\rA}
    \proofstepwidestar[1,2]{\exists n\in\mathbb N\,(a+n=b)}{%
      \FormulaRefAuto{m\in\mathbb{N}\dsep n\in\mathbb{N}\vdash
        m\leq n\leftrightarrow
        \exists k\in\mathbb{N}\,(m+k=n)}{1,2}}
    \proofstepwidestar[4]{n\in\mathbb N\land a+n=b}{\rA}
    \proofstepwidestar[1,4]{a+n=b+0}{%
      \rIE{\rAEb{4},\FormulaRefAuto{PeanoAddRightZero}{1}}}
    \proofstepwidestar[1,4]{\IntClass{a}{b}=\IntClass{0}{n}}{%
      \FormulaRefAuto{IntegerClassEquality}{1,4,5}}
    \proofstepwidestar[4]{\IntEmbed(n)=\IntClass{n}{0}}{%
      \FormulaRefAuto{NaturalIntegerEmbedding}{4}}
    \proofstepwidestar[4]{-\IntEmbed(n)=\IntClass{0}{n}}{%
      \FormulaRefAuto{IntegerNegationDef}{4,7}}
    \proofstepwidestar[1,4]{\IntClass{a}{b}=-\IntEmbed(n)}{%
      \FormulaRefAuto{a=b\dsep c=b\vdash a=c}{6,8}}
    \proofstepwidestar[1,2]{\exists n\in\mathbb N\,
      \IntClass{a}{b}=-\IntEmbed(n)}{\rEE{3,4,9}}
  \closeproofpartwideindR

  \proofpartwide{Schluss}
    \proofstepwidestar[1]{z\in\mathbb Z}{\rA}
    \proofstepwidestar[1]{\exists a,b\in\mathbb N\,
      z=\IntClass{a}{b}}{\FormulaRefAuto{IntegerPairRepresentative}{1}}
    \proofstepwidestar[3]{a,b\in\mathbb N\land z=\IntClass{a}{b}}{\rA}
    \proofstepwidestar[3]{b\leq a\lor a\leq b}{%
      \FormulaRefAuto{PeanoOrderComparison}{3}}
    \proofstepwidestar[5]{b\leq a}{\rA}
    \proofstepwidestar[3,5]{\exists n\in\mathbb N\,
      \IntClass{a}{b}=\IntEmbed(n)}{%
      \FormulaRefAuto{IntegerNormalFormNonnegative}{3,5}}
    \proofstepwidestar[3,5]{\exists n\in\mathbb N\,
      (z=\IntEmbed(n)\lor z=-\IntEmbed(n))}{\rIE{3,6}}
    \proofstepwidestar[8]{a\leq b}{\rA}
    \proofstepwidestar[3,8]{\exists n\in\mathbb N\,
      \IntClass{a}{b}=-\IntEmbed(n)}{%
      \FormulaRefAuto{IntegerNormalFormNonpositive}{3,8}}
    \proofstepwidestar[3,8]{\exists n\in\mathbb N\,
      (z=\IntEmbed(n)\lor z=-\IntEmbed(n))}{\rIE{3,9}}
    \proofstepwidestar[3]{\exists n\in\mathbb N\,
      (z=\IntEmbed(n)\lor z=-\IntEmbed(n))}{%
      \rOE{4,5,7,8,10}}
    \proofstepwidestar[1]{\exists n\in\mathbb N\,
      (z=\IntEmbed(n)\lor z=-\IntEmbed(n))}{\rEE{2,3,11}}
  \closeproofpartwide
\end{tabproofsplitwide}

\chapter{Ordnung der ganzen Zahlen}

\section{Translationsinvarianz der natürlichen Ordnung}

\FormulaThmDeltaK[Translationsinvarianz der natürlichen Ordnung]{%
  a,b,c\in\mathbb N\vdash
  a\leq b\leftrightarrow a+c\leq b+c%
}{NaturalOrderTranslationForIntegers}{%
  \DeltaRow{Mengen}{a\dsep b\dsep c\dsep k}%
}
\begin{tabproofsplitwide}
  \proofpartwideindR[Erhaltung unter Translation]{%
    a,b,c\in\mathbb N\dsep a\leq b
    \vdash a+c\leq b+c%
  }{%
    a,b,c\in\mathbb N\dsep a\leq b
    \vdash a+c\leq b+c%
  }[NaturalOrderTranslationForward]
    \proofstepwidestar[1]{a,b,c\in\mathbb N}{\rA}
    \proofstepwidestar[2]{a\leq b}{\rA}
    \proofstepwidestar[1,2]{\exists k\in\mathbb N\,(a+k=b)}{%
      \FormulaRefAuto{m\in\mathbb{N}\dsep n\in\mathbb{N}\vdash
        m\leq n\leftrightarrow
        \exists k\in\mathbb{N}\,(m+k=n)}{1,2}}
    \proofstepwidestar[4]{k\in\mathbb N\land a+k=b}{\rA}
    \proofstepwide[1,4]{(a+c)+k}{=}{(a+k)+c}{%
      \FormulaRefAuto{PeanoAddInterchange}{1}}
    \proofstepwide[1,4]{(a+k)+c}{=}{b+c}{\rIE{\rAEb{4}}}
    \proofstepwidestar[1,4]{(a+c)+k=b+c}{%
      \FormulaRefAuto{a=b\dsep b=c\vdash a=c}{5,6}}
    \proofstepwidestar[1,4]{a+c\leq b+c}{%
      \FormulaRefAuto{m\in\mathbb{N}\dsep n\in\mathbb{N}\vdash
        m\leq n\leftrightarrow
        \exists k\in\mathbb{N}\,(m+k=n)}{1,\rEI{\rAI{\rAEa{4},7}}}}
    \proofstepwidestar[1,2]{a+c\leq b+c}{\rEE{3,4,8}}
  \closeproofpartwideindR

  \proofpartwideindR[Kürzen einer Translation]{%
    a,b,c\in\mathbb N\dsep a+c\leq b+c
    \vdash a\leq b%
  }{%
    a,b,c\in\mathbb N\dsep a+c\leq b+c
    \vdash a\leq b%
  }[NaturalOrderTranslationBackward]
    \proofstepwidestar[1]{a,b,c\in\mathbb N}{\rA}
    \proofstepwidestar[2]{a+c\leq b+c}{\rA}
    \proofstepwidestar[1,2]{\exists k\in\mathbb N\,((a+c)+k=b+c)}{%
      \FormulaRefAuto{m\in\mathbb{N}\dsep n\in\mathbb{N}\vdash
        m\leq n\leftrightarrow
        \exists k\in\mathbb{N}\,(m+k=n)}{1,2}}
    \proofstepwidestar[4]{k\in\mathbb N\land(a+c)+k=b+c}{\rA}
    \proofstepwide[1,4]{(a+k)+c}{=}{(a+c)+k}{%
      \FormulaRefAuto{PeanoAddInterchange}{1}}
    \proofstepwide[1,4]{(a+c)+k}{=}{b+c}{\rAEb{4}}
    \proofstepwidestar[1,4]{(a+k)+c=b+c}{%
      \FormulaRefAuto{a=b\dsep b=c\vdash a=c}{5,6}}
    \proofstepwidestar[1,4]{a+k=b}{%
      \FormulaRefAuto{PeanoAddRightCancellation}{%
        \FormulaRefAuto{PeanoAddClosure}{1,4},1,1,7}}
    \proofstepwidestar[1,4]{a\leq b}{%
      \FormulaRefAuto{m\in\mathbb{N}\dsep n\in\mathbb{N}\vdash
        m\leq n\leftrightarrow
        \exists k\in\mathbb{N}\,(m+k=n)}{1,\rEI{\rAI{\rAEa{4},8}}}}
    \proofstepwidestar[1,2]{a\leq b}{\rEE{3,4,9}}
  \closeproofpartwideindR

  \proofpartwide{Schluss}
    \proofstepwidestar[1]{a,b,c\in\mathbb N}{\rA}
    \proofstepwidestar[1]{a\leq b\rightarrow a+c\leq b+c}{%
      \FormulaRefAuto{NaturalOrderTranslationForward}{1}}
    \proofstepwidestar[1]{a+c\leq b+c\rightarrow a\leq b}{%
      \FormulaRefAuto{NaturalOrderTranslationBackward}{1}}
    \proofstepwidestar[1]{a\leq b\leftrightarrow a+c\leq b+c}{%
      \rLRI{2,3}}
  \closeproofpartwide
\end{tabproofsplitwide}

\FormulaThmDeltaK[Addition natürlicher Ungleichungen]{%
  a,b,c,d\in\mathbb N\dsep a\leq b\dsep c\leq d
  \vdash a+c\leq b+d%
}{NaturalOrderAdditionForIntegers}{%
  \DeltaRow{Mengen}{a\dsep b\dsep c\dsep d}%
}
\begin{tabproofsplitwide}
  \proofpartwideindR[Translation der ersten Ungleichung]{%
    a,b,c\in\mathbb N\dsep a\leq b
    \vdash a+c\leq b+c%
  }{%
    a,b,c\in\mathbb N\dsep a\leq b
    \vdash a+c\leq b+c%
  }[NaturalOrderAdditionFirstTranslation]
    \proofstepwidestar[1]{a,b,c\in\mathbb N}{\rA}
    \proofstepwidestar[2]{a\leq b}{\rA}
    \proofstepwidestar[1,2]{a+c\leq b+c}{%
      \FormulaRefAuto{NaturalOrderTranslationForward}{1,2}}
  \closeproofpartwideindR

  \proofpartwideindR[Translation der zweiten Ungleichung]{%
    b,c,d\in\mathbb N\dsep c\leq d
    \vdash b+c\leq b+d%
  }{%
    b,c,d\in\mathbb N\dsep c\leq d
    \vdash b+c\leq b+d%
  }[NaturalOrderAdditionSecondTranslation]
    \proofstepwidestar[1]{b,c,d\in\mathbb N}{\rA}
    \proofstepwidestar[2]{c\leq d}{\rA}
    \proofstepwidestar[1,2]{c+b\leq d+b}{%
      \FormulaRefAuto{NaturalOrderTranslationForward}{1,2}}
    \proofstepwidestar[1]{b+c=c+b}{%
      \FormulaRefAuto{PeanoAddCommutative}{1}}
    \proofstepwidestar[1]{d+b=b+d}{%
      \FormulaRefAuto{PeanoAddCommutative}{1}}
    \proofstepwidestar[1,2]{b+c\leq b+d}{\rIE{4,3,5}}
  \closeproofpartwideindR

  \proofpartwideindR[Verknüpfung durch Transitivität]{%
    a,b,c,d\in\mathbb N\dsep a\leq b\dsep c\leq d
    \vdash a+c\leq b+d%
  }{%
    a,b,c,d\in\mathbb N\dsep a\leq b\dsep c\leq d
    \vdash a+c\leq b+d%
  }[NaturalOrderAdditionTransitiveConclusion]
    \proofstepwidestar[1]{a,b,c,d\in\mathbb N}{\rA}
    \proofstepwidestar[2]{a\leq b}{\rA}
    \proofstepwidestar[3]{c\leq d}{\rA}
    \proofstepwidestar[1,2]{a+c\leq b+c}{%
      \FormulaRefAuto{NaturalOrderAdditionFirstTranslation}{1,2}}
    \proofstepwidestar[1,3]{b+c\leq b+d}{%
      \FormulaRefAuto{NaturalOrderAdditionSecondTranslation}{1,3}}
    \proofstepwidestar[1]{a+c\in\mathbb N}{%
      \FormulaRefAuto{PeanoAddClosure}{1}}
    \proofstepwidestar[1]{b+c\in\mathbb N}{%
      \FormulaRefAuto{PeanoAddClosure}{1}}
    \proofstepwidestar[1]{b+d\in\mathbb N}{%
      \FormulaRefAuto{PeanoAddClosure}{1}}
    \proofstepwidestar[1,2,3]{a+c\leq b+d}{%
      \FormulaRefAuto{%
        k\in\mathbb{N}\dsep m\in\mathbb{N}\dsep n\in\mathbb{N}\dsep
        k\leq m\dsep m\leq n\vdash k\leq n%
      }{6,7,8,4,5}}
  \closeproofpartwideindR
\end{tabproofsplitwide}

\section{Quotientenordnung}

\FormulaThmDeltaK[Repräsentantenunabhängigkeit der Ganzzahlordnung]{%
  \begin{aligned}[t]
    &a,b,a',b',c,d,c',d'\in\mathbb N\dsep
      \IntClass{a}{b}=\IntClass{a'}{b'}\dsep
      \IntClass{c}{d}=\IntClass{c'}{d'}\vdash{}\\
    &(a+d\leq c+b)\leftrightarrow(a'+d'\leq c'+b')
  \end{aligned}%
}{IntegerOrderCompatible}{%
  \DeltaRow{Mengen}{a\dsep b\dsep a'\dsep b'\dsep
    c\dsep d\dsep c'\dsep d'}%
}
\begin{tabproofsplitwide}
  \proofpartwideindR[Hinrichtung]{%
    \begin{aligned}[t]
      &a,b,a',b',c,d,c',d'\in\mathbb N\dsep
      a+b'=b+a'\dsep c+d'=d+c'\dsep a+d\leq c+b\vdash{}\\
      &a'+d'\leq c'+b'
    \end{aligned}
  }{%
    a,b,a',b',c,d,c',d'\in\mathbb N\dsep
    a+b'=b+a'\dsep c+d'=d+c'\dsep a+d\leq c+b
    \vdash a'+d'\leq c'+b'
  }[IntegerOrderCompatibleForward]
    \proofstepwidestar[1]{a,b,a',b',c,d,c',d'\in\mathbb N}{\rA}
    \proofstepwidestar[2]{a+b'=b+a'}{\rA}
    \proofstepwidestar[3]{c+d'=d+c'}{\rA}
    \proofstepwidestar[4]{a+d\leq c+b}{\rA}
    \proofstepwidestar[1,4]{(a+d)+(b'+c')\leq(c+b)+(b'+c')}{%
      \FormulaRefAuto{NaturalOrderTranslationForward}{1,4}}
    \proofstepwidestar[1,2,3]{%
      (a+d)+(b'+c')=(a'+d')+(b+c)}{%
      \FormulaRefAuto{PeanoAddInterchange}{1,2,3}}
    \proofstepwidestar[1]{%
      (c+b)+(b'+c')=(c'+b')+(b+c)}{%
      \FormulaRefAuto{PeanoAddInterchange}{1}}
    \proofstepwidestar[1,2,3,4]{%
      (a'+d')+(b+c)\leq(c'+b')+(b+c)}{\rIE{5,6,7}}
    \proofstepwidestar[1,2,3,4]{a'+d'\leq c'+b'}{%
      \FormulaRefAuto{NaturalOrderTranslationBackward}{1,8}}
  \closeproofpartwideindR

  \proofpartwideindR[Rückrichtung]{%
    \begin{aligned}[t]
      &a,b,a',b',c,d,c',d'\in\mathbb N\dsep
      a+b'=b+a'\dsep c+d'=d+c'\dsep a'+d'\leq c'+b'\vdash{}\\
      &a+d\leq c+b
    \end{aligned}
  }{%
    a,b,a',b',c,d,c',d'\in\mathbb N\dsep
    a+b'=b+a'\dsep c+d'=d+c'\dsep a'+d'\leq c'+b'
    \vdash a+d\leq c+b
  }[IntegerOrderCompatibleBackward]
    \proofstepwidestar[1]{a,b,a',b',c,d,c',d'\in\mathbb N}{\rA}
    \proofstepwidestar[2]{a+b'=b+a'}{\rA}
    \proofstepwidestar[3]{c+d'=d+c'}{\rA}
    \proofstepwidestar[4]{a'+d'\leq c'+b'}{\rA}
    \proofstepwidestar[1,2]{a'+b=b'+a}{%
      \FormulaRefAuto{a=b\vdash b=a}{%
        \FormulaRefAuto{PeanoAddCommutative}{2}}}
    \proofstepwidestar[1,3]{c'+d=d'+c}{%
      \FormulaRefAuto{a=b\vdash b=a}{%
        \FormulaRefAuto{PeanoAddCommutative}{3}}}
    \proofstepwidestar[1,2,3,4]{a+d\leq c+b}{%
      \FormulaRefAuto{IntegerOrderCompatibleForward}{1,5,6,4}}
  \closeproofpartwideindR

  \proofpartwide{Schluss}
    \proofstepwidestar[1]{a,b,a',b',c,d,c',d'\in\mathbb N}{\rA}
    \proofstepwidestar[2]{\IntClass{a}{b}=\IntClass{a'}{b'}}{\rA}
    \proofstepwidestar[3]{\IntClass{c}{d}=\IntClass{c'}{d'}}{\rA}
    \proofstepwidestar[1,2]{a+b'=b+a'}{%
      \FormulaRefAuto{IntegerClassEquality}{1,2}}
    \proofstepwidestar[1,3]{c+d'=d+c'}{%
      \FormulaRefAuto{IntegerClassEquality}{1,3}}
    \proofstepwidestar[1,2,3]{%
      a+d\leq c+b\rightarrow a'+d'\leq c'+b'}{%
      \FormulaRefAuto{IntegerOrderCompatibleForward}{1,4,5}}
    \proofstepwidestar[1,2,3]{%
      a'+d'\leq c'+b'\rightarrow a+d\leq c+b}{%
      \FormulaRefAuto{IntegerOrderCompatibleBackward}{1,4,5}}
    \proofstepwidestar[1,2,3]{%
      (a+d\leq c+b)\leftrightarrow(a'+d'\leq c'+b')}{\rLRI{6,7}}
  \closeproofpartwide
\end{tabproofsplitwide}

\FormulaDefDeltaK[Ordnungsrelation der ganzen Zahlen]{%
  \leq_{\mathbb Z}\coloneqq
  \bigl\{(\IntClass{a}{b},\IntClass{c}{d})\in
  \mathbb Z\times\mathbb Z\mid a+d\leq c+b\bigr\}%
}{IntegerOrderDef}{%
  \DeltaRow{Mengen}{a\dsep b\dsep c\dsep d}%
  \DeltaRow{Relation}{\leq_{\mathbb Z}}%
  \DeltaRow{Neue Symbole}{\leq_{\mathbb Z}}%
}

\FormulaDefDeltaK[Notation der Ganzzahlordnung]{%
  a,b,c,d\in\mathbb N\vdash
  \IntClass{a}{b}\leq\IntClass{c}{d}
  \coloneqq a+d\leq c+b%
}{IntegerOrderNotation}{%
  \DeltaRow{Mengen}{a\dsep b\dsep c\dsep d}%
  \DeltaRow{Relation}{\leq}%
}

\FormulaThmDeltaK[Trägerbedingung der Ganzzahlordnung]{%
  \leq_{\mathbb Z}\subseteq\mathbb Z\times\mathbb Z%
}{IntegerOrderCarrier}{%
  \DeltaRow{Mengen}{x\dsep a\dsep b\dsep c\dsep d}%
  \DeltaRow{Relation}{\leq_{\mathbb Z}}%
}
\begin{tabproof}
  \proofstep{1}{x\in\leq_{\mathbb Z}}{\rA}
  \proofstep{}{%
    \leq_{\mathbb Z}=
    \bigl\{(\IntClass{a}{b},\IntClass{c}{d})\in
    \mathbb Z\times\mathbb Z\mid a+d\leq c+b\bigr\}}{%
    \FormulaRefAuto{IntegerOrderDef}}
  \proofstep{1}{x\in\mathbb Z\times\mathbb Z}{%
    \FormulaRefAuto{x \in \{x \in A \mid P(x)\}\vdash x\in A}{2,1}}
  \proofstep{}{\leq_{\mathbb Z}\subseteq\mathbb Z\times\mathbb Z}{%
    \FormulaRefAuto{A\subseteq B \coloneq
      \forall x\,(x\in A\rightarrow x\in B)}{\rUI{\rRI{1,3}}}}
\end{tabproof}

\FormulaThmDeltaK[Totale Ordnung der ganzen Zahlen]{%
  \TotOrd{\leq_{\mathbb Z},\mathbb Z}%
}{IntegerOrderTotal}{%
  \DeltaRow{Mengen}{z\dsep w\dsep u}%
  \DeltaRow{Relation}{\leq_{\mathbb Z}}%
}
\begin{tabproofsplitwide}
  \proofpartwideindR[Reflexivität]{%
    z\in\mathbb Z\vdash z\leq z%
  }{%
    z\in\mathbb Z\vdash z\leq z%
  }[IntegerOrderReflexive]
    \proofstepwidestar[1]{z\in\mathbb Z}{\rA}
    \proofstepwidestar[1]{\exists a,b\in\mathbb N\,
      z=\IntClass{a}{b}}{\FormulaRefAuto{IntegerPairRepresentative}{1}}
    \proofstepwidestar[3]{a,b\in\mathbb N\land z=\IntClass{a}{b}}{\rA}
    \proofstepwidestar[3]{a+b\leq a+b}{%
      \FormulaRefAuto{PeanoLeqReflexive}{%
        \FormulaRefAuto{PeanoAddClosure}{3}}}
    \proofstepwidestar[3]{\IntClass{a}{b}\leq\IntClass{a}{b}}{%
      \FormulaRefAuto{IntegerOrderNotation}{3,4}}
    \proofstepwidestar[3]{z\leq z}{\rIE{3,5}}
    \proofstepwidestar[1]{z\leq z}{\rEE{2,3,6}}
  \closeproofpartwideindR

  \proofpartwideindR[Antisymmetrie]{%
    z,w\in\mathbb Z\dsep z\leq w\dsep w\leq z
    \vdash z=w%
  }{%
    z,w\in\mathbb Z\dsep z\leq w\dsep w\leq z
    \vdash z=w%
  }[IntegerOrderAntisymmetric]
    \proofstepwidestar[1]{z,w\in\mathbb Z}{\rA}
    \proofstepwidestar[2]{z\leq w}{\rA}
    \proofstepwidestar[3]{w\leq z}{\rA}
    \proofstepwidestar[1]{z\in\mathbb Z}{\rAEa{1}}
    \proofstepwidestar[1]{w\in\mathbb Z}{\rAEb{1}}
    \proofstepwidestar[1]{\exists a,b\in\mathbb N\,
      z=\IntClass{a}{b}}{\FormulaRefAuto{IntegerPairRepresentative}{4}}
    \proofstepwidestar[1]{\exists c,d\in\mathbb N\,
      w=\IntClass{c}{d}}{\FormulaRefAuto{IntegerPairRepresentative}{5}}
    \proofstepwidestar[8]{a,b,c,d\in\mathbb N\land
      z=\IntClass{a}{b}\land w=\IntClass{c}{d}}{\rA}
    \proofstepwidestar[2,8]{a+d\leq c+b}{%
      \FormulaRefAuto{IntegerOrderNotation}{8,2}}
    \proofstepwidestar[3,8]{c+b\leq a+d}{%
      \FormulaRefAuto{IntegerOrderNotation}{8,3}}
    \proofstepwidestar[8]{a+d\in\mathbb N}{%
      \FormulaRefAuto{PeanoAddClosure}{8}}
    \proofstepwidestar[8]{c+b\in\mathbb N}{%
      \FormulaRefAuto{PeanoAddClosure}{8}}
    \proofstepwidestar[2,3,8]{a+d=c+b}{%
      \FormulaRefAuto{%
        m\in\mathbb{N}\dsep n\in\mathbb{N}\dsep
        m\leq n\dsep n\leq m\vdash m=n%
      }{11,12,9,10}}
    \proofstepwidestar[2,3,8]{%
      \IntClass{a}{b}=\IntClass{c}{d}}{%
      \FormulaRefAuto{IntegerClassEquality}{8,13}}
    \proofstepwidestar[2,3,8]{z=w}{\rIE{8,14}}
    \proofstepwidestar[1,2,3]{z=w}{\rEE{6,7,8,15}}
  \closeproofpartwideindR

  \proofpartwideindR[Transitivität]{%
    z,w,u\in\mathbb Z\dsep z\leq w\dsep w\leq u
    \vdash z\leq u%
  }{%
    z,w,u\in\mathbb Z\dsep z\leq w\dsep w\leq u
    \vdash z\leq u%
  }[IntegerOrderTransitive]
    \proofstepwidestar[1]{z,w,u\in\mathbb Z}{\rA}
    \proofstepwidestar[2]{z\leq w}{\rA}
    \proofstepwidestar[3]{w\leq u}{\rA}
    \proofstepwidestar[1]{z\in\mathbb Z}{\rAEa{1}}
    \proofstepwidestar[1]{w\in\mathbb Z}{\rAEa{\rAEb{1}}}
    \proofstepwidestar[1]{u\in\mathbb Z}{\rAEb{\rAEb{1}}}
    \proofstepwidestar[4]{\exists a,b\in\mathbb N\,
      z=\IntClass{a}{b}}{%
      \FormulaRefAuto{IntegerPairRepresentative}{4}}
    \proofstepwidestar[5]{\exists c,d\in\mathbb N\,
      w=\IntClass{c}{d}}{%
      \FormulaRefAuto{IntegerPairRepresentative}{5}}
    \proofstepwidestar[6]{\exists e,f\in\mathbb N\,
      u=\IntClass{e}{f}}{%
      \FormulaRefAuto{IntegerPairRepresentative}{6}}
    \proofstepwidestar[10]{a,b,c,d,e,f\in\mathbb N\land
      z=\IntClass{a}{b}\land w=\IntClass{c}{d}\land
      u=\IntClass{e}{f}}{\rA}
    \proofstepwidestar[2,10]{a+d\leq c+b}{%
      \FormulaRefAuto{IntegerOrderNotation}{10,2}}
    \proofstepwidestar[3,10]{c+f\leq e+d}{%
      \FormulaRefAuto{IntegerOrderNotation}{10,3}}
    \proofstepwidestar[2,3,10]{%
      (a+d)+(c+f)\leq(c+b)+(e+d)}{%
      \FormulaRefAuto{NaturalOrderAdditionForIntegers}{10,11,12}}
    \proofstepwidestar[10]{%
      (a+d)+(c+f)=(a+f)+(c+d)}{%
      \FormulaRefAuto{PeanoAddInterchange}{10}}
    \proofstepwidestar[10]{%
      (c+b)+(e+d)=(b+e)+(c+d)}{%
      \FormulaRefAuto{PeanoAddInterchange}{10}}
    \proofstepwidestar[2,3,10]{(a+f)+(c+d)\leq(b+e)+(c+d)}{%
      \rIE{13,14,15}}
    \proofstepwidestar[2,3,10]{a+f\leq b+e}{%
      \FormulaRefAuto{NaturalOrderTranslationBackward}{10,16}}
    \proofstepwidestar[2,3,10]{z\leq u}{%
      \FormulaRefAuto{IntegerOrderNotation}{10,17}}
    \proofstepwidestar[1,2,3]{z\leq u}{\rEE{7,8,9,10,18}}
  \closeproofpartwideindR

  \proofpartwideindR[Totalität]{%
    z,w\in\mathbb Z\vdash z\leq w\lor w\leq z%
  }{%
    z,w\in\mathbb Z\vdash z\leq w\lor w\leq z%
  }[IntegerOrderComparable]
    \proofstepwidestar[1]{z,w\in\mathbb Z}{\rA}
    \proofstepwidestar[1]{z\in\mathbb Z}{\rAEa{1}}
    \proofstepwidestar[1]{w\in\mathbb Z}{\rAEb{1}}
    \proofstepwidestar[2]{\exists a,b\in\mathbb N\,
      z=\IntClass{a}{b}}{%
      \FormulaRefAuto{IntegerPairRepresentative}{2}}
    \proofstepwidestar[3]{\exists c,d\in\mathbb N\,
      w=\IntClass{c}{d}}{%
      \FormulaRefAuto{IntegerPairRepresentative}{3}}
    \proofstepwidestar[6]{a,b,c,d\in\mathbb N\land
      z=\IntClass{a}{b}\land w=\IntClass{c}{d}}{\rA}
    \proofstepwidestar[6]{a+d\leq c+b\lor c+b\leq a+d}{%
      \FormulaRefAuto{PeanoOrderComparison}{%
        \FormulaRefAuto{PeanoAddClosure}{6},%
        \FormulaRefAuto{PeanoAddClosure}{6}}}
    \proofstepwidestar[6]{z\leq w\lor w\leq z}{%
      \FormulaRefAuto{IntegerOrderNotation}{6,7}}
    \proofstepwidestar[1]{z\leq w\lor w\leq z}{\rEE{4,5,6,8}}
  \closeproofpartwideindR

  \proofpartwide{Schluss}
    \proofstepwidestar[]{\leq_{\mathbb Z}\subseteq
      \mathbb Z\times\mathbb Z}{%
      \FormulaRefAuto{IntegerOrderCarrier}}
    \proofstepwidestar[]{\forall z\in\mathbb Z\,(z\leq z)}{%
      \FormulaRefAuto{IntegerOrderReflexive}}
    \proofstepwidestar[]{%
      z,w,u\in\mathbb Z\dsep z\leq w\dsep w\leq u\vdash z\leq u}{%
      \FormulaRefAuto{IntegerOrderTransitive}}
    \proofstepwidestar[]{%
      z,w\in\mathbb Z\dsep z\leq w\dsep w\leq z\vdash z=w}{%
      \FormulaRefAuto{IntegerOrderAntisymmetric}}
    \proofstepwidestar[]{\PartOrd{\leq_{\mathbb Z},\mathbb Z}}{%
      \FormulaRefAuto{Partielle Ordnung}{1,2,3,4}}
    \proofstepwidestar[]{z,w\in\mathbb Z\vdash z\leq w\lor w\leq z}{%
      \FormulaRefAuto{IntegerOrderComparable}}
    \proofstepwidestar[]{\TotOrd{\leq_{\mathbb Z},\mathbb Z}}{%
      \FormulaRefAuto{Totale Ordnung}{5,6}}
  \closeproofpartwide
\end{tabproofsplitwide}

\FormulaDefDeltaK[Strikte Ordnung der ganzen Zahlen]{%
  z,w\in\mathbb Z\vdash
  z<w\coloneqq z\leq w\land z\neq w%
}{IntegerStrictOrderDef}{%
  \DeltaRow{Mengen}{z\dsep w}%
  \DeltaRow{Relation}{<}%
}

\FormulaThmDeltaK[Jede ganze Zahl ist kleiner als ihr Nachfolger]{%
  z\in\mathbb Z\vdash z<\IntSucc(z)%
}{IntegerSuccessorPositive}{%
  \DeltaRow{Mengen}{z\dsep a\dsep b}%
}
\begin{tabproofsplitwide}
  \proofpartwideindR[Ordnungsteil]{%
    z\in\mathbb Z\vdash z\leq\IntSucc(z)%
  }{%
    z\in\mathbb Z\vdash z\leq\IntSucc(z)%
  }[IntegerSuccessorLeq]
    \proofstepwidestar[1]{z\in\mathbb Z}{\rA}
    \proofstepwidestar[1]{\exists a,b\in\mathbb N\,
      z=\IntClass{a}{b}}{\FormulaRefAuto{IntegerPairRepresentative}{1}}
    \proofstepwidestar[3]{a,b\in\mathbb N\land z=\IntClass{a}{b}}{\rA}
    \proofstepwidestar[3]{\IntSucc(z)=\IntClass{\Succ(a)}{b}}{%
      \FormulaRefAuto{IntegerSuccessorClassValue}{3}}
    \proofstepwidestar[3]{a\leq\Succ(a)}{%
      \FormulaRefAuto{PeanoLeqAddOne}{3}}
    \proofstepwidestar[3]{a+b\leq\Succ(a)+b}{%
      \FormulaRefAuto{NaturalOrderTranslationForward}{3,5}}
    \proofstepwidestar[3]{z\leq\IntSucc(z)}{%
      \FormulaRefAuto{IntegerOrderNotation}{3,4,6}}
    \proofstepwidestar[1]{z\leq\IntSucc(z)}{\rEE{2,3,7}}
  \closeproofpartwideindR

  \proofpartwideindR[Ungleichheitsteil]{%
    z\in\mathbb Z\vdash z\neq\IntSucc(z)%
  }{%
    z\in\mathbb Z\vdash z\neq\IntSucc(z)%
  }[IntegerSuccessorDistinct]
    \proofstepwidestar[1]{z\in\mathbb Z}{\rA}
    \proofstepwidestar[2]{z=\IntSucc(z)}{\rA}
    \proofstepwidestar[1]{\exists a,b\in\mathbb N\,
      z=\IntClass{a}{b}}{\FormulaRefAuto{IntegerPairRepresentative}{1}}
    \proofstepwidestar[4]{a,b\in\mathbb N\land z=\IntClass{a}{b}}{\rA}
    \proofstepwidestar[4]{\IntSucc(z)=\IntClass{\Succ(a)}{b}}{%
      \FormulaRefAuto{IntegerSuccessorClassValue}{4}}
    \proofstepwidestar[2,4]{\IntClass{a}{b}=\IntClass{\Succ(a)}{b}}{\rIE{4,5,2}}
    \proofstepwidestar[2,4]{a+b=b+\Succ(a)}{%
      \FormulaRefAuto{IntegerClassEquality}{4,6}}
    \proofstepwidestar[4]{b+\Succ(a)=\Succ(a)+b}{%
      \FormulaRefAuto{PeanoAddCommutative}{4}}
    \proofstepwidestar[2,4]{a+b=\Succ(a)+b}{\rIE{7,8}}
    \proofstepwidestar[2,4]{a=\Succ(a)}{%
      \FormulaRefAuto{PeanoAddRightCancellation}{4,9}}
    \proofstepwidestar[4]{a<\Succ(a)}{%
      \FormulaRefAuto{n\in\mathbb{N}\vdash n<\Succ(n)}{4}}
    \proofstepwidestar[2,4]{\bot}{%
      \FormulaRefAuto{PeanoStrictEqualityImpossible}{4,11,10}}
    \proofstepwidestar[1,2]{\bot}{\rEE{3,4,12}}
    \proofstepwidestar[1]{z\neq\IntSucc(z)}{\rCI{2,13}}
  \closeproofpartwideindR

  \proofpartwide{Schluss}
    \proofstepwidestar[1]{z\in\mathbb Z}{\rA}
    \proofstepwidestar[1]{z\leq\IntSucc(z)}{%
      \FormulaRefAuto{IntegerSuccessorLeq}{1}}
    \proofstepwidestar[1]{z\neq\IntSucc(z)}{%
      \FormulaRefAuto{IntegerSuccessorDistinct}{1}}
    \proofstepwidestar[1]{z<\IntSucc(z)}{%
      \FormulaRefAuto{IntegerStrictOrderDef}{1,\rAI{2,3}}}
  \closeproofpartwide
\end{tabproofsplitwide}

\chapter{Erweiterte Peano-Axiome der ganzen Zahlen}

\section{Herleitung der zweiseitigen Induktion}

\FormulaThmDeltaK[Zweiseitige Induktion im Quotientenmodell]{%
  \begin{aligned}[t]
    &P(0_{\mathbb Z})\dsep
    \forall z\in\mathbb Z\,
      (P(z)\rightarrow P(\IntSucc(z)))\dsep{}\\
    &\forall z\in\mathbb Z\,
      (P(z)\rightarrow P(\IntPred(z)))
    \vdash \forall z\in\mathbb Z\,P(z)
  \end{aligned}%
}{IntegerTwoSidedInductionModel}{%
  \DeltaRow{Einstelliges Prädikat}{P}%
  \DeltaRow{Mengen}{z\dsep n}%
}
\begin{tabproofsplitwide}
  \proofpartwideindR[Induktion auf dem nichtnegativen Strahl]{%
    \begin{aligned}[t]
      &P(0_{\mathbb Z})\dsep
      \forall z\in\mathbb Z\,
      (P(z)\rightarrow P(\IntSucc(z)))\vdash{}\\
      &\forall n\in\mathbb N\,P(\IntEmbed(n))
    \end{aligned}
  }{%
    P(0_{\mathbb Z})\dsep
    \forall z\in\mathbb Z\,
    (P(z)\rightarrow P(\IntSucc(z)))
    \vdash\forall n\in\mathbb N\,P(\IntEmbed(n))
  }[IntegerTwoSidedInductionPositiveRay]
    \proofstepwidestar[1]{P(0_{\mathbb Z})}{\rA}
    \proofstepwidestar[2]{\forall z\in\mathbb Z\,
      (P(z)\rightarrow P(\IntSucc(z)))}{\rA}
    \proofstepwidestar[]{0_{\mathbb Z}=\IntEmbed(0)}{%
      \FormulaRefAuto{IntegerZeroDef}}
    \proofstepwidestar[1]{P(\IntEmbed(0))}{\rIE{3,1}}

    \proofstepwidestar[5]{n\in\mathbb N}{\rA}
    \proofstepwidestar[6]{P(\IntEmbed(n))}{\rA}
    \proofstepwidestar[5]{\IntEmbed(n)\in\mathbb Z}{%
      \FormulaRefAuto{F\colon A\to B\dsep x\in A\vdash F(x)\in B}{%
        \FormulaRefAuto{NaturalIntegerEmbeddingFunction},5}}
    \proofstepwidestar[2,5]{%
      P(\IntEmbed(n))\rightarrow
      P(\IntSucc(\IntEmbed(n)))}{\rRE{\rUE{2},7}}
    \proofstepwidestar[2,5,6]{P(\IntSucc(\IntEmbed(n)))}{\rRE{8,6}}
    \proofstepwidestar[5]{%
      \IntSucc(\IntEmbed(n))=\IntEmbed(\Succ(n))}{%
      \FormulaRefAuto{IntegerSuccessorOnEmbedding}{5}}
    \proofstepwidestar[2,5,6]{P(\IntEmbed(\Succ(n)))}{\rIE{10,9}}
    \proofstepwidestar[2,5]{%
      P(\IntEmbed(n))\rightarrow P(\IntEmbed(\Succ(n)))}{\rRI{6,11}}
    \proofstepwidestar[2]{\forall n\in\mathbb N\,
      (P(\IntEmbed(n))\rightarrow P(\IntEmbed(\Succ(n))))}{%
      \rUI{\rRI{5,12}}}
    \proofstepwidestar[1,2]{\forall n\in\mathbb N\,
      P(\IntEmbed(n))}{%
      \FormulaRefAuto{PeanoInductionAxiom}{4,13}}
  \closeproofpartwideindR

  \proofpartwideindR[Induktion auf dem nichtpositiven Strahl]{%
    \begin{aligned}[t]
      &P(0_{\mathbb Z})\dsep
      \forall z\in\mathbb Z\,
      (P(z)\rightarrow P(\IntPred(z)))\vdash{}\\
      &\forall n\in\mathbb N\,P(-\IntEmbed(n))
    \end{aligned}
  }{%
    P(0_{\mathbb Z})\dsep
    \forall z\in\mathbb Z\,
    (P(z)\rightarrow P(\IntPred(z)))
    \vdash\forall n\in\mathbb N\,P(-\IntEmbed(n))
  }[IntegerTwoSidedInductionNegativeRay]
    \proofstepwidestar[1]{P(0_{\mathbb Z})}{\rA}
    \proofstepwidestar[2]{\forall z\in\mathbb Z\,
      (P(z)\rightarrow P(\IntPred(z)))}{\rA}
    \proofstepwidestar[]{\IntEmbed(0)=0_{\mathbb Z}}{%
      \FormulaRefAuto{a=b\vdash b=a}{\FormulaRefAuto{IntegerZeroDef}}}
    \proofstepwidestar[]{-0_{\mathbb Z}=0_{\mathbb Z}}{%
      \FormulaRefAuto{IntegerNegativeZero}}
    \proofstepwidestar[1]{P(-\IntEmbed(0))}{\rIE{3,4,1}}

    \proofstepwidestar[6]{n\in\mathbb N}{\rA}
    \proofstepwidestar[7]{P(-\IntEmbed(n))}{\rA}
    \proofstepwidestar[6]{-\IntEmbed(n)\in\mathbb Z}{%
      \FormulaRefAuto{IntNegClosure}{%
        \FormulaRefAuto{F\colon A\to B\dsep x\in A\vdash F(x)\in B}{%
          \FormulaRefAuto{NaturalIntegerEmbeddingFunction},6}}}
    \proofstepwidestar[2,6]{%
      P(-\IntEmbed(n))\rightarrow
      P(\IntPred(-\IntEmbed(n)))}{\rRE{\rUE{2},8}}
    \proofstepwidestar[2,6,7]{P(\IntPred(-\IntEmbed(n)))}{\rRE{9,7}}
    \proofstepwidestar[6]{%
      \IntPred(-\IntEmbed(n))=-\IntEmbed(\Succ(n))}{%
      \FormulaRefAuto{IntegerPredecessorOnNegativeEmbedding}{6}}
    \proofstepwidestar[2,6,7]{P(-\IntEmbed(\Succ(n)))}{\rIE{11,10}}
    \proofstepwidestar[2,6]{%
      P(-\IntEmbed(n))\rightarrow P(-\IntEmbed(\Succ(n)))}{\rRI{7,12}}
    \proofstepwidestar[2]{\forall n\in\mathbb N\,
      (P(-\IntEmbed(n))\rightarrow P(-\IntEmbed(\Succ(n))))}{%
      \rUI{\rRI{6,13}}}
    \proofstepwidestar[1,2]{\forall n\in\mathbb N\,
      P(-\IntEmbed(n))}{%
      \FormulaRefAuto{PeanoInductionAxiom}{5,14}}
  \closeproofpartwideindR

  \proofpartwide{Schluss mittels Normalform}
    \proofstepwidestar[1]{P(0_{\mathbb Z})}{\rA}
    \proofstepwidestar[2]{\forall z\in\mathbb Z\,
      (P(z)\rightarrow P(\IntSucc(z)))}{\rA}
    \proofstepwidestar[3]{\forall z\in\mathbb Z\,
      (P(z)\rightarrow P(\IntPred(z)))}{\rA}
    \proofstepwidestar[1,2]{\forall n\in\mathbb N\,
      P(\IntEmbed(n))}{%
      \FormulaRefAuto{IntegerTwoSidedInductionPositiveRay}{1,2}}
    \proofstepwidestar[1,3]{\forall n\in\mathbb N\,
      P(-\IntEmbed(n))}{%
      \FormulaRefAuto{IntegerTwoSidedInductionNegativeRay}{1,3}}
    \proofstepwidestar[6]{z\in\mathbb Z}{\rA}
    \proofstepwidestar[6]{\exists n\in\mathbb N\,
      (z=\IntEmbed(n)\lor z=-\IntEmbed(n))}{%
      \FormulaRefAuto{IntegerSignedNormalForm}{6}}
    \proofstepwidestar[8]{n\in\mathbb N\land
      (z=\IntEmbed(n)\lor z=-\IntEmbed(n))}{\rA}
    \proofstepwidestar[4,8]{P(\IntEmbed(n))}{\rRE{\rUE{4},\rAEa{8}}}
    \proofstepwidestar[5,8]{P(-\IntEmbed(n))}{\rRE{\rUE{5},\rAEa{8}}}
    \proofstepwidestar[11]{z=\IntEmbed(n)}{\rA}
    \proofstepwidestar[4,8,11]{P(z)}{\rIE{11,9}}
    \proofstepwidestar[13]{z=-\IntEmbed(n)}{\rA}
    \proofstepwidestar[5,8,13]{P(z)}{\rIE{13,10}}
    \proofstepwidestar[4,5,8]{P(z)}{%
      \rOE{\rAEb{8},11,12,13,14}}
    \proofstepwidestar[4,5,6]{P(z)}{\rEE{7,8,15}}
    \proofstepwidestar[1,2,3]{\forall z\in\mathbb Z\,P(z)}{%
      \rUI{\rRI{6,16}}}
  \closeproofpartwide
\end{tabproofsplitwide}

\section{Axiomatische Zusammenfassung}

Die folgenden Axiome erweitern die Peano-Struktur um einen überall
definierten Vorgänger. Anders als auf \(\mathbb N\) ist der Nachfolger auf
\(\mathbb Z\) bijektiv. Die totale Ordnung und die strikte
Nachfolgerorientierung verhindern zyklische Nachfolgermodelle; die
zweiseitige Induktion verbindet den positiven und den negativen Strahl.

Metalogisch ist dabei zwischen zwei Lesarten zu unterscheiden: Als
erststufiges Axiomenschema ist die Beschreibung nicht kategorisch und
besitzt nichtstandardmäßige Modelle. Unter voller zweitstufiger Lesart der
Induktion beschreibt sie dagegen die intendierte zweiseitige, von
\(0_{\mathbb Z}\) durch Nachfolger und Vorgänger erzeugte Kette.

\FormulaDefDeltaK[Erweiterte Peano-Axiome der ganzen Zahlen]{%
  \IntegerPeano{%
    \mathbb Z,0_{\mathbb Z},\IntSucc,\IntPred,\leq_{\mathbb Z}}%
}{IntegerPeanoAxioms}{%
  \DeltaRow{\textbf{Axiome}}{}%
  \DeltaRow{Null}{0_{\mathbb Z}\in\mathbb Z}
    [\FormulaRefAuto{IntegerPeanoZero}]%
  \DeltaRow{Nachfolger}{\IntSucc\colon\mathbb Z\to\mathbb Z}
    [\FormulaRefAuto{IntegerPeanoSuccFunction}]%
  \DeltaRow{Vorgänger}{\IntPred\colon\mathbb Z\to\mathbb Z}
    [\FormulaRefAuto{IntegerPeanoPredFunction}]%
  \DeltaRow{Vorgänger nach Nachfolger}
    {z\in\mathbb Z\vdash\IntPred(\IntSucc(z))=z}
    [\FormulaRefAuto{IntegerPeanoPredSucc}]%
  \DeltaRow{Nachfolger nach Vorgänger}
    {z\in\mathbb Z\vdash\IntSucc(\IntPred(z))=z}
    [\FormulaRefAuto{IntegerPeanoSuccPred}]%
  \DeltaRow{Totale Ordnung}
    {\TotOrd{\leq_{\mathbb Z},\mathbb Z}}
    [\FormulaRefAuto{IntegerPeanoOrder}]%
  \DeltaRow{Nachfolgerorientierung}
    {z\in\mathbb Z\vdash z<\IntSucc(z)}
    [\FormulaRefAuto{IntegerPeanoSuccPositive}]%
  \DeltaRow{Zweiseitige Induktion}
    {\begin{aligned}[t]
       &P(0_{\mathbb Z})\dsep
       \forall z\in\mathbb Z\,
         (P(z)\to P(\IntSucc(z)))\dsep{}\\
       &\forall z\in\mathbb Z\,
         (P(z)\to P(\IntPred(z)))
       \vdash\forall z\in\mathbb Z\,P(z)
     \end{aligned}}
    [\FormulaRefAuto{IntegerPeanoInduction}]%
  \DeltaRow{\textbf{Neue Symbole}}{}%
  \DeltaRow{Ganze Zahlen}{\mathbb Z}%
  \DeltaRow{Null}{0_{\mathbb Z}}%
  \DeltaRow{Nachfolger}{\IntSucc}%
  \DeltaRow{Vorgänger}{\IntPred}%
  \DeltaRow{Ordnung}{\leq_{\mathbb Z}}%
}

\FormulaThmDeltaK[Das Quotientenmodell erfüllt die Ganzzahlaxiome]{%
  \begin{aligned}[t]
    &\mathbb Z=\QuotientSet{D_{\mathbb Z}}{\IntRel}\ \land{}\\
    &\IntegerPeano{%
      \mathbb Z,0_{\mathbb Z},\IntSucc,\IntPred,\leq_{\mathbb Z}}
  \end{aligned}%
}{IntQuotientModelsIntegerPeano}{%
  \DeltaRow{Mengen}{z\dsep D_{\mathbb Z}}%
  \DeltaRow{Funktionen}{\IntSucc\dsep\IntPred}%
  \DeltaRow{Relationen}{\IntRel\dsep\leq_{\mathbb Z}}%
}
\begin{tabproofsplitwide}
  \proofpartwideindR[Träger und Abbildungen]{%
    \begin{aligned}[t]
      &0_{\mathbb Z}\in\mathbb Z\ \land\
      \IntSucc\colon\mathbb Z\to\mathbb Z\ \land{}\\
      &\IntPred\colon\mathbb Z\to\mathbb Z
    \end{aligned}
  }{%
    0_{\mathbb Z}\in\mathbb Z\land
    \IntSucc\colon\mathbb Z\to\mathbb Z\land
    \IntPred\colon\mathbb Z\to\mathbb Z
  }[IntegerPeanoModelCarrierMaps]
    \proofstepwidestar[]{0_{\mathbb Z}\in\mathbb Z}{%
      \FormulaRefAuto{IntZeroCarrier}}
    \proofstepwidestar[]{\IntSucc\colon\mathbb Z\to\mathbb Z}{%
      \FormulaRefAuto{IntegerSuccessorFunction}}
    \proofstepwidestar[]{\IntPred\colon\mathbb Z\to\mathbb Z}{%
      \FormulaRefAuto{IntegerPredecessorFunction}}
    \proofstepwidestar[]{0_{\mathbb Z}\in\mathbb Z\land
      \IntSucc\colon\mathbb Z\to\mathbb Z\land
      \IntPred\colon\mathbb Z\to\mathbb Z}{\rAI{1,\rAI{2,3}}}
  \closeproofpartwideindR

  \proofpartwideindR[Gegenseitige Inversen]{%
    \begin{aligned}[t]
      &\forall z\in\mathbb Z\,
        \IntPred(\IntSucc(z))=z\ \land{}\\
      &\forall z\in\mathbb Z\,
        \IntSucc(\IntPred(z))=z
    \end{aligned}
  }{%
    \forall z\in\mathbb Z\,\IntPred(\IntSucc(z))=z
    \land
    \forall z\in\mathbb Z\,\IntSucc(\IntPred(z))=z
  }[IntegerPeanoModelInverses]
    \proofstepwidestar[]{\forall z\in\mathbb Z\,
      \IntPred(\IntSucc(z))=z}{%
      \FormulaRefAuto{IntegerPredAfterSucc}}
    \proofstepwidestar[]{\forall z\in\mathbb Z\,
      \IntSucc(\IntPred(z))=z}{%
      \FormulaRefAuto{IntegerSuccAfterPred}}
    \proofstepwidestar[]{%
      \forall z\in\mathbb Z\,\IntPred(\IntSucc(z))=z
      \land
      \forall z\in\mathbb Z\,\IntSucc(\IntPred(z))=z}{\rAI{1,2}}
  \closeproofpartwideindR

  \proofpartwideindR[Geordnete Nachfolgerstruktur]{%
    \TotOrd{\leq_{\mathbb Z},\mathbb Z}
    \land\forall z\in\mathbb Z\,(z<\IntSucc(z))%
  }{%
    \TotOrd{\leq_{\mathbb Z},\mathbb Z}
    \land\forall z\in\mathbb Z\,(z<\IntSucc(z))%
  }[IntegerPeanoModelOrder]
    \proofstepwidestar[]{\TotOrd{\leq_{\mathbb Z},\mathbb Z}}{%
      \FormulaRefAuto{IntegerOrderTotal}}
    \proofstepwidestar[]{\forall z\in\mathbb Z\,
      (z<\IntSucc(z))}{\FormulaRefAuto{IntegerSuccessorPositive}}
    \proofstepwidestar[]{\TotOrd{\leq_{\mathbb Z},\mathbb Z}
      \land\forall z\in\mathbb Z\,(z<\IntSucc(z))}{\rAI{1,2}}
  \closeproofpartwideindR

  \proofpartwideindR[Induktionsschema]{%
    \begin{aligned}[t]
      &P(0_{\mathbb Z})\dsep
      \forall z\in\mathbb Z\,
      (P(z)\to P(\IntSucc(z)))\dsep{}\\
      &\forall z\in\mathbb Z\,
      (P(z)\to P(\IntPred(z)))
      \vdash\forall z\in\mathbb Z\,P(z)
    \end{aligned}
  }{%
    P(0_{\mathbb Z})\dsep
    \forall z\in\mathbb Z\,(P(z)\to P(\IntSucc(z)))\dsep
    \forall z\in\mathbb Z\,(P(z)\to P(\IntPred(z)))
    \vdash\forall z\in\mathbb Z\,P(z)
  }[IntegerPeanoModelInduction]
    \proofstepwidestar[1]{P(0_{\mathbb Z})}{\rA}
    \proofstepwidestar[2]{\forall z\in\mathbb Z\,
      (P(z)\to P(\IntSucc(z)))}{\rA}
    \proofstepwidestar[3]{\forall z\in\mathbb Z\,
      (P(z)\to P(\IntPred(z)))}{\rA}
    \proofstepwidestar[1,2,3]{\forall z\in\mathbb Z\,P(z)}{%
      \FormulaRefAuto{IntegerTwoSidedInductionModel}{1,2,3}}
  \closeproofpartwideindR

  \proofpartwide{Schluss}
    \proofstepwidestar[]{%
      \mathbb Z=\QuotientSet{D_{\mathbb Z}}{\IntRel}}{%
      \FormulaRefAuto{IntegersAsQuotient}}
    \proofstepwidestar[]{0_{\mathbb Z}\in\mathbb Z}{%
      \rAEa{\FormulaRefAuto{IntegerPeanoModelCarrierMaps}}}
    \proofstepwidestar[]{\IntSucc\colon\mathbb Z\to\mathbb Z}{%
      \rAEa{\rAEb{\FormulaRefAuto{IntegerPeanoModelCarrierMaps}}}}
    \proofstepwidestar[]{\IntPred\colon\mathbb Z\to\mathbb Z}{%
      \rAEb{\rAEb{\FormulaRefAuto{IntegerPeanoModelCarrierMaps}}}}
    \proofstepwidestar[]{z\in\mathbb Z\vdash
      \IntPred(\IntSucc(z))=z}{%
      \rAEa{\FormulaRefAuto{IntegerPeanoModelInverses}}}
    \proofstepwidestar[]{z\in\mathbb Z\vdash
      \IntSucc(\IntPred(z))=z}{%
      \rAEb{\FormulaRefAuto{IntegerPeanoModelInverses}}}
    \proofstepwidestar[]{\TotOrd{\leq_{\mathbb Z},\mathbb Z}}{%
      \rAEa{\FormulaRefAuto{IntegerPeanoModelOrder}}}
    \proofstepwidestar[]{z\in\mathbb Z\vdash z<\IntSucc(z)}{%
      \rAEb{\FormulaRefAuto{IntegerPeanoModelOrder}}}
    \proofstepwidestar[]{%
      P(0_{\mathbb Z})\dsep
      \forall z\in\mathbb Z(P(z)\to P(\IntSucc(z)))\dsep
      \forall z\in\mathbb Z(P(z)\to P(\IntPred(z)))
      \vdash\forall z\in\mathbb Z\,P(z)}{%
      \FormulaRefAuto{IntegerPeanoModelInduction}}
    \proofstepwidestar[]{%
      \IntegerPeano{%
        \mathbb Z,0_{\mathbb Z},\IntSucc,\IntPred,\leq_{\mathbb Z}}}{%
      \FormulaRefAuto{IntegerPeanoAxioms}{2,3,4,5,6,7,8,9}}
    \proofstepwidestar[]{%
      \mathbb Z=\QuotientSet{D_{\mathbb Z}}{\IntRel}
      \land
      \IntegerPeano{%
        \mathbb Z,0_{\mathbb Z},\IntSucc,\IntPred,\leq_{\mathbb Z}}}{%
      \rAI{1,10}}
  \closeproofpartwide
\end{tabproofsplitwide}

\FormulaAxiomAutoR[Null-Axiom der ganzen Zahlen]{%
  0_{\mathbb Z}\in\mathbb Z%
}{IntegerPeanoZero}
\par\noindent\emph{Herkunft aus der Quotientenkonstruktion:}
\FormulaRefAuto{IntZeroCarrier}[thm].\par

\FormulaAxiomAutoR[Nachfolgerfunktion der ganzen Zahlen]{%
  \IntSucc\colon\mathbb Z\to\mathbb Z%
}{IntegerPeanoSuccFunction}
\par\noindent\emph{Herkunft aus der Quotientenkonstruktion:}
\FormulaRefAuto{IntegerSuccessorFunction}[thm].\par

\FormulaAxiomAutoR[Vorgängerfunktion der ganzen Zahlen]{%
  \IntPred\colon\mathbb Z\to\mathbb Z%
}{IntegerPeanoPredFunction}
\par\noindent\emph{Herkunft aus der Quotientenkonstruktion:}
\FormulaRefAuto{IntegerPredecessorFunction}[thm].\par

\FormulaAxiomAutoR[Vorgänger nach Nachfolger]{%
  z\in\mathbb Z\vdash\IntPred(\IntSucc(z))=z%
}{IntegerPeanoPredSucc}
\par\noindent\emph{Herkunft aus der Quotientenkonstruktion:}
\FormulaRefAuto{IntegerPredAfterSucc}[thm].\par

\FormulaAxiomAutoR[Nachfolger nach Vorgänger]{%
  z\in\mathbb Z\vdash\IntSucc(\IntPred(z))=z%
}{IntegerPeanoSuccPred}
\par\noindent\emph{Herkunft aus der Quotientenkonstruktion:}
\FormulaRefAuto{IntegerSuccAfterPred}[thm].\par

\FormulaAxiomAutoR[Totale Ordnung der ganzen Zahlen]{%
  \TotOrd{\leq_{\mathbb Z},\mathbb Z}%
}{IntegerPeanoOrder}
\par\noindent\emph{Herkunft aus der Quotientenkonstruktion:}
\FormulaRefAuto{IntegerOrderTotal}[thm].\par

\FormulaAxiomAutoR[Strikte Nachfolgerorientierung]{%
  z\in\mathbb Z\vdash z<\IntSucc(z)%
}{IntegerPeanoSuccPositive}
\par\noindent\emph{Herkunft aus der Quotientenkonstruktion:}
\FormulaRefAuto{IntegerSuccessorPositive}[thm].\par

\FormulaAxiomAutoR[Zweiseitiges Induktionsaxiom]{%
  \begin{aligned}[t]
    &P(0_{\mathbb Z})\dsep
    \forall z\in\mathbb Z\,
      (P(z)\to P(\IntSucc(z)))\dsep{}\\
    &\forall z\in\mathbb Z\,
      (P(z)\to P(\IntPred(z)))
    \vdash\forall z\in\mathbb Z\,P(z)
  \end{aligned}%
}{IntegerPeanoInduction}
\par\noindent\emph{Herkunft aus der Quotientenkonstruktion:}
\FormulaRefAuto{IntegerTwoSidedInductionModel}[thm].\par

\end{document}
